\documentclass[11pt]{amsart}
\usepackage[T1]{fontenc}
\usepackage{lmodern,amsmath,amssymb,amsthm,mathtools,microtype,booktabs,array,longtable}
\usepackage[margin=1.06in]{geometry}
\usepackage{tikz}
\usepackage{xr-hyper}
\usepackage[hidelinks]{hyperref}
\hypersetup{pdftitle={Finite-Use Discrimination Geometry of Ordered Quantum Measurements},pdfauthor={Qian Qi}}
\emergencystretch=2em
\numberwithin{equation}{section}
\newtheorem{theorem}{Theorem}[section]
\newtheorem{lemma}[theorem]{Lemma}
\newtheorem{proposition}[theorem]{Proposition}
\newtheorem{corollary}[theorem]{Corollary}
\theoremstyle{definition}\newtheorem{definition}[theorem]{Definition}
\theoremstyle{remark}\newtheorem{remark}[theorem]{Remark}
\newcommand{\Q}{\mathbb Q}\newcommand{\R}{\mathbb R}\newcommand{\C}{\mathbb C}\newcommand{\N}{\mathbb N}\newcommand{\Z}{\mathbb Z}\newcommand{\E}{\mathbb E}
\newcommand{\ip}[2]{\langle #1,#2\rangle}\newcommand{\norm}[1]{\lVert #1\rVert}
\newcommand{\epsc}{\varepsilon_{\mathrm c}}
\DeclareMathOperator{\osc}{osc}\DeclareMathOperator{\spanop}{span}\DeclareMathOperator{\End}{End}\DeclareMathOperator{\diag}{diag}\DeclareMathOperator{\tr}{tr}\DeclareMathOperator{\rad}{rad}\DeclareMathOperator{\TV}{TV}
\title[Finite-use discrimination geometry]{Finite-Use Discrimination Geometry of Ordered Quantum Measurements}
\author{Qian Qi}
\date{6 October 2026}
\subjclass[2020]{Primary 81P15, 81P47; Secondary 94A17, 81P70}
\externaldocument{build/xref-supplement}[BINARY_SUPPLEMENT.pdf]
\begin{document}
\raggedbottom
\begin{abstract}
We study ordered quantum measurements with separate acquisition and
receiver resources.  For every finite two-call decision experiment,
the unit-reset optimum is an attained common-barycenter concave-envelope
problem; its dual consists of Hermitian majorants.  Rank-restricted
atom domains resolve the dimensions of both retained receiver registers
without importing a purification across a counted cut.  Pure atoms give
the classical value, and the resulting certificates characterize when
retained receiver memory is advantageous.
On a finite shared-basis experiment, the complete hierarchy has exact
equal-prior value $1/2+t^2(dr-1)/(2dr(d+1))$, where $r$ is the smaller
retained dimension.  We identify all initial Grams attaining this
value: they are precisely the density matrices $\rho\preceq I/r$.
Each has an attaining decomposition into at most $d$ commuting
normalized rank-$r$ projections.  Bounding the initial reference as
well replaces $r$ by the smallest of the three cut dimensions, and
fixed-subspace preparations attain the bound without feedback.
An explicit centered-swap deficit estimate gives common-projection
rigidity, an approximate initial decomposition, and a quadratic
spectral loss certificate.  An exact pinching identity treats the
exceptional qubit face.  The finite hierarchy persists under full-rank
deformation and calibrated changes of channels, priors and latent law.
The complementary local theory classifies fixed quadratic tangent
tubes by the complete support-covariance trichotomy: a profile with
$N=\sum n_j$ and $V=\sum n_j^2$ has trace orders
$\min\{1,\sqrt Ns\}$, $\min\{1,\sqrt Vs\}$ or $\min\{1,Ns\}$.
Hard reset-policy square budget $Q$ gives coherent order
$\min\{1,\sqrt Qs\}$.  The finite remainders and unknown-remainder
independence of the lower readouts are retained.  The exact variational
principle is general; the solved spectrum concerns the stated finite
family, and the local constants depend on the fixed tube.
\end{abstract}
\maketitle
\section{Introduction}\label{sec:introduction84}
A measurement used as an unknown device returns a classical label, but
its discrimination need not be classical.  Correlations may be created
inside an acquisition, retained in an external receiver, or transmitted
coherently into a later acquisition.  These are distinct resources.
This article gives a local geometric classification with sharp allocation
laws, a general variational principle for reset decisions, and exact
finite experiments separating receiver storage from coherent acquisition.  These results concern the same classical-output
interface and retain the order of its labels.

\subsection{Local geometry and exact experiments}
An ordered measurement is a tuple $E=(E_1,\ldots,E_k)$ of positive
$d\times d$ matrices with $\sum_jE_j=I$.  Its channel is
$\mathcal M_E(\rho)=\sum_j\tr(E_j\rho)|j\rangle\langle j|$.
No post-measurement device system is available.  Known preparations,
ancillary systems and controls are allowed.  Trace norms and diamond
norms below are unhalved; success probabilities are ordinary probabilities.

The first result resolves the local geometry for a prescribed allocation.
Fix $E$, a nonzero Hermitian zero-sum tuple $H$ in the one-sided positive
tangent cone, and a remainder allowance $\Lambda\ge0$.  For
$\|R\|_\Sigma=\sum_j\|R_j\|_{\rm op}$, consider every legal $F$ with
\begin{equation}\label{eq:fronttube87}
 \|F-E-sH\|_\Sigma\le\Lambda s^2.
\end{equation}
Write $P_j=\operatorname{supp}E_j$, $Q_j=I-P_j$,
$J_j=Q_jH_jQ_j$, $\Gamma_E(H)=\frac i2\sum_j[P_j,H_j]$, and
$\mathcal W_E=\sum_jP_j\mathcal H_dP_j$, where $\mathcal H_d$ is the
real Hermitian matrix space.  The cone condition is $J_j\succeq0$.
For positive integer acquisition sizes $\boldsymbol n=(n_1,\ldots,n_\ell)$,
put $N=\sum_rn_r$ and $V=\sum_rn_r^2$.
Theorem~\ref{thm:allocation89} proves, respectively, the orders
\begin{equation}\label{eq:frontprofile89}
 \min\{1,\sqrt N s\},\qquad \min\{1,\sqrt V s\},\qquad \min\{1,Ns\}
\end{equation}
when $H\in\operatorname{Ran}\mathcal C_E$, when all $J_j=0$ and
$\Gamma_E(H)\notin\mathcal W_E$, and when some $J_j\ne0$.
The support-kernel identity shows these cases exhaust the nonzero cone.
The same orders hold for independent complete acquisition groups and
for prescribed-profile reset feedback.  At a reset cut the entire fresh
probe, reference and within-block memory is conditionally independent
of the receiver's entire old memory.  The receiver may keep arbitrarily
large quantum registers and act on them jointly; it cannot import them
coherently into a fresh acquisition.  Preparations and controls are one
common rule on every formal history, including null histories.  The
conditional old receiver states need not coincide.

For a history-dependent policy $\pi$, hard pathwise reservations define
$N_\pi=\sup_\zeta\sum_{h\in\zeta}n(h)$ and
$Q_\pi=\sup_\zeta\sum_{h\in\zeta}n(h)^2$.
Theorem~\ref{thm:quadraticbudget89} proves that the optimum over
$N_\pi\le N$, $Q_\pi\le Q$, $N\le Q\le N^2$, has coherent order
$\min\{1,\sqrt Qs\}$; its other two rows are unchanged.
The constants and the interval $0<s\le s_0$ depend on the fixed
$E,H,\Lambda$, but not on the profile, budget or unknown remainder.
A common bounded readout gives the lower without learning that remainder.
The individual-policy upper retains both $\alpha\sqrt{Q_\pi}$ and
$\sqrt{rN_\pi}$; a large cost is not a lower information certificate
for a policy which erases its output.

The second result concerns exact success, rather than local orders.
For every $d\ge2$, a finite family of ordered bases with simultaneous
Haar second moments exists.  Compare the scalar measurement $C_y=I/d$
with a once-chosen member $E_y^b(t)=(1-t)I/d+tP_y^b$, using that same
device twice.  The decision is scalar versus the family.  Put
$L=2(d+1)-t^2$ and assign these alternatives priors
$(d+1-t^2)/L$ and $(d+1)/L$.  Theorem~\ref{thm:operationalhierarchy90}
gives the exact values
\begin{equation}\label{eq:fronthierarchy90}
\begin{split}
 P_{\rm ca}&=(d+1)/L,\\
 P_{\rm reset}&=(d+1)/L+t^2(d-1)/(2dL),\\
 P_{\rm all}&=(d+1+t^2)/L.
\end{split}
\end{equation}
Here the first class completely measures and reprepares between calls;
the second allows arbitrary unit-reset receiver feedback; the third
allows all adaptive controls.  For $t>0$ the inequalities are strict,
including the full-rank interval $0<t<1$.  Two independent maximally
entangled pairs attain the middle value; an antisymmetric two-input
state attains the last.  A swap-filter spectral identity bounds every
reset policy, including all intervening receiver instruments.
For the scalar qubit measurement and the six ordered Pauli measurements,
the exact values are $3/5$, $13/20$, and $4/5$.
This finite ensemble result neither asserts a gap for every fixed pair
nor identifies a coherent-memory-dimension hierarchy.

\subsection{A general optimization principle and equal priors}
Theorem~\ref{thm:resetvariational91} applies to every finite two-call
ordered-measurement decision experiment, with arbitrary priors and
rewards.  If $g_y(\rho)$ is the best fresh-reference value after a first
label without receiver feedback and $c_y$ is its upper concave envelope
on all density matrices, then
\[
 P_{\rm reset}=\max_{\rho}\sum_yc_y(\rho).
\]
A common first marginal is essential.  At most
$(\operatorname{rank}\rho)^2$ instrument outcomes per first label
suffice, and the quantum receiver has dimension at most $d_1d_2$.
Theorem~\ref{thm:resetdual91} supplies exact Hermitian majorants and
contact conditions.  Mixed density-matrix atoms cannot be replaced
silently by pure states.  This is a finite-dimensional characterization
of optimal values, not a polynomial-time optimization claim.

For the same finite basis devices, Theorem~\ref{thm:equalprior91}
removes the special prior in \eqref{eq:fronthierarchy90}.  With equal
hypothesis priors its exact values are
\[
 \frac12+\frac{t^2(d-1)}{2d(d+1)},\qquad
 \frac12+\frac{t^2(d-1)}{2d^2},\qquad
 \frac12+\frac{t^2}{2d}.
\]
For the seven qubit devices at $t=1$ these become
$7/12<5/8<3/4$.  The middle upper permits all receiver feedback,
but its attaining strategy uses none.  The alternative index is chosen
once for both calls; independent redraws give no signal.  The explicit
qubit calculation and a channel/prior/law stability estimate follow the
theorem.

Theorem~\ref{thm:classicalroof91} identifies the classical optimum by
restricting the envelope atoms to pure states.  Thus two explicit
convexifications distinguish complete classical readout from quantum
receiver storage in every finite experiment.
Corollary~\ref{cor:gapcertificate91} gives necessary and sufficient
primal--dual certificates for a strict gap, and contact certificates
for equality.  Their majorization inequalities are universal; no
sampling procedure is asserted to certify them.

These statements share one operational interface and one precise
fresh-preparation cut, but have different quantifiers.  The local profile
law classifies shrinking fixed-tube signals at arbitrary hard resources.
The variational principle treats exact two-call values for arbitrary
finite experiments.  The two basis theorems solve particular experiments
inside that general class.  Neither exact calculation supplies uniform
constants for the local theorem, and the local theorem is not a premise
of either exact score proof.

\subsection{Geometry underlying the allocation law}
The positive covariance operator $\mathcal C_E$ acts on zero-sum
Hermitian tuples and incorporates the coupled measurement normalization.
For $M=(E+F)/2$ and $H=F-E$, the global certificate is
\begin{equation}\label{eq:frontupper84}
 D_N(E,F)\le\min\left\{2,2\sqrt{k+1}
 \sqrt{N\langle H,(\mathcal C_M+N^{-1}\operatorname{Id})^{-1}H\rangle}
 \right\}.
\end{equation}
Theorem~\ref{thm:supportkernel84} identifies its complete kernel in
support coordinates.  This is an upper certificate for arbitrary
pairs, not an asserted global matching metric.  On a nonstationary
fixed unitary orbit $E_j(t)=e^{itG}E_je^{-itG}$, the finite-angle theorem is
\begin{equation}\label{eq:frontorbit84}
 D_N(E,E(t))\asymp_{E,G}
 \begin{cases}
 \min\{1,\sqrt N|t|\},&G\in\mathcal W_E,\\
 \min\{1,N|t|\},&G\notin\mathcal W_E.
 \end{cases}
\end{equation}
The support span is the Hermitian Kraus-product span.  The channel
metrological criterion and correction principle are credited to
Zhou--Jiang \cite{ZJ84}; the finite-pair estimates retain explicit
quadratic channel errors.  Normalized tangent realizations extend
these estimates to the whole legal tube, including unknown remainders
and one-sided support openings.

The width-cap consequence, for $1\le b\le N$, is
\begin{equation}\label{eq:frontwidth87}
 D_N^{[b]}(E,F)\asymp_{E,H,\Lambda}
 \begin{cases}
 \min\{1,\sqrt Ns\},&H\in\operatorname{Ran}\mathcal C_E,\\
 \min\{1,\sqrt{Nb}\,s\},&J_j=0\ (\forall j),
                          \Gamma_E(H)\notin\mathcal W_E,\\
 \min\{1,Ns\},&J_j\ne0\ (\exists j).
 \end{cases}
\end{equation}
The complete profile refines this cap through $V$, rather than through
its largest group alone.  The upper uses conditional fidelity, not
unconditional tensorization after feedback.  The lower first pays the
full $m_r\kappa s^2$ block errors and then combines unequal signals.
These distinctions are essential to uniformity in the unknown remainder.

The variational and equality statements also have quantitative forms.
Corollary~\ref{cor:contactdefect91} decomposes a certified reward loss
into nonnegative spectral, majorization, and leaf-decision defects.
Lemma~\ref{lem:stableswap91} and Corollary~\ref{cor:resetstability91}
control, at the relative scale $\varepsilon/t^2$, the deviation of
near-optimal reset Gram matrices from scalar matrices.  This concerns
the selected normal form, not unique dilations or physical calibration.


\subsection{A dimension-resolved reset hierarchy}
The two receiver resources can be resolved separately rather than
compared only at their classical and unrestricted endpoints.
At the first recording cut let the retained old receiver have dimension
at most $k$, and let the independently prepared second reference have
dimension at most $\ell$.  The initial reference and temporary
workspaces before that cut are not counted by these thresholds.
Theorem~\ref{thm:rankroof92} expresses the exact reward as the same
common-barycenter optimization, with atom rank at most $k$ and fresh
Gram rank at most $\ell$.  Both primal and dual are attained.
A dimension-preserving Kraus and mixture refinement, rather than a
quantum dilation retained across the cut, proves the operational
reduction.  Thus every finite experiment has an exact monotone
hierarchy whose endpoint is the full reset optimum.

For the fixed shared-basis family, this hierarchy is explicitly solved
at every pair of thresholds.  With $r=\min\{k,\ell\}$,
Theorem~\ref{thm:rankspectrum92} gives the equal-prior value
\[
 P_{k,\ell}^{\rm eq}=\frac12+\frac{t^2(dr-1)}{2dr(d+1)}.
\]
The consecutive gains are $t^2/[2d(d+1)r(r+1)]$.  They are strictly
positive for $t>0$, throughout the same full-rank deformation.
Cyclic rank-$r$ compression instruments attain them; matching upper
bounds cover all mixed, randomized and history-dependent competitors
in the specified class.  A corresponding exact formula holds for the
previous biased priors.  The saturation is at the full reset value,
not at the unrestricted coherent-acquisition value.  This is a
complete hierarchy for the two stated recording cuts and is not a
classification of auxiliary workspaces at every time in a general
multi-time process.

\subsection{Optimal initial spectra and quantitative rigidity}
The value hierarchy has a geometric equality set.  On the same
shared-basis family, let $r=\min\{k,\ell\}$ and fix the initial
Gram matrix $\rho$.  Theorem~\ref{thm:initialspectra93} shows that
its rank-$r$ optimum is attainable precisely when $\rho\preceq I/r$.
Every such matrix is a convex combination of at most $d$ commuting
rank-$r$ projections divided by $r$.  A finite interval construction
provides the recorded instrument, so the assertion is both necessary
and sufficient.  If the first reference is also bounded by $q$, the
exact value depends on $\min\{q,k,\ell\}$ and is attained without
feedback on a fixed subspace.  This counts three specified cuts,
not the maximum workspace at every time.

Away from the exceptional qubit rank-one face, the centered-swap
deficit controls squared trace distance of both Gram factors to a
common normalized rank-$r$ projection.  The resulting reward-loss
estimate supplies an approximate projection decomposition of the
initial Gram matrix and a spectral obstruction proportional to
$t^2[\tr(\rho-I/r)_+]^2$.  The constants are explicit, and the
relative scale is the score loss divided by $t^2$.
At the exceptional face the exact equality condition is commutation,
and an exact pinching identity replaces common-projection rigidity.
No uniqueness of instruments or uncalibrated device certification is
inferred from these statements.

\subsection{Notation, organization and scope}
For positive trace-class $R,S$, $f(R,S)=\|\sqrt R\sqrt S\|_1$ is root
fidelity.  For states, $d_B(R,S)^2=2(1-f(R,S))$; the inherited symbol
$\beta$ for a state Bures distance denotes this same distance when used
with two state arguments.  Named coefficients such as $\beta_h$ are
bound coefficients, not an additional distance.  $D(\pi;E,F)$ denotes
one policy's separation; a resource subscript denotes a supremum.
$V/N$ is an effective allocation width, not an integer or memory dimension.

The proof proceeds from covariance and its support kernel to normalized
curves, tangent neighborhoods, complete groups, conditional reset
accumulation, and hard quadratic budgets.  Section~\ref{sec:operationalmemory90}
then gives the exact finite hierarchy and its stability.  Balanced-interior
geometry, coding, learning and affine repair are retained with their full
proofs in the current linked supplement.  The complete research edition
also retains every earlier mathematical section.  The independent
structural article is not a premise of either main result.

The fixed-tube constants may deteriorate at changing support floors,
vanishing tangents or closing correction gaps.  Neither local theorem
is a uniform full-boundary metric, entropy or learning theorem.  The
finite hierarchy supplies a different, exact, all-dimension statement on
its specified ensemble family.  The historical analytic programme and
independent specialist priority assessment are separate obligations.
Our comparison with the memory-dimension hierarchy of Zonnios--Binder
\cite{ZB90} and with accessible post-measurement states \cite{EQ90}
keeps these interfaces and resources distinct.

\section{A covariance principle on the closed measurement body}
\label{sec:covariance83}

Let $\mathcal P_{d,k}$ denote the complete body of ordered $k$-outcome
measurements on $\C^d$: its elements are Hermitian tuples
$E=(E_1,\ldots,E_k)$ with $E_j\succeq0$ and $\sum_jE_j=I_d$.
Here $d\ge1$, $k\ge2$, and $N\ge1$ are integers.  Zero effects,
changing ranks, and noncommuting supports are allowed throughout
this section.  The operational distance $D_N$ uses unhalved trace
norm, retained references, common adaptive control, and at most $N$
calls with bounded public stopping.  The device consumes its input
and returns only its ordered classical label.

The normalized tangent space and its inner product are
\[
 \mathcal T_{d,k}=\{H_j=H_j^*: \textstyle\sum_jH_j=0\},\qquad
 \langle H,K\rangle=\sum_j\tr(H_jK_j).
\]
All operator inequalities below on tuple spaces refer to this real
Hilbert structure, unless stated otherwise.  For any Hermitian tuple
$K$, not necessarily of zero sum, put
\begin{align}
 S_E(K)&=\sum_jE_jK_j,\label{eq:covariances83}\\
 (\mathcal C_EK)_j
  &=\tfrac12\bigl(E_jK_j+K_jE_j-E_jS_E(K)-S_E(K)^*E_j\bigr).
                                      \label{eq:covarianceoperator83}
\end{align}
The matrix $S_E(K)$ need not be Hermitian.  Keeping its adjoint in
\eqref{eq:covarianceoperator83} is essential.  The operator
$\mathcal C_E$ is polynomial in the effects, annihilates constant
tuples, and has range in $\mathcal T_{d,k}$.  In particular its
restriction to $\mathcal T_{d,k}$ is well defined.

For $\tau=N^{-1}$ set
\begin{align}
 g_{N,E}(H)^2
  &=N\langle H,(\mathcal C_E+\tau\operatorname{Id})^{-1}H\rangle,
                                      \label{eq:covarianceform83}\\
 \mathcal Q_N(E,F)&=g_{N,(E+F)/2}(F-E).
                                      \label{eq:covariancemidpoint83}
\end{align}
The inverse acts on $\mathcal T_{d,k}$.  Positivity proved below
makes it nonsingular even on the spectral boundary.  We do not
identify $\mathcal Q_N$ with an operational distance or assume a
triangle inequality for it.

\begin{theorem}[Closed-body adaptive bound]
\label{thm:closedcovariance83}
For every $E,F\in\mathcal P_{d,k}$,
\begin{equation}\label{eq:closedcovariance83}
 D_N(E,F)\le
 \min\{2,\,2\sqrt{k+1}\,\mathcal Q_N(E,F)\}.
\end{equation}
More generally, every continuously differentiable path $E(t)$ in
$\mathcal P_{d,k}$ satisfies
\begin{equation}\label{eq:covariancepath83}
 D_N(E(0),E(1))\le
 \sqrt{k+1}\int_0^1g_{N,E(t)}(E'(t))\,dt.
\end{equation}
These statements are uniform over all ranks and support patterns.
For $k=2$, $E=(A,I_d-A)$ and $H=(B,-B)$, the restriction is exactly
\begin{equation}\label{eq:binarycovariance83}
 (\mathcal C_EH)_1=A(I_d-A)B+BA(I_d-A),\qquad
 \mathcal Q_N(E,F)^2=2Q_N(A,F_1)^2,
\end{equation}
where $Q_N$ is the binary comparison modulus in
\eqref{eq:matrixmodulus76}.
\end{theorem}

Thus the binary variance operator has a normalized-tuple extension,
not an independently regularized list of effects.  The theorem gives
a global upper certificate.  The two-sided balanced-interior
classification and its entropy theorem remain
Theorems~\ref{thm:multioutcomemetric82} and
\ref{thm:multioutcomeentropy82}; a matching general lower bound for
\eqref{eq:covariancemidpoint83} is a separate question.

\subsection{Positive covariance and horizontal decomposition}

\begin{lemma}[Covariance identities]
\label{lem:covarianceidentities83}
On all Hermitian tuples, $\mathcal C_E$ is self-adjoint and positive
semidefinite, and
\begin{align}
 \langle K,\mathcal C_EK\rangle
 &=\sum_j\tr(E_jK_j^2)-\|S_E(K)\|_{\rm HS}^2\notag\\
 &=\sum_j\|E_j^{1/2}(K_j-S_E(K))\|_{\rm HS}^2.
                                      \label{eq:covariancevariance83}
\end{align}
For $E_t=(1-t)E+tF$, $0\le t\le1$,
\begin{equation}\label{eq:covarianceconcavity83}
 \langle K,[\mathcal C_{E_t}-(1-t)\mathcal C_E-t\mathcal C_F]K\rangle
 =t(1-t)\|S_F(K)-S_E(K)\|_{\rm HS}^2.
\end{equation}
The kernel consists precisely of tuples satisfying
$E_j^{1/2}(K_j-S_E(K))=0$ for every $j$.
Moreover $\mathcal C_E=0$ on $\mathcal T_{d,k}$ if and only if the
$E_j$ are mutually orthogonal projections, with zero projections
permitted.
\end{lemma}
\begin{proof}
Write $A_j=E_j^{1/2}$ and let $A$ be the vertical stack of the $A_j$.
Then $A^*A=I_d$, so $P=I-AA^*$ is an orthogonal projection.  On
complex matrix tuples with real Hilbert--Schmidt inner product,
define
\[
 (T_AB)_j=A_j^*B_j+B_j^*A_j.
\]
Its real adjoint is $(T_A^*K)_j=2A_jK_j$.  Direct multiplication gives
\begin{equation}\label{eq:covariancefactor83}
 \mathcal C_E=\tfrac14 T_APT_A^*.
\end{equation}
This proves self-adjointness, positivity, and
\eqref{eq:covariancevariance83}.  It also proves the kernel claim.
Constant tuples are killed and the displayed formula has zero sum,
so no extra projection onto $\mathcal T_{d,k}$ is implicit.

In the first line of \eqref{eq:covariancevariance83}, the first term
is linear in $E$ and $S_E(K)$ is linear in $E$.  Expanding its squared
norm gives \eqref{eq:covarianceconcavity83}.

If the restriction of $\mathcal C_E$ to $\mathcal T_{d,k}$ vanishes,
subtracting the arithmetic mean from any tuple shows that its form
vanishes on all tuples.  Take $K_j=I_d$ at one index and zero
elsewhere.  Its form is $\tr(E_j-E_j^2)=0$.  Since $0\preceq E_j
\preceq I_d$, this forces $E_j^2=E_j$.  Their sum being $I_d$, they
are mutually orthogonal: each positive term in
$\sum_{i\ne j}E_jE_iE_j=0$ must vanish.  Conversely, for orthogonal
projections the summands $E_jK_j$ have mutually orthogonal ranges,
and $\|\sum_jE_jK_j\|_{\rm HS}^2=\sum_j\|E_jK_j\|_{\rm HS}^2$.
Their covariance form is zero.
\end{proof}

\begin{lemma}[Regularized horizontal energy]
\label{lem:horizontalenergy83}
For $H\in\mathcal T_{d,k}$ and $\tau>0$,
\begin{equation}\label{eq:energymin83}
 \langle H,(\mathcal C_E+\tau\operatorname{Id})^{-1}H\rangle
 =\min_{\substack{A^*B=0,\ R\in\mathcal T_{d,k}\\T_AB+R=H}}
       \left(4\|B\|_{\rm HS}^2+\tau^{-1}\|R\|_{\rm HS}^2\right).
\end{equation}
If $K=(\mathcal C_E+\tau\operatorname{Id})^{-1}H$, a minimizer is
\begin{equation}\label{eq:energyoptimizer83}
 B_j=\tfrac12 A_j(K_j-S_E(K)),\qquad R=\tau K.
\end{equation}
\end{lemma}
\begin{proof}
The constraint $A^*B=0$ means $PB=B$.  Thus
$B=PT_A^*K/4$ gives $T_AB=\mathcal C_EK$ and is feasible with
$R=\tau K$.  Its objective is
$\langle K,\mathcal C_EK\rangle+\tau\|K\|^2=\langle H,K\rangle$.
For any feasible perturbation $(\Delta B,\Delta R)$, its linear
objective term is
\[
 8\langle B,\Delta B\rangle+2\tau^{-1}\langle R,\Delta R\rangle
 =2\langle K,T_A\Delta B+\Delta R\rangle=0.
\]
The remaining terms are nonnegative.  This proves minimality,
including at singular $E$; no inverse of an effect was used.
\end{proof}

\subsection{Adaptive differentiation and boundary passage}

\begin{proof}[Proof of Theorem~\ref{thm:closedcovariance83}]
First suppose every effect at the point of differentiation is
positive definite.  The tangent $H_0=T_AB$ in
\eqref{eq:energyoptimizer83} has zero sum.  For small real $s$,
$E_j+s(H_0)_j$ is positive and the tuple remains normalized.
If $R_j(s)$ is its positive square root, then
\[
 U_j=(B_j-R_j'(0))A_j^{-1}
\]
is anti-Hermitian.  Indeed, multiplication of $U_j+U_j^*$ on the
left and right by $A_j$ gives
$A_jB_j+B_j^*A_j-A_jR_j'(0)-R_j'(0)A_j=0$.
The factors $\exp(sU_j)R_j(s)$ therefore realize the derivative
$B_j$ while preserving the effects exactly.

Use the measurement isometry with an observed label and an
inaccessible duplicate label, together with the factor environment.
At the point under consideration it obeys
\[
 W^*\dot W=\sum_jA_j^*B_j=0,\qquad
 \|\dot W\|_{\rm op}\le\|B\|_{\rm HS}.
\]
Purify any fixed common $N$-slot tester.  In the inner product of
terms differentiated at distinct slots, cancel the common later
isometries; the later differentiated slot contains $W^*\dot W$ or
its adjoint and vanishes.  Thus the observed output derivative in
the direction $H_0$ has trace norm at most
$2\sqrt N\|B\|_{\rm HS}$.  This holds for every reference and
memory dimension.  The environments used for comparison are not
returned to the tester.

The output derivative of the same tester is linear in the channel
tangent.  For the residual tuple $R$, one call on a joint state
$\rho$ has reference blocks $\tr_A[(R_j\otimes I)\rho]$.
Duality against Hermitian contractions bounds the trace norm of the
$j$th block by $\|R_j\|_{\rm op}$.  Inserting the residual at each
slot and using trace-norm contraction hence costs at most
\[
 N\sum_j\|R_j\|_{\rm op}\le N\sqrt{k}\|R\|_{\rm HS}.
\]
This is a differential decomposition, not a proposed physical mixture
of measurements.  With $\tau=N^{-1}$, Cauchy--Schwarz and
\eqref{eq:energymin83} now give
\begin{align*}
 2\sqrt N\|B\|_{\rm HS}+N\sqrt{k}\|R\|_{\rm HS}
 &\le\sqrt{k+1}
       \sqrt{4N\|B\|_{\rm HS}^2+N^2\|R\|_{\rm HS}^2}\\
 &=\sqrt{k+1}\,g_{N,E}(H).
\end{align*}
Integrate for the fixed tester, then take the supremum.  Padding a
publicly stopped experiment with discarded dummy calls makes its
original stopped output a common postprocessing of an $N$-slot
experiment, so the same bound applies.

For a boundary path, use $E^\epsilon(t)=(1-\epsilon)E(t)+\epsilon U$,
where $U_j=I_d/k$ and $0<\epsilon<1$.  The covariance is polynomial
and positive semidefinite, so its regularized inverse is continuous.
Moreover $g_{N,E}(H)\le N\|H\|_{\rm HS}$, uniformly on the closed
body.  Dominated convergence applies to the integrals.  The hybrid
bound $D_N(E,F)\le N\sum_j\|E_j-F_j\|_{\rm op}$ gives convergence
of the endpoint distances.  This proves \eqref{eq:covariancepath83}
on the entire body.

For the segment $E(t)=(1-t)E+tF$, put $M=(E+F)/2$ and $H=F-E$.
Concavity and positivity imply, for $0<t\le1/2$,
\[
 \mathcal C_{E(t)}+\tau\operatorname{Id}
 \succeq2t(\mathcal C_M+\tau\operatorname{Id}).
\]
Inverse order bounds $g_{N,E(t)}(H)$ by $(2t)^{-1/2}g_{N,M}(H)$;
use $2(1-t)$ on the second half.  The two scalar integrals sum to
$2$, proving \eqref{eq:closedcovariance83}.  Finally, substituting
$K=(B,-B)$ into \eqref{eq:covarianceoperator83} gives
\eqref{eq:binarycovariance83}.  The factor two in the quadratic form
comes from the two equal Hilbert--Schmidt contributions.
\end{proof}

\subsection{Classical postprocessing}

\begin{proposition}[Data processing of horizontal energy]
\label{prop:covarianceprocessing83}
Let $T=(T_{aj})$ be a column-stochastic classical channel and
$F_a=\sum_jT_{aj}E_j$.  Write $(T^*L)_j=\sum_aT_{aj}L_a$.
Then, on all Hermitian tuples,
\begin{equation}\label{eq:covarianceprocessing83}
 \langle L,\mathcal C_FL\rangle
 \ge\langle T^*L,\mathcal C_E(T^*L)\rangle.
\end{equation}
Consequently the extended quadratic energy
\[
 \mathcal I_E(H)=\sup_K
       \{2\langle H,K\rangle-\langle K,\mathcal C_EK\rangle\}
\]
satisfies $\mathcal I_F(TH)\le\mathcal I_E(H)$.  The value may be
infinite at a boundary; when finite it is
$\langle H,\mathcal C_E^\dagger H\rangle$.
\end{proposition}
\begin{proof}
The two $S$ matrices coincide.  The difference of the covariance
forms is therefore
\[
 \sum_j\tr\!\left[E_j\left(\sum_aT_{aj}L_a^2
                   -(\sum_aT_{aj}L_a)^2\right)\right]\ge0.
\]
The bracket equals the positive sum
$\tfrac12\sum_{a,b}T_{aj}T_{bj}(L_a-L_b)^2$.
Substitution in the variational definition and enlargement from
$K=T^*L$ to all $K$ prove the energy inequality.  The pseudoinverse
statement follows by diagonalizing the real self-adjoint covariance;
a component along its kernel makes the supremum infinite.
\end{proof}

\subsection{A uniform noise crossover, sharp in horizon and noise}

\begin{theorem}[Regularization of arbitrary support patterns]
\label{thm:noisecrossover83}
For $0\le\epsilon\le1$ put $E^\epsilon=(1-\epsilon)E+\epsilon U$.
On $\mathcal T_{d,k}$,
\begin{equation}\label{eq:covariancenoisefloor83}
 \mathcal C_{E^\epsilon}\succeq
       (1-\epsilon)\mathcal C_E+\frac\epsilon k\operatorname{Id}.
\end{equation}
In particular, for arbitrary $E,F\in\mathcal P_{d,k}$,
\begin{equation}\label{eq:noisecrossoverupper83}
 D_N(E^\epsilon,F^\epsilon)\le
 \min\left\{2,
 \frac{2\sqrt{k+1}(1-\epsilon)N}
      {\sqrt{1+N\epsilon/k}}\|F-E\|_{\rm HS}\right\}.
\end{equation}
The dependence on $N$ and $\epsilon$ is attained, up to constants
depending only on $k$, by a two-dimensional family.  Namely, let
$P=|0\rangle\langle0|$, $P_\theta=|u_\theta\rangle\langle u_\theta|$,
$u_\theta=\cos\theta|0\rangle+\sin\theta|1\rangle$, and
\[
 E=(P,I_2-P,0,\ldots,0),\qquad
 F=(P_\theta,I_2-P_\theta,0,\ldots,0).
\]
For $0\le\theta\le\pi/2$, $0\le\epsilon\le1/2$, and
$z=(1-\epsilon)N\sin\theta/\sqrt{1+N\epsilon}$,
\begin{equation}\label{eq:noisecrossoverlower83}
 \frac{\min\{1,z\}}{16384}
 \le D_N^{\rm na}(E^\epsilon,F^\epsilon)
 \le D_N(E^\epsilon,F^\epsilon)
 \le4\sqrt{k(k+1)}\min\{1,z\}.
\end{equation}
\end{theorem}
\begin{proof}
For $K\in\mathcal T_{d,k}$, $S_U(K)=0$ and
$\mathcal C_UK=K/k$.  Concavity proves
\eqref{eq:covariancenoisefloor83}.  Apply it at the midpoint and use
Theorem~\ref{thm:closedcovariance83} to obtain
\eqref{eq:noisecrossoverupper83}.

For the displayed family, keep labels one and two and send every
other label to an independent fair bit.  This common postprocessing
produces the binary effects
$B=\beta P+\epsilon I_2/2$ and
$B_\theta=\beta P_\theta+\epsilon I_2/2$, where $\beta=1-\epsilon$.
The equality holds on subnormalized reference states.  It therefore
preserves the nonadaptive lower witnesses as well as channel
postprocessing.  The midpoint eigenvalues are
$(1\pm\beta\cos\theta)/2$, and direct substitution in
\eqref{eq:matrixcoordinates76} gives
\[
 Q_N(B,B_\theta)^2
 =\frac{4N^2\beta^2\sin^2\theta}
             {2+N(1-\beta^2\cos^2\theta)}
 \ge\frac{4z^2}{2+z^2}.
\]
Here $1-\beta^2\le2\epsilon$ and $N\ge1$ were used.
It follows that $\min\{1,Q_N(B,B_\theta)\}\ge\min\{1,z\}$.
The binary lower theorem in dimension two gives the first inequality
of \eqref{eq:noisecrossoverlower83}.  For the upper,
$\|F-E\|_{\rm HS}=2\sin\theta$ and
$(1+N\epsilon)/(1+N\epsilon/k)\le k$.
Combine \eqref{eq:noisecrossoverupper83} with $D_N\le2$.
\end{proof}

The family passes from projections and zero effects at $\epsilon=0$
to strictly positive effects at every $\epsilon>0$.  Thus a bound
uniform over the closed body must accommodate both the $N$ scale
and the $\sqrt{N/\epsilon}$ scale.  This conclusion does not assert a
matching entropy law on all multi-outcome boundary strata.

\subsection{Exact evaluation without effect square roots}

\begin{proposition}[Polynomial-size rational certificate]
\label{prop:covarianceexact83}
For Gaussian-rational $E,F\in\mathcal P_{d,k}$ and integer $N\ge1$,
$\mathcal Q_N(E,F)^2$ is rational and can be evaluated by exact
linear algebra on a positive definite real system of dimension
$p=(k-1)d^2$.  The bit complexity is polynomial in $d,k$, the total
input bit length, and $\log N$.  Legality is decidable by exact
positive-semidefinite tests.  No effect square root, ODE integration,
rank decision by tolerance, or enumeration of a dictionary is required.
\end{proposition}
\begin{proof}
Use the rational Hermitian matrix basis consisting of diagonal
units and real and imaginary off-diagonal pairs.  For each of the
first $k-1$ slots, insert one matrix basis element there and its
negative in the last slot.  These tuples $B_a$ form a rational basis
of $\mathcal T_{d,k}$; they need not be orthonormal.  Form
\[
 G_{ab}=\langle B_a,B_b\rangle,\quad
 L_{ab}=\langle B_a,\mathcal C_{(E+F)/2}B_b\rangle,\quad
 b_a=\langle B_a,F-E\rangle.
\]
All entries are rational, $G$ is positive definite, and $L\succeq0$.
Solve $(L+N^{-1}G)x=b$.  Then
$\mathcal Q_N^2=N b^{\mathsf T}x$.  Omitting the Gram matrix would
compute a different regularization.  Fraction-free elimination
performs polynomially many rational arithmetic operations; clearing
input denominators and applying determinant bounds gives polynomial
bit length for all minors and for the output.  The same reasoning
applies to exact symmetric elimination for each effect's positivity.
The complexity statement concerns this certificate evaluator, not
the exhaustive public code or collective-readout synthesis.
\end{proof}

\section{A finite product estimate}\label{sec:product85}
All classical distances in this paper are unhalved $\ell^1$ distances.
The following elementary estimate is the Bernoulli-product input to
Theorems~\ref{thm:orbittrichotomy84} and \ref{thm:multioutcomemetric82}.
We include its proof so that the support and discrimination theorems
can be checked without the binary supplement.

\begin{lemma}[A uniform finite Bernoulli lower bound]\label{lem:bernoulli85}
For $p,q\in[0,1]$ and every integer $N\ge1$,
\begin{equation}\label{eq:bernoulli85}
 \|\operatorname{Ber}(p)^{\otimes N}-
       \operatorname{Ber}(q)^{\otimes N}\|_1
 \ge \frac1{128}\min\{1,\sqrt N|p-q|\}.
\end{equation}
In particular, if a differentiable two-sided probability curve satisfies
$p'(0)\ne0$, then, on a sufficiently small two-sided interval,
its product experiment separates $p(0)$ from $p(t)$ by at least
$c\min\{1,\sqrt N|t|\}$ for every $N$, with $c>0$ depending
on that curve.  Necessarily $0<p(0)<1$.
\end{lemma}
\begin{proof}
Set $h=|p-q|$.  First let $m\ge1$ and $\sqrt m h\le1/4$.
For $u=m^{-1/2}$ define $z(a)=1-a+ae^{iu}$, $0\le a\le1$.
The argument of $z(a)$ has a continuous choice since its real part
is positive.  Direct differentiation and $0<u\le1$ give
\[
 |z(a)|^2\ge1-\frac1{4m},\qquad
 \frac u2\le\frac{d}{da}\arg z(a)
       =\frac{\sin u}{|z(a)|^2}\le\frac{4u}3.
\]
Bernoulli's inequality gives $|z(a)|^m\ge\sqrt{3/4}>3/4$.
The arguments of $z(p)^m$ and $z(q)^m$ differ by an angle
$\vartheta\in[\sqrt m h/2,4\sqrt m h/3]\subset[0,1/3]$.
For two complex numbers of moduli at least $3/4$,
\[
 |z(p)^m-z(q)^m|\ge\tfrac32\sin(\vartheta/2)
 \ge\frac{3}{2\pi}\vartheta\ge\frac1{16}\sqrt m h.
\]
These complex numbers are the expectations of
$\exp(iu\sum_{a=1}^m X_a)$ under the two Bernoulli product laws.
This statistic has modulus one, so their $\ell^1$ distance is at least
its expectation difference.

If $h\le1/4$ and $\sqrt N h\le1/4$, use $m=N$.
If $h\le1/4$ and $\sqrt N h>1/4$, use
$m=\lfloor(16h^2)^{-1}\rfloor\le N$.
Then $1/(4\sqrt2)\le\sqrt m h\le1/4$, since
$\lfloor x\rfloor\ge x/2$ for $x\ge1$.
Discarding the other observations is a common processing and gives
a lower bound $1/(64\sqrt2)>1/128$.
For $h>1/4$ the one-copy distance is $2h>1/2$.
The case $h=0$ is immediate.  Finally differentiability gives
$|p(t)-p(0)|\ge|p'(0)||t|/2$ near zero.  A two-sided probability
curve cannot have a nonzero first derivative at either endpoint.
\end{proof}

\section{Support geometry of the covariance kernel}\label{sec:support84}
Let $\mathcal H_d$ be the real Hilbert space of Hermitian $d\times d$
matrices, with inner product $\tr(AB)$.  Throughout this section,
$E=(E_1,\ldots,E_k)$ is an arbitrary ordered measurement:
$E_j\succeq0$ and $\sum_jE_j=I_d$.  Put
\[
 P_j=\operatorname{supp}E_j,\qquad Q_j=I_d-P_j,\qquad r_j=\operatorname{rank}E_j.
\]
The support of a zero effect is the zero projection.  The real subspaces
\begin{equation}\label{eq:supportspan84}
 \mathcal W_E=\sum_{j=1}^k P_j\mathcal H_dP_j,
 \qquad \mathcal A_E=\mathcal W_E^\perp
   =\{A\in\mathcal H_d:P_jAP_j=0\text{ for every }j\}
\end{equation}
will be called the \emph{support span} and its orthogonal complement.
The sum in \eqref{eq:supportspan84} is a sum of real linear subspaces,
not an assertion that $\mathcal W_E$ is an algebra.  In particular,
for higher-rank effects it need not equal the span of the effects themselves.
We use the zero-sum tuple space
$\mathcal T=\{K\in\mathcal H_d^k:\sum_jK_j=0\}$ with the sum
Hilbert--Schmidt inner product, and the covariance of
Lemma~\ref{lem:covarianceidentities83}.

\begin{theorem}[Complete support description of the kernel]\label{thm:supportkernel84}
For $A\in\mathcal A_E$ and $Z_j\in Q_j\mathcal H_dQ_j$, set
\begin{equation}\label{eq:kernelparam84}
 X_j=i[P_j,A]+Z_j,\qquad
 K_j=X_j-\frac1k\sum_{\ell=1}^kX_\ell.
\end{equation}
This is a real linear bijection from
$\mathcal A_E\oplus\bigoplus_jQ_j\mathcal H_dQ_j$ onto
$\ker(\mathcal C_E|_{\mathcal T})$.  Consequently
\begin{equation}\label{eq:kerneldimension84}
 \dim\ker(\mathcal C_E|_{\mathcal T})
 =d^2-\dim\mathcal W_E+\sum_j(d-r_j)^2.
\end{equation}
In particular the kernel, and not just its dimension, depends only
on the ordered support projections.
\end{theorem}
\begin{proof}
Write $S=\sum_jE_jK_j$.  The nonnegative variance decomposition in
Lemma~\ref{lem:covarianceidentities83} shows that $K$ is in the kernel
exactly when
\begin{equation}\label{eq:supportkernelcondition84}
 P_jK_j=P_jS\quad\text{for every }j.
\end{equation}
Indeed multiplication by $\sqrt{E_j}$ has precisely the same kernel
as multiplication by $P_j$; no inverse at a singular effect is used.
Decompose $S=B+iA$, where $A,B$ are Hermitian.  Since
$P_jK_jP_j$ and $P_jBP_j$ are Hermitian, compressing
\eqref{eq:supportkernelcondition84} gives $P_jAP_j=0$.
Furthermore $P_j(K_j-B)=iP_jA$, with the adjoint identity on the
other side.  These determine the two off-diagonal support blocks
and the support block of $K_j-B$; its remaining block is arbitrary.
Thus
\[
 K_j=B+i[P_j,A]+Z_j,\qquad Z_j=Q_jZ_jQ_j=Z_j^*.
\]
The condition $\sum_jK_j=0$ forces
$B=-k^{-1}\sum_j(i[P_j,A]+Z_j)$, proving necessity.

Conversely, let $A,Z_j$ be as stated.  Because $P_jAP_j=0$,
$E_jAP_j=0$, and hence
\[
 \sum_jE_jX_j=i\sum_jE_jA=iA.
\]
For the mean-subtracted tuple, therefore, $S=iA-\overline X$,
where $\overline X=k^{-1}\sum_jX_j$ is Hermitian.
We have $P_jX_j=iP_jA$, which proves
\eqref{eq:supportkernelcondition84}.  To check injectivity, suppose
that all $K_j$ vanish.  Then all $X_j$ equal a common Hermitian
matrix $D$, so the last display gives $D=iA$.  A matrix that is
both Hermitian and anti-Hermitian is zero.  Therefore $A=D=0$ and
all $Z_j=0$.  Finally $\dim(Q_j\mathcal H_dQ_j)=(d-r_j)^2$;
rank--nullity in \eqref{eq:supportspan84} proves
\eqref{eq:kerneldimension84}.
\end{proof}

\begin{proposition}[Unitary directions and the support span]\label{prop:unitaryrange84}
Let $G\in\mathcal H_d$ and $H_j=i[G,E_j]$.  Then
\begin{equation}\label{eq:unitaryrange84}
 H\in\operatorname{Ran}(\mathcal C_E|_{\mathcal T})
 \quad\Longleftrightarrow\quad G\in\mathcal W_E.
\end{equation}
More precisely, the pairing of $H$ with the kernel tuple
\eqref{eq:kernelparam84} is $-2\tr(GA)$.
\end{proposition}
\begin{proof}
We have $Q_jH_jQ_j=0$, so the $Z_j$ terms have zero pairing with $H$;
its zero tuple sum also removes $\overline X$.  Trace cyclicity gives
\[
 \sum_j\tr\bigl(H_ji[P_j,A]\bigr)
 =\tr\left(A\sum_ji[H_j,P_j]\right).
\]
Expanding in the displayed order yields
\[
 \sum_ji[H_j,P_j]
 =-2G+\sum_j(E_jGP_j+P_jGE_j).
\]
Every summand in the last sum belongs to
$P_j\mathcal H_dP_j$; it is therefore orthogonal to $A$.
The pairing is $-2\tr(GA)$ as asserted.  A self-adjoint linear
operator in finite dimension has range equal to the orthogonal
complement of its kernel.  Theorem~\ref{thm:supportkernel84} now
proves \eqref{eq:unitaryrange84}.
\end{proof}

\begin{corollary}[Regularized tangent scales]\label{cor:tangentscales84}
Let $\Pi_0$ be the orthogonal projection onto
$\ker(\mathcal C_E|_{\mathcal T})$, and let $u_\ell$ be an orthonormal
basis of its complement with $\mathcal C_Eu_\ell=\lambda_\ell u_\ell$,
$\lambda_\ell>0$.  For every $H\in\mathcal T$,
\begin{equation}\label{eq:tangentspectral84}
 g_{N,E}(H)^2
 =N^2\|\Pi_0H\|_{\mathrm{HS}}^2+
  N\sum_\ell\frac{|\langle H,u_\ell\rangle|^2}{\lambda_\ell+N^{-1}}.
\end{equation}
Thus a nonzero horizontal direction has order $\sqrt N$, whereas a
direction with nonzero kernel component has order $N$, at a fixed
measurement.  For $H=i[G,E]$ the first alternative is exactly
$G\in\mathcal W_E$ and $H\ne0$.
\end{corollary}
\begin{proof}
Apply the finite-dimensional spectral theorem to
$N(\mathcal C_E+N^{-1}\operatorname{Id})^{-1}$ and use
Proposition~\ref{prop:unitaryrange84}.
\end{proof}
The corollary concerns a comparison form, not yet an operational
lower bound.  The distinction matters at the boundary; the next
section proves the corresponding finite-pair statement on entire
unitary orbits.

\begin{proposition}[Exact represented-input decision]\label{prop:supportexact84}
On Gaussian-rational effects and a Gaussian-rational Hermitian
generator, membership in $\mathcal W_E$, stationarity of the orbit,
and the kernel dimension \eqref{eq:kerneldimension84} are decidable
by exact rational linear algebra in polynomial bit complexity.
The same calculation gives the orthogonal projection of $G$ onto
$\mathcal W_E$.
\end{proposition}
\begin{proof}
First use exact positive-semidefinite and normalization tests to
validate the input.  Select independent columns $V_j$ of $E_j$ by
rational elimination.  For a nonzero effect its support is
\[
 P_j=V_j(V_j^*V_j)^{-1}V_j^*;
\]
for a zero effect use $P_j=0$.  Compress a rational Hermitian basis
by all $P_j$ and select a real-linearly independent subfamily
$W_1,\ldots,W_q$.  The positive Gram system
$\sum_b\tr(W_aW_b)c_b=\tr(W_aG)$ gives
$\Pi_{\mathcal W_E}G=\sum_bc_bW_b$.  This tests membership exactly.
Commutators test stationarity.  Formula~\eqref{eq:kerneldimension84}
then gives the dimension.  There are at most $kd^2$ candidate
matrices and $d^2$ independent matrices.  Clearing denominators
and applying fraction-free elimination gives polynomial bit
bounds by the determinant bounds used in
Proposition~\ref{prop:covarianceexact83}.  This is a decision about
supports and a specified tangent; no optimum over testers or
quantum-control synthesis is computed.
\end{proof}

\section{Finite-angle distinguishability on unitary orbits}\label{sec:orbit84}
Fix an ordered measurement $E$ and a Hermitian generator $G$, and write
\begin{equation}\label{eq:unitorbit84}
 E_j(t)=e^{itG}E_je^{-itG},\qquad H_j=i[G,E_j].
\end{equation}
Both $E$ and $G$ are known in this pair-comparison problem.  A tester
must use the same controls under the two hypotheses $E$ and $E(t)$;
its design may depend on this known pair.  It has access to the
classical device labels and its own reference and memory, never to
an unobserved measurement environment.

The Hamiltonian-not-in-Kraus-span criterion and its quantum-error-
correction interpretation are established in quantum channel metrology
\cite[Theorems~1--3]{ZJ84}.  The support criterion below is its
measurement-channel specialization.  We prove a finite-angle,
finite-use trace-distance statement directly, including an explicit
accumulated error estimate.  Thus no lower bound is inferred solely
from an infinitesimal Fisher-information scaling.

\begin{theorem}[Finite-use trichotomy on every fixed unitary orbit]\label{thm:orbittrichotomy84}
Let $d,N\ge1$, $k\ge2$, and let $E$ and $G$ be as in
\eqref{eq:unitorbit84}, with no rank or strict-positivity assumption.
The following alternatives are exhaustive.
\begin{enumerate}
\item If $[G,E_j]=0$ for every $j$, then $D_N(E,E(t))=0$ for all
$N$ and all $t$.
\item If $H\ne0$ and $G\in\mathcal W_E$, there are constants
$c,C,t_0>0$, depending on $E,G$, such that
\begin{equation}\label{eq:orbitsql84}
 c\min\{1,\sqrt N|t|\}\le D_N(E,E(t))
 \le C\min\{1,\sqrt N|t|\},\qquad |t|\le t_0.
\end{equation}
The lower bound is attained by a product-input, nonadaptive experiment.
\item If $G\notin\mathcal W_E$, there are constants
$c,C,t_0>0$, depending on $E,G$, such that
\begin{equation}\label{eq:orbithl84}
 c\min\{1,N|t|\}\le D_N(E,E(t))
 \le C\min\{1,N|t|\},\qquad |t|\le t_0.
\end{equation}
The lower bound admits a common adaptive experiment carrying a logical
qubit between calls, with a reference of dimension at most $2d$ at
an individual call and ideal trusted controls.
\end{enumerate}
The constants are not uniform as the measurement or generator changes.
All upper bounds allow arbitrary retained reference and adaptive memory,
and bounded public stopping.  This theorem classifies a fixed unitary
orbit; it does not assert two-sided equivalence of the midpoint
covariance modulus on arbitrary pairs in the complete measurement body.
\end{theorem}

We first record the error-correction statement in the form needed for
finite angles.  Its exact-correction premise is the usual
Knill--Laflamme condition \cite{KL84}; a proof is included to specify
which systems are available to the recovery.

\begin{lemma}[A corrected logical channel and its quadratic error]\label{lem:correctedchannel84}
Let $\mathcal M$ be a channel on an input space $\mathsf A$, and
let $J:\C^2\longrightarrow\mathsf A\otimes\mathsf R$ be an isometry.
Suppose a channel $\mathcal R$ from the actual output of
$\mathcal M\otimes\operatorname{Id}_{\mathsf R}$ to $\C^2$ satisfies
\[
 \mathcal R\circ(\mathcal M\otimes\operatorname{Id}_{\mathsf R})
          \circ\mathcal J=\operatorname{Id}_2,
 \qquad \mathcal J(X)=JXJ^*.
\]
For a Hermitian $G$ on $\mathsf A$, put
$G_L=J^*(G\otimes I)J$, $\mathcal U_t(X)=e^{-itG_L}Xe^{itG_L}$,
and
\[
 \Phi_t=\mathcal R\circ(\mathcal M\otimes\operatorname{Id}_{\mathsf R})
       \circ\operatorname{Ad}(e^{-it(G\otimes I)})\circ\mathcal J.
\]
In unhalved diamond norm,
\begin{equation}\label{eq:logicalerror84}
 \|\Phi_t-\mathcal U_t\|_\diamond
 \le4t^2\|G\|_{\mathrm{op}}^2,
 \qquad
 \|\Phi_t^{\circ m}-\mathcal U_t^{\circ m}\|_\diamond
 \le4mt^2\|G\|_{\mathrm{op}}^2.
\end{equation}
\end{lemma}
\begin{proof}
Let $\mathcal T=\mathcal R\circ(\mathcal M\otimes\operatorname{Id})$
and choose Kraus operators $V_\ell$ for $\mathcal T$.  Since
$\mathcal T\circ\mathcal J$ is the identity channel, whose Choi
operator has rank one, $V_\ell J=c_\ell I_2$ with
$\sum_\ell|c_\ell|^2=1$.  Trace preservation on the entire input
space gives
\[
 \sum_\ell c_\ell^*V_\ell
 =\sum_\ell(V_\ell J)^*V_\ell
 =J^*\sum_\ell V_\ell^*V_\ell=J^*.
\]
It follows by differentiating that $\Phi_0=\operatorname{Id}_2$
and $\Phi'_0=-i[G_L,\,\cdot\,]$.
For any Hermitian $B$, the generator of $\operatorname{Ad}(e^{-itB})$
has diamond norm at most $2\|B\|_{\mathrm{op}}$.  Its second
 derivative has diamond norm at most $4\|B\|_{\mathrm{op}}^2$,
since unitary channels have diamond norm one.  Taylor's integral
remainder and contractivity under the surrounding channels bound
the remainder for $\Phi_t$ by $2t^2\|G\|_{\mathrm{op}}^2$.
The remainder for $\mathcal U_t$ has the same bound because
$\|G_L\|_{\mathrm{op}}\le\|G\|_{\mathrm{op}}$.  Their zeroth
and first derivatives agree, proving the first inequality.
Telescoping the two $m$-fold compositions of channels proves the second.
\end{proof}

\begin{lemma}[A reference-preserving code from the support complement]\label{lem:supportcode84}
Suppose $G\notin\mathcal W_E$ and put
\begin{equation}\label{eq:signstates84}
 A=G-\Pi_{\mathcal W_E}G,\qquad
 s=\tr A_+=\tr A_->0,\qquad \rho_\pm=A_\pm/s.
\end{equation}
There is a code $J:\C^2\to\C^d\otimes\mathsf R$ with
$\dim\mathsf R\le2d$, corrected exactly by a channel acting only
on the classical label and $\mathsf R$ after a call to $E$.
The logical generator $G_L=J^*(G\otimes I)J$ is diagonal with gap
\begin{equation}\label{eq:logicalgap84}
 \Delta=\tr\bigl(G(\rho_+-\rho_-)\bigr)
        =\|A\|_{\mathrm{HS}}^2/s>0.
\end{equation}
\end{lemma}
\begin{proof}
The identity belongs to $\mathcal W_E$ because $\sum_jE_j=I$;
thus $A$ is traceless.  Its nonzero positive and negative parts have
equal trace.  Also
\[
 P_j\rho_+P_j=P_j\rho_-P_j\quad\text{for every }j.
\]
Purify $\rho_+$ and $\rho_-$ in $\C^d\otimes\C^d$ and attach
orthogonal flags to the second factor.  Map the two logical basis
vectors to these flagged purifications.  They are orthonormal,
and the flags make all off-diagonal code matrix elements of
$B\otimes I$ vanish.  The two diagonal elements are equal for
every $B\in P_jM_d(\C)P_j$ by the compression identity above.

Choose measurement Kraus operators
$L_{j,a}=|j\rangle\langle a|\sqrt{E_j}$, $1\le a\le d$.
Products with different labels are zero; for one label their span
is $P_jM_d(\C)P_j$.  Hence, with $\mu=(j,a)$,
\begin{equation}\label{eq:kl84}
 J^*(L_\mu^*L_\nu\otimes I)J=\lambda_{\mu\nu}I_2
\end{equation}
for a positive matrix $\lambda$ of trace one.  Diagonalize
$\lambda$ by a unitary change of Kraus representation.  If its
positive eigenvalues are $\lambda_b$, the corresponding operators
$(\widetilde L_b\otimes I)J/\sqrt{\lambda_b}$ are isometries
with mutually orthogonal ranges.  A channel projects onto those
ranges and applies their inverse isometries; on the orthogonal
complement it prepares a fixed logical state.  This is completely
positive and trace preserving on the entire output space, and it
corrects the code exactly.  Its input space is just the actual
classical-label register tensor $\mathsf R$.  It does not contain
the lost input or a dilation environment.  Equivalently, since the
actual state is label diagonal, the same channel may first read the
label and then act on the retained reference.

The logical generator is diagonal by the flags, with entries
$\tr(G\rho_\pm)$.  Orthogonality of $A$ to
$\Pi_{\mathcal W_E}G$ gives
$\tr(GA)=\tr(A^2)$, proving \eqref{eq:logicalgap84}.
\end{proof}

\begin{proof}[Proof of Theorem~\ref{thm:orbittrichotomy84}]
The first assertion follows directly from \eqref{eq:unitorbit84}.
For the second, Proposition~\ref{prop:unitaryrange84} gives
\[
 h^2=\langle H,\mathcal C_E^\dagger H\rangle<\infty.
\]
Conjugating all tuple components by $e^{itG}$ intertwines the
covariances and their tangent directions, so this value is constant
along the orbit.  Inverse spectral comparison gives
$g_{N,E(t)}(E'(t))\le\sqrt N h$.  The complete-body path theorem
therefore yields
\begin{equation}\label{eq:orbitupper84}
 D_N(E,E(t))\le\min\{2,\sqrt{k+1}\,h\sqrt N|t|\}.
\end{equation}
There is an index $j$ with $H_j\ne0$.  Choose a unit vector $v$
with $\langle v,H_jv\rangle\ne0$.  The two Bernoulli parameters
obtained by input $v$ and the event ``label $j$'' have gap at least
$c_0|t|$ on a sufficiently small two-sided interval.  The finite
Bernoulli-product bound used in
Theorem~\ref{thm:multioutcomemetric82} gives
$(1/128)\min\{1,c_0\sqrt N|t|\}$.  This proves
\eqref{eq:orbitsql84}, including all finite $N$.

For the third alternative, use Lemma~\ref{lem:supportcode84}.
The equality
$\mathcal M_{E(t)}=\mathcal M_E\circ\operatorname{Ad}(e^{-itG})$
holds also with a retained reference.  Encode, make one device call,
recover, and repeat $m\le N$ times.  Under $E$ the logical channel
is the identity.  Under $E(t)$ it is $\Phi_t^{\circ m}$ from
Lemma~\ref{lem:correctedchannel84}.  Start in the equal superposition
of the two logical basis states.  The ideal logical unitary produces
unhalved trace distance $2|\sin(mt\Delta/2)|$ from that state.
Consequently, for every real $t$,
\begin{equation}\label{eq:finiteanglelower84}
 D_N(E,E(t))\ge
 \max_{1\le m\le N}
 \left(2|\sin(mt\Delta/2)|-4mt^2\|G\|_{\mathrm{op}}^2\right)_+.
\end{equation}
This bound concerns one specified common tester for each $m$,
not an inaccessible choice of a different recovery under the two hypotheses.

For $t\ne0$ choose
\[
 |t|\le t_0:=\min\left\{\frac1{2\Delta},
                       \frac{\Delta}{4\pi\|G\|_{\mathrm{op}}^2}\right\},
 \qquad
 m=\min\left\{N,\left\lfloor\frac1{|t|\Delta}\right\rfloor\right\}.
\]
Then $x=m|t|\Delta\le1$ and
$x\ge\tfrac12\min\{1,N|t|\Delta\}$.
Since $2\sin(x/2)\ge2x/\pi$ and the error in
\eqref{eq:finiteanglelower84} is at most $x/\pi$, we obtain
\begin{equation}\label{eq:explicitlocalhl84}
 D_N(E,E(t))\ge\frac1{2\pi}\min\{1,N|t|\Delta\}.
\end{equation}
The elementary hybrid estimate for the input unitaries gives
$D_N(E,E(t))\le\min\{2,2N|t|\|G\|_{\mathrm{op}}\}$,
which proves \eqref{eq:orbithl84}.  The zero-angle case is immediate.
The protocol uses deterministic $m\le N$ calls, hence belongs to
the class allowing bounded public stopping.  The upper estimates
hold for that entire class.  Finally if the first alternative held
with $G\notin\mathcal W_E$, then $H=0$ would contradict
Proposition~\ref{prop:unitaryrange84}; the alternatives are disjoint.
\end{proof}

\begin{corollary}[Full-rank, rank-one and projective measurements]\label{cor:orbitexamples84}
If one effect is positive definite, every nonconstant unitary orbit
has the square-root law \eqref{eq:orbitsql84}.  If all nonzero
effects have rank one, every nonconstant orbit has that law precisely
when the measurement is informationally complete, meaning that its
effects span $\mathcal H_d$.  If such a rank-one measurement is not
informationally complete, some generator has the linear law
\eqref{eq:orbithl84}.  For a projection-valued measurement every
orbit is either constant or has the linear law.
\end{corollary}
\begin{proof}
A full-rank support contributes $\mathcal H_d$ to $\mathcal W_E$.
For rank-one supports, $P_j\mathcal H_dP_j=\R E_j$, so that
$\mathcal W_E$ is exactly the effect span.  If this span is proper,
a nonzero element of its orthogonal complement is a generator
outside $\mathcal W_E$ and cannot commute with all effects by
Proposition~\ref{prop:unitaryrange84}.  For a projection-valued
measurement $\mathcal W_E$ is the block-diagonal Hermitian space
of the orthogonal supports, which is exactly the space of generators
commuting with every effect.  Apply Theorem~\ref{thm:orbittrichotomy84}.
\end{proof}

\begin{proposition}[Support spans under output processing]\label{prop:supportprocessing84}
Let $T_{aj}\ge0$ satisfy $\sum_aT_{aj}=1$, and put
$F_a=\sum_jT_{aj}E_j$.  Then $\mathcal W_E\subseteq\mathcal W_F$.
Thus classical processing of the ordered output cannot turn a
nonconstant square-root orbit into a linear orbit.  A reversible
splitting of labels preserves the support span and the classification.
\end{proposition}
\begin{proof}
For positive matrices, the kernel of a positive weighted sum is the
intersection of the kernels of its positively weighted summands.
Thus the support of $F_a$ contains $P_j$ whenever $T_{aj}>0$.
Every column has such a row, giving
$P_j\mathcal H_dP_j\subseteq\mathcal W_F$.
This proves the inclusion.  The orbit commutes with classical
processing, so Theorem~\ref{thm:orbittrichotomy84} applies with the
same generator.  A reverse stochastic map gives the opposite
inclusion for a reversible splitting.
\end{proof}

\begin{remark}[What the lower construction does not assert]\label{rem:orbitresources84}
The code and recovery depend on the known pair through $E,G$.
They do not constitute a common estimator of an unknown measurement.
Their existence is proved in finite-dimensional quantum mechanics;
polynomial gate synthesis, physical execution, and parameter-uniform
bounds for approximate controls are not inferred.  Formula
\eqref{eq:finiteanglelower84} quantifies the finite-angle channel
remainder under ideal controls.  It is not a claim of perfect
finite-use discrimination.  The full-body entropy and growing-$k$
learning questions are separate from this orbit classification.
\end{remark}

\section{Finite-use classification of smooth boundary curves}
\label{sec:curves85}
The support criterion extends beyond unitary conjugation.  Here
$E(t)=(E_1(t),\ldots,E_k(t))$, $|t|<a$, is a fixed two-sided
$C^2$ curve of ordered measurements, with $E=E(0)$ and
$H=E'(0)\ne0$.  The curve may change ranks at zero and need not
have a differentiable exact Kraus representation there.  All constants
below belong to this fixed curve, not uniformly to the measurement body.

Fix $0<a_0<a$.  Whenever a remainder below uses $L$, it means
\begin{equation}\label{eq:curvature86}
 L=\sup_{|u|\le a_0}\sum_{j=1}^k\|E_j''(u)\|_{\mathrm{op}}<\infty.
\end{equation}
This is the supremum of the sum, on the stated fixed interval.

\begin{lemma}[The support generator of a two-sided tangent]
\label{lem:curvegenerator85}
Let $P_j=\operatorname{supp}E_j$ and $Q_j=I-P_j$.  Then
$Q_jH_jQ_j=0$.  Conversely every Hermitian zero-sum tuple with these
vanishing blocks is the derivative of a real-analytic two-sided
measurement curve.  Define the Hermitian matrix
\begin{equation}\label{eq:curvegenerator85}
 \Gamma_E(H)=\frac i2\sum_j[P_j,H_j].
\end{equation}
On the real zero-sum tuple space,
\begin{equation}\label{eq:curverange85}
 H\in\operatorname{Ran}\mathcal C_E
 \quad\Longleftrightarrow\quad \Gamma_E(H)\in\mathcal W_E.
\end{equation}
For a unitary tangent $H_j=i[G,E_j]$,
$\Gamma_E(H)-G\in\mathcal W_E$.  In particular every stationary
unitary generator belongs to $\mathcal W_E$.
\end{lemma}
\begin{proof}
For $v\in\operatorname{Ran}Q_j$, the nonnegative function
$\langle v,E_j(t)v\rangle$ has a two-sided minimum at zero;
its derivative vanishes.  Polarization on that subspace gives
$Q_jH_jQ_j=0$.  The sum of the $H_j$ vanishes by normalization.
Pairing with the complete kernel tuple in
Theorem~\ref{thm:supportkernel84}, both the mean and the
missing-support blocks disappear, leaving
\[
 \langle H,K(A,Z)\rangle
 =\tr\left(A\sum_ji[H_j,P_j]\right)
 =-2\tr(A\Gamma_E(H)).
\]
Finite-dimensional self-adjointness proves \eqref{eq:curverange85}.
For a unitary tangent, expansion gives
\[
 \Gamma_E(H)=G-\tfrac12\sum_j(E_jGP_j+P_jGE_j),
\]
and the last sum belongs to $\mathcal W_E$.  If the tangent is zero,
this identity also proves the final assertion.  For the converse
realizability statement, the normalized-factor construction in the
proof of Lemma~\ref{lem:curvelogical85} is real analytic near zero
and has precisely the prescribed value and derivative.
\end{proof}

\begin{theorem}[Two-sided smooth-curve classification]\label{thm:curveclassification85}
For the fixed $C^2$ curve above there are $c,C,t_0>0$ such that,
for all integers $N\ge1$ and $|t|\le t_0$,
\begin{equation}\label{eq:curvedichotomy85}
 c\min\{1,N^\alpha|t|\}\le D_N(E,E(t))
       \le C\min\{1,N^\alpha|t|\},\qquad
 \alpha=
 \begin{cases}
  1/2,&\Gamma_E(H)\in\mathcal W_E,\\
  1,&\Gamma_E(H)\notin\mathcal W_E.
 \end{cases}
\end{equation}
The first lower bound uses repeated product inputs.  The second has a
known-pair adaptive realization with ideal encoding and recovery,
a logical qubit, and a per-call external reference of dimension at
most $2d$.  The same controls are used under both hypotheses.
Constants may degenerate with the tangent, nonzero support eigenvalues,
curvature, or support-complement gap.
\end{theorem}

The exponents and the error-correction mechanism have the established
channel-metrology antecedents \cite{ZJ84,KL84}.  The content here is a
direct finite-pair theorem from $C^2$ effect coordinates: the support
formula decides its branch even at second-order rank openings, without
assuming a smooth exact Kraus family or inferring a finite-angle lower
from Fisher information.  It is a theorem for each fixed nonzero
first-order curve, not a uniform arbitrary-pair midpoint equivalence.

\begin{lemma}[A horizontal curve with prescribed first derivative]
\label{lem:horizontalcurve85}
If $H\in\operatorname{Ran}\mathcal C_E$, there are factors
$A_j=\sqrt{E_j}$ and $B_j$ satisfying
\[
 A^*A=I,\quad A^*B=0,\quad
 A_j^*B_j+B_j^*A_j=H_j,
\]
where $A,B$ denote the vertically stacked matrices.  Put
$b=\|B\|_{\mathrm{op}}$ and $S=B^*B$.  The factors
\begin{equation}\label{eq:horizontalcurve85}
 A(t)=(A+tB)(I+t^2S)^{-1/2}
\end{equation}
define a normalized measurement $\widetilde E_j(t)=A_j(t)^*A_j(t)$
with the same zeroth and first derivatives as $E(t)$, and
\[
 A(t)^*A'(t)=0,\qquad
 \|A'(t)\|_{\mathrm{op}}\le b,\qquad
 \|A''(t)\|_{\mathrm{op}}\le7b^2\quad(|t|b\le1).
\]
If $L=\sup_{|s|\le a_0}\sum_j\|E_j''(s)\|_{\mathrm{op}}$ on a
closed subinterval of the given curve, then, for $|t|\le a_0$ and
$|t|b\le1$,
\begin{equation}\label{eq:curvesqlupper85}
 D_N(E,E(t))\le
 \min\left\{2,\;2b\sqrt N|t|+
             \tfrac12(L+16b^2)Nt^2\right\}.
\end{equation}
\end{lemma}
\begin{proof}
The factorization
$\mathcal C_E=T_A(I-AA^*)T_A^*/4$ in
Lemma~\ref{lem:covarianceidentities83} gives a horizontal solution:
for $K=\mathcal C_E^\dagger H$ take
$B=(I-AA^*)T_A^*K/4$.  Thus $T_AB=H$ and $A^*B=0$.
In fact $4\|B\|_{\mathrm{HS}}^2=
\langle H,\mathcal C_E^\dagger H\rangle$.
Formula \eqref{eq:horizontalcurve85} is normalized because
$(A+tB)^*(A+tB)=I+t^2S$.  Its derivative is
\[
 A'(t)=(B-tAS)(I+t^2S)^{-3/2},\qquad
 A'(t)^*A'(t)=S(I+t^2S)^{-2}.
\]
These identities prove horizontality and the first derivative bound.
Differentiation once more gives, in the displayed order,
\[
 A''(t)=-AS(I+t^2S)^{-3/2}
       -3t(B-tAS)S(I+t^2S)^{-5/2}.
\]
Since $\|A\|=1$, $\|S\|=b^2$, and $|t|b\le1$, its norm is at
most $b^2+3|t|(b+|t|b^2)b^2\le7b^2$.

Duplicate each output label into an inaccessible environment to form
a Stinespring isometry $W(t)$ from $A(t)$.  It has exactly the same
operator derivative bounds and $W(t)^*W'(t)=0$.
For one fixed purified adaptive tester, differentiating its final
vector gives a sum over slots.  Distinct terms are orthogonal: after
cancelling later common isometries, the later differentiated slot
contains $W^*W'$ or its adjoint.  Each diagonal term has squared
norm at most $b^2$.  The final vector derivative is at most
$\sqrt N b$, and the density derivative at most $2\sqrt N b$.
Integration followed by common partial trace proves
$D_N(E,\widetilde E(t))\le2b\sqrt N|t|$.
Padding stopped branches by discarded fixed-input calls includes
bounded public stopping; the original stopped output is a common
postprocessing of the padded experiment.

In the following application $\sum_jR_j=0$.  Thus the map is a
Hermiticity-preserving, trace-annihilating channel difference or
derivative, not a positive channel.  Lemma~\ref{lem:channelnorm86}
also proves its bound on arbitrary amplified trace-class inputs by
adjoint duality, without a state-optimization assumption.

For a Hermitian tuple $R$, the measurement derivative has diamond norm
at most $\sum_j\|R_j\|_{\mathrm{op}}$.
Indeed each reference output block is the partial trace against $R_j$;
its trace norm on a state is at most $\|R_j\|_{\mathrm{op}}$ by
duality.  The Hermiticity-preserving diamond norm may be optimized
over states with a reference.  The Stinespring second derivative
bound similarly gives
$\|\mathcal M_{\widetilde E(t)}''\|_\diamond
\le2\|W''(t)\|+2\|W'(t)\|^2\le16b^2$.
Taylor's integral formula, with common value and derivative at zero,
therefore bounds the one-call channel difference by
$(L+16b^2)t^2/2$.  Hybrid telescoping over $N$ calls proves
\eqref{eq:curvesqlupper85}.
\end{proof}

\begin{lemma}[A logical derivative for a general measurement curve]
\label{lem:curvelogical85}
Suppose $\Gamma=\Gamma_E(H)\notin\mathcal W_E$.
Use the support-complement code of Lemma~\ref{lem:supportcode84}
with $\Gamma$ in place of $G$, and let $\mathcal R$ be its fixed
CPTP recovery.  Put
\[
 \Phi_t=\mathcal R\circ(\mathcal M_{E(t)}\otimes\operatorname{Id})
             \circ\mathcal J,\qquad
 \Gamma_L=J^*(\Gamma\otimes I)J.
\]
Here $L$ is exactly the local curvature quantity
$\sup_{|u|\le a_0}\sum_j\|E_j''(u)\|_{\mathrm{op}}$ from
\eqref{eq:curvature86}; the remainder does not infer it from $E,H$.
Then $\Phi_0=\operatorname{Id}_2$, $\Phi'_0=-i[\Gamma_L,\cdot]$,
and, for $|t|\le a_0$,
\begin{equation}\label{eq:curvelogical85}
 \|\Phi_t-\operatorname{Ad}(e^{-it\Gamma_L})\|_\diamond
 \le\kappa t^2,\qquad
 \kappa=\tfrac12L+2\|\Gamma\|_{\mathrm{op}}^2.
\end{equation}
Its logical gap is
$\Delta=\|Z\|_{\mathrm{HS}}^2/\tr Z_+>0$,
where $Z=\Gamma-\Pi_{\mathcal W_E}\Gamma$.
\end{lemma}
\begin{proof}
To compute the derivative, not to factor the unknown curve itself,
put
\[
 A_j=\sqrt{E_j},\qquad
 B_j=(\sqrt{E_j})^+H_j(I-P_j/2).
\]
The inverse here is only on the fixed support; it is zero on its
complement.  Since $Q_jH_jQ_j=0$,
$A_j^*B_j+B_j^*A_j=H_j$ and
\[
 \sum_j A_j^*B_j=\tfrac12\sum_j[P_j,H_j]=-i\Gamma.
\]
In particular this sum is anti-Hermitian.  The normalized factors
$(A+tB)(I+t^2B^*B)^{-1/2}$ have derivative $B$ and induce the
same channel value and derivative as $E(t)$, irrespective of its
second-order rank changes.

Write $L_\mu=|j\rangle\langle a|A_j$ and
$\dot L_\mu=|j\rangle\langle a|B_j$ for $\mu=(j,a)$,
with identity on the reference understood.  If $R_\ell$ are Kraus
operators of the completed recovery, exact correction gives
$R_\ell L_\mu J=c_{\ell\mu}I_2$.
The left first-order term of the logical channel is therefore
\[
 \sum_{\ell,\mu}\overline{c_{\ell\mu}}
                  R_\ell\dot L_\mu J
 =J^*\sum_\mu L_\mu^*\dot L_\mu J
 =-i\Gamma_L,
\]
where completeness of the recovery is used in the first equality.
The adjoint term proves the commutator derivative.  The actual
logical second derivative is at most $L$ in diamond norm, and that
of the logical unitary is at most $4\|\Gamma\|_{\mathrm{op}}^2$.
Taylor's integral remainder proves \eqref{eq:curvelogical85}.
The support-code lemma gives the gap and the available systems.
\end{proof}

\begin{proof}[Proof of Theorem~\ref{thm:curveclassification85}]
If $H$ is in the covariance range, let $x=\sqrt N|t|$.
For $x\le1$, \eqref{eq:curvesqlupper85} is at most
$[2b+(L+16b^2)/2]x$; for $x\ge1$ use the bound two.
For the lower, choose $j,v$ with $\langle v,H_jv\rangle\ne0$.
The fixed input $v$ and the event ``label $j$'' have a first-order
probability gap, and Lemma~\ref{lem:bernoulli85} proves the result.

If $H$ is outside the range, apply Lemma~\ref{lem:curvelogical85}.
Starting from the equal logical superposition and telescoping $m$
corrected calls gives, for $|t|\le a_0$,
\begin{equation}\label{eq:curvefinitelower85}
 D_N(E,E(t))\ge
 \max_{1\le m\le N}\bigl(2|\sin(mt\Delta/2)|-m\kappa t^2\bigr)_+.
\end{equation}
Choose
\[
 t_0\le\min\{a_0,(2\Delta)^{-1},\Delta/(2\pi\kappa)\},\qquad
 m=\min\{N,\lfloor(|t|\Delta)^{-1}\rfloor\}.
\]
For $0<|t|\le t_0$ the floor is at least two, and
$x=m|t|\Delta$ satisfies
$\tfrac12\min\{1,N|t|\Delta\}\le x\le1$.
The sine term is at least $2x/\pi$ and the error at most $x/(2\pi)$.
This gives a positive constant times $\min\{1,N|t|\}$.
For the upper, the finite constant
$L_1=\sup_{|s|\le a_0}\sum_j\|E_j'(s)\|_{\mathrm{op}}$
and hybrid telescoping give $D_N\le\min\{2,NL_1|t|\}$.
The point $t=0$ is immediate.
\end{proof}

\begin{corollary}[Rank opening without a unitary orbit]\label{cor:rankopening85}
Let $P=|0\rangle\langle0|$ on $\C^2$,
$G=\left(\begin{smallmatrix}0&-i\\i&0\end{smallmatrix}\right)$,
$U_t=e^{itG}$, and for $|t|<1/2$ set
\[
 E_1(t)=(1-t^2)U_tPU_t^*+\tfrac12t^2I_2,
 \qquad E_2(t)=I_2-E_1(t).
\]
Both effects are full rank for $t\ne0$ and projective at zero.
Nevertheless $D_N(E(0),E(t))\asymp\min\{1,N|t|\}$ on a fixed
small interval, uniformly over $N\ge1$.
\end{corollary}
\begin{proof}
At zero the derivative is $i[G,(P,I-P)]\ne0$ and the support span
is diagonal, while $G$ is not.  Lemma~\ref{lem:curvegenerator85}
and Theorem~\ref{thm:curveclassification85} apply.  The eigenvalues
of the first effect are $t^2/2$ and $1-t^2/2$, proving the rank assertion.
\end{proof}
A vanishing first derivative is intentionally excluded from the
classification: it does not imply stationarity of a general curve.
For example the scalar curve $(t^2,1-t^2)$ has distance
$2[1-(1-t^2)^N]$ from $(0,1)$, of order $\min\{1,Nt^2\}$.
Likewise a one-sided rank opening can have a nonzero missing-support
first-order block and is not covered by the two-sided tangent lemma.

\section{Certified errors in the discrimination controls}
\label{sec:controlstability85}
The preceding lower bounds first concern ideal known-pair controls.
Their stability has a direct finite formulation.  An error bound on a
control channel is a bound in unhalved diamond norm, valid on every
input including its memory reference.  It is not a bound measured only
on one calibration state.

\begin{theorem}[A finite control-error ledger]\label{thm:controlstability85}
Use the code and recovery from Lemma~\ref{lem:curvelogical85}, with
gap $\Delta$ and remainder constant $\kappa$.
In this statement
$L=\sup_{|u|\le a_0}\sum_j\|E_j''(u)\|_{\mathrm{op}}$ and
$\kappa=L/2+2\|\Gamma_E(H)\|_{\mathrm{op}}^2$, as in
\eqref{eq:curvature86} and \eqref{eq:curvelogical85}.
Suppose an implemented initial logical state differs from the ideal
one in trace norm by at most $e_0$.  At cycle $r$, suppose its
implemented logical channel differs from the corresponding ideal
corrected channel by at most $\nu_r$ under each hypothesis.
Suppose the implemented final readout differs from the ideal
Helstrom readout by at most $e_f$ in diamond norm.
For every integer $m\le N$ and $|t|\le a_0$, the actual classical
output laws obey
\begin{equation}\label{eq:controlledger85}
 \|\widetilde p_{0,m}-\widetilde p_{t,m}\|_1
 \ge\left(2|\sin(mt\Delta/2)|-m\kappa t^2
              -2e_0-2e_f-2\sum_{r=1}^m\nu_r\right)_+.
\end{equation}
Here one may choose the Helstrom readout for the ideal unrotated and
logical-unitary states; it depends on the known pair and on $m$.
The cycle error may be bounded by the sum of the encoding and recovery
errors, because the unknown measurement itself is a channel.

In particular, put $z=\min\{1,N|t|\Delta\}$ and choose $m,t_0$
as in the proof of Theorem~\ref{thm:curveclassification85}.
For $0<|t|\le t_0$, the sufficient conditions
\begin{equation}\label{eq:controlprecision85}
 e_0+e_f\le z/(16\pi),\qquad
 \nu_r\le |t|\Delta/(8\pi)\quad(1\le r\le m)
\end{equation}
imply
$\|\widetilde p_{0,m}-\widetilde p_{t,m}\|_1\ge z/(2\pi)$.
All errors are deterministic uniform channel bounds; no independence
of control errors is required.
\end{theorem}
\begin{proof}
Choose the ideal two-outcome Helstrom readout for the equal logical
superposition and its image under $e^{-imt\Gamma_L}$.
Its classical $\ell^1$ separation is $2|\sin(mt\Delta/2)|$.
Replacing the logical unitary by the actual ideal corrected channel
loses at most $m\kappa t^2$, by channel telescoping and
Lemma~\ref{lem:curvelogical85}.  Under either hypothesis,
telescoping the implemented preparation, cycles and readout changes
its output law in $\ell^1$ by at most
$e_0+e_f+\sum_r\nu_r$.  The triangle inequality proves
\eqref{eq:controlledger85}.  Replacing encoding and recovery around
one unchanged unknown channel costs no more than the sum of their
diamond errors, which proves the stated decomposition of $\nu_r$.

For the indicated choice set $x=m|t|\Delta$, so $z/2\le x\le1$.
The ideal sine term minus curvature is at least $3x/(2\pi)$.
Conditions \eqref{eq:controlprecision85} give
$e_0+e_f\le x/(8\pi)$ and $\sum_r\nu_r\le x/(8\pi)$.
The total loss is therefore at most $x/(2\pi)$, leaving
$x/\pi\ge z/(2\pi)$.
\end{proof}

For a unitary orbit, Lemma~\ref{lem:correctedchannel84} gives the
same result with $\kappa=4\|G\|_{\mathrm{op}}^2$ and the original
support-code gap.  An error allowance fixed independently of the
angle need not meet \eqref{eq:controlprecision85} as $t\to0$.
The theorem does not certify a particular piece of hardware, synthesize
a gate circuit, or allow target-dependent advice in common learning.
It specifies exactly which certified approximation errors suffice
to preserve the finite-use lower bound.

\begin{proposition}[Exact first-order decisions]\label{prop:curveexact85}
For Gaussian-rational $E$ and a Gaussian-rational Hermitian zero-sum
$H$, two-sided first-order realizability and the branch of
Theorem~\ref{thm:curveclassification85} for $H\ne0$ are decidable
in polynomial represented-input bit complexity.  The decision gives
$\Gamma_E(H)$, its projection on $\mathcal W_E$, and its orthogonal
residual.  A zero tangent is reported as a higher-order question,
not as stationarity of an unspecified curve.
\end{proposition}
\begin{proof}
The exact support projectors in Proposition~\ref{prop:supportexact84}
are rational.  Test $Q_jH_jQ_j=0$ for every $j$ and compute
\eqref{eq:curvegenerator85}.  Necessity and sufficiency of the test
follow from Lemma~\ref{lem:curvegenerator85} and the normalized-factor
construction in Lemma~\ref{lem:curvelogical85}: for any such tangent
it gives a real-analytic measurement curve with that derivative.
Real coordinate rank selection and the Hilbert--Schmidt Gram system
then project $\Gamma_E(H)$ exactly.  Rational matrix arithmetic and
fraction-free elimination have the bit bounds already proved in
Proposition~\ref{prop:supportexact84}.  Only these finite operations
are performed.  A numerical curvature bound for an unspecified
curve, an optimal tester, and a compiled recovery are not outputs
of this decision.
\end{proof}

\section{One-sided tangent neighborhoods}
\label{sec:tangentneighborhood86}
The preceding curve results admit a finite-pair formulation.  The
base measurement and its direction are fixed, but the second-order
perturbation may vary arbitrarily within a prescribed bound.  This
also permits first-order opening of a missing support.  Write
$\|R\|_{\Sigma}=\sum_j\|R_j\|_{\mathrm{op}}$ for a Hermitian tuple.
All trace and diamond norms in this section are unhalved.

\begin{lemma}[A reference-uniform channel bound]\label{lem:channelnorm86}
For a Hermitian tuple $R$, let
$\mathcal D_R(X)=\sum_j\tr(R_jX)|j\rangle\langle j|$.
Then
\begin{equation}\label{eq:channelnorm86}
 \|\mathcal D_R\|_\diamond\le\|R\|_{\Sigma}.
\end{equation}
The map is Hermiticity preserving, and is trace annihilating when
$\sum_jR_j=0$.  In particular this applies to measurement-channel
differences, derivatives, and Taylor remainders.
\end{lemma}
\begin{proof}
At every reference dimension the output is block diagonal.  In trace
norm duality it therefore suffices to use a block-diagonal test
$Y=\bigoplus_jY_j$, with $\|Y_j\|_{\mathrm{op}}\le1$.  Its adjoint
image is $\sum_jR_j\otimes Y_j$, of norm at most
$\sum_j\|R_j\|_{\mathrm{op}}$.  Duality bounds the amplified map on
\emph{every} trace-class input $X$, not only on states.  Taking the
supremum over reference dimensions proves \eqref{eq:channelnorm86}.
The remaining assertions follow by taking adjoints and traces.
\end{proof}

\begin{lemma}[The one-sided cone and its lineality]\label{lem:onesidedcone86}
Fix a measurement $E$, put $P_j=\operatorname{supp}E_j$ and $Q_j=I-P_j$,
and let $H$ be a Hermitian zero-sum tuple.  It is the right derivative
at zero of a one-sided real-analytic measurement curve if and only if
\begin{equation}\label{eq:onesidedcone86}
 J_j:=Q_jH_jQ_j\succeq0\quad\hbox{for every }j.
\end{equation}
The largest vector subspace of this cone is given by $J_j=0$ for all
$j$.  For every Hermitian zero-sum tuple, not necessarily in the cone,
\begin{equation}\label{eq:fullrange86}
 H\in\operatorname{Ran}\mathcal C_E
 \quad\Longleftrightarrow\quad
 J_j=0\ (1\le j\le k),\qquad
 \Gamma_E(H)\in\mathcal W_E.
\end{equation}
Here $\Gamma_E(H)=\frac i2\sum_j[P_j,H_j]$ and
$\mathcal W_E=\sum_jP_j\mathcal H_dP_j$ are real Hermitian objects.
\end{lemma}
\begin{proof}
Right differentiation of $\langle v,E_j(t)v\rangle\ge0$ for
$v\in\operatorname{Ran}Q_j$ proves necessity.  For sufficiency, for
each nonzero support let $\lambda_j>0$ be the smallest eigenvalue
of $E_j$ on that support.  Choose
\[
 c\ge1+\max_{j:P_j\ne0}2\|H_j\|_{\mathrm{op}}^2/\lambda_j.
\]
For all sufficiently small $s\ge0$, the $P_j$ block of
$E_j+sH_j+cs^2I$ is at least $\lambda_jP_j/2$.
Its Schur complement on $Q_j$ is bounded below by
\[
 sJ_j+s^2\left(c-2\|H_j\|_{\mathrm{op}}^2/\lambda_j\right)Q_j
 \succeq0.
\]
If $P_j=0$, condition \eqref{eq:onesidedcone86} is $H_j\succeq0$
and positivity is immediate.  Thus the rational matrix curve
\begin{equation}\label{eq:conecurve86}
 E_j^{\mathrm c}(s)=\frac{E_j+sH_j+cs^2I}{1+kcs^2}
\end{equation}
is positive, sums exactly to $I$, and has right derivative $H_j$.
This proves sufficiency, including singular $J_j$ with nonzero cross
blocks.  The cone and its negative intersect precisely at $J_j=0$.

Pair $H$ with the complete kernel in
Theorem~\ref{thm:supportkernel84}.  Taking only its independent
missing-support parameters $Z_j$ shows that orthogonality to the
kernel forces $Q_jH_jQ_j=0$; the tuple-mean term vanishes because
$\sum_jH_j=0$.  Once those blocks vanish, the remaining pairing is
$-2\tr(A\Gamma_E(H))$ for every $A\in\mathcal W_E^\perp$.
Self-adjointness of $\mathcal C_E$ proves \eqref{eq:fullrange86}.
\end{proof}

\begin{lemma}[A finite normalized-factor remainder]\label{lem:factorremainder86}
Let $A_j^*A_j=E_j$, $A^*A=I$, and
$A_j^*B_j+B_j^*A_j=H_j$, with $\sum_jH_j=0$.
Set $b=\|B\|_{\mathrm{op}}$, $h=\max_j\|H_j\|_{\mathrm{op}}$ and
\[
 R_s=(I+s^2B^*B)^{-1/2},\qquad
 G_j(s)=R_s(A_j+sB_j)^*(A_j+sB_j)R_s.
\]
This defines an analytic measurement near zero, and, for $0\le s\le1$,
\begin{equation}\label{eq:factorremainder86}
 \|G(s)-E-sH\|_{\Sigma}\le K_b s^2,
 \qquad K_b=kb^2(2+h).
\end{equation}
When $sb\le1$, its stacked factors differ from $A$ in operator norm
by at most $3bs/2$.  If $A^*B=0$, then also
\begin{equation}\label{eq:horizontalpair86}
 D_N(E,G(s))\le\min\{2,2b\sqrt N s\}.
\end{equation}
\end{lemma}
\begin{proof}
The zero tuple sum gives $A^*B+B^*A=0$, so the stacked normalization
is $I+s^2B^*B$.  This matrix is positive definite for all real $s$;
its inverse square root is analytic near zero.  Functional calculus
gives $\|R_s\|\le1$ and $\|R_s-I\|\le s^2b^2/2$.
Expanding the effect as
$R_s(E_j+sH_j+s^2B_j^*B_j)R_s$ bounds its remainder by
$s^2b^2(1+sh)+s^2\|B_j\|^2$.  Sum over $j$, use
$\sum_j\|B_j\|^2\le kb^2$, and obtain \eqref{eq:factorremainder86}.
The factor difference is at most $sb+s^2b^2/2$.
For horizontal $B$, the ordered identities in
Lemma~\ref{lem:horizontalcurve85} give $A(s)^*A'(s)=0$ and
$\|A'(s)\|\le b$.  Copying the classical label into a discarded
environment is an isometric embedding of the stacked factors, so
it preserves these derivative norms.  The orthogonal-slot argument
of that lemma gives \eqref{eq:horizontalpair86}, including bounded
public stopping.  The environment is not supplied to the tester.
\end{proof}

\begin{definition}[A quadratic tangent neighborhood]\label{def:tangenttube86}
Fix $E$, a Hermitian zero-sum $H\ne0$ satisfying
\eqref{eq:onesidedcone86}, and a number $\Lambda\ge0$.
For $s>0$, the neighborhood consists of all legal measurements $F$
with
\begin{equation}\label{eq:tangenttube86}
 \|F-E-sH\|_{\Sigma}\le\Lambda s^2.
\end{equation}
The number $\Lambda$ is a bound on a \emph{finite remainder}, not an
inferred curvature bound.  No curve, or differentiability of $F$ as
a function of $s$, is part of this definition.
\end{definition}

\begin{theorem}[Finite-pair classification in a fixed tangent neighborhood]
\label{thm:tubedichotomy86}
For each fixed $E,H,\Lambda$ in Definition~\ref{def:tangenttube86}
there are $c,C,s_0>0$ such that, for every integer $N\ge1$, every
$0<s\le s_0$, and every legal $F$ satisfying
\eqref{eq:tangenttube86},
\begin{equation}\label{eq:tubedichotomy86}
 c\min\{1,N^\alpha s\}\le D_N(E,F)
       \le C\min\{1,N^\alpha s\},\qquad
 \alpha=\begin{cases}
  1/2,&H\in\operatorname{Ran}\mathcal C_E,\\
  1,&H\notin\operatorname{Ran}\mathcal C_E.
 \end{cases}
\end{equation}
The constants are uniform over the stated remainder neighborhood,
but not over $E,H$ or $\Lambda$.  If some $J_j\ne0$, the linear
lower is attained by repeated product inputs and a classical event.
If all $J_j=0$ but $\Gamma_E(H)\notin\mathcal W_E$, it is attained
by the known-direction correction code on the actual label and a
retained reference, with ideal controls.  Neither construction is
a common unknown-measurement learner.
\end{theorem}

There are two different reasons for a linear scale.  First-order
support opening creates an event impossible at the base measurement.
Within the lineality space of the cone, coherent correction produces
the established Kraus-span metrological mechanism \cite{ZJ84,KL84}.
The result above makes this distinction in support coordinates and
controls finite remainders uniformly at a fixed base and direction.
It is not a uniform arbitrary-pair equivalence for the midpoint
covariance certificate.

\begin{proof}
Put $h=\max_j\|H_j\|_{\mathrm{op}}>0$ and
$h_1=\|H\|_{\Sigma}$.  Choose a unit eigenvector of one $H_j$
with expectation of magnitude $h$.  On that fixed input the event
``label $j$'' has probability gap at least $hs-\Lambda s^2$.
For $s$ sufficiently small this is at least $hs/2$.
Lemma~\ref{lem:bernoulli85} therefore gives a uniform product lower
of order $\min\{1,\sqrt N s\}$, for every $F$ in the neighborhood.
All pairs also satisfy the elementary upper
\begin{equation}\label{eq:tubehybrid86}
 D_N(E,F)\le\min\{2,N(h_1s+\Lambda s^2)\},
\end{equation}
by Lemma~\ref{lem:channelnorm86} and channel telescoping.

Suppose first that $H\in\operatorname{Ran}\mathcal C_E$.
The horizontal solution of Lemma~\ref{lem:horizontalcurve85}
provides factors $B$ with $A^*B=0$ and $T_AB=H$.
Lemma~\ref{lem:factorremainder86} and \eqref{eq:tangenttube86} give
\begin{equation}\label{eq:tubefiniteupper86}
 D_N(E,F)\le\min\{2,2b\sqrt N s+(\Lambda+K_b)Ns^2\}.
\end{equation}
For $x=\sqrt N s\le1$, the right side is at most
$(2b+\Lambda+K_b)x$; for $x\ge1$ it is at most two.
This proves the upper in the first branch without discarding its
finite $Ns^2$ remainder before comparison.

Next suppose some $J_j\ne0$.  It is positive semidefinite; choose a
unit vector $v\in\operatorname{Ran}Q_j$ with
$\lambda=\langle v,J_jv\rangle>0$.
The base probability of label $j$ is zero, while that under $F$ is
at least $\lambda s-\Lambda s^2\ge\lambda s/2$ after decreasing
$s_0$.  Repeating $v$ gives the exact impossible-event bound
\begin{equation}\label{eq:leakagelower86}
 D_N(E,F)\ge2\left[1-(1-\lambda s/2)^N\right]
 \ge2(1-e^{-1})\min\{1,N\lambda s/2\}.
\end{equation}
Here $s_0$ is chosen also so that $\lambda s/2\le1$.
Together with \eqref{eq:tubehybrid86} this proves the linear law.
The vector and the decision event are the same for every allowed $F$.

It remains to consider $J_j=0$ for all $j$ and
$\Gamma:=\Gamma_E(H)\notin\mathcal W_E$.
Use the completed recovery and logical code of
Lemma~\ref{lem:supportcode84}, applied to $\Gamma$.
The normalized first-order factors in Lemma~\ref{lem:curvelogical85}
show algebraically that the corrected derivative $\mathcal D_H$
is $-i[\Gamma_L,\cdot]$.  This calculation needs only $E,H$, not
a smooth factorization of $F$.
For $\Phi_F=\mathcal R\circ(\mathcal M_F\otimes\operatorname{Id})
\circ\mathcal J$, \eqref{eq:tangenttube86} gives
\[
 \|\Phi_F-\operatorname{Id}_2+is[\Gamma_L,\cdot]\|_\diamond
 \le\Lambda s^2.
\]
The integral Taylor remainder for the logical unitary is at most
$2\|\Gamma\|_{\mathrm{op}}^2s^2$.  Consequently
\begin{equation}\label{eq:tubelogical86}
 \|\Phi_F-\operatorname{Ad}(e^{-is\Gamma_L})\|_\diamond
 \le\kappa_{\mathrm{tube}}s^2,
 \qquad \kappa_{\mathrm{tube}}=\Lambda+2\|\Gamma\|_{\mathrm{op}}^2.
\end{equation}
The gap is $\Delta=\|Z\|_{\mathrm{HS}}^2/\tr Z_+>0$, where
$Z=\Gamma-\Pi_{\mathcal W_E}\Gamma$.
For every $m\le N$, equal logical amplitudes and channel telescoping
give a lower bound
$\bigl(2|\sin(ms\Delta/2)|-m\kappa_{\mathrm{tube}}s^2\bigr)_+$.
Choose
\[
 s_0\le\min\{1,(2\Delta)^{-1},
              \Delta/(2\pi\kappa_{\mathrm{tube}})\},\qquad
 m=\min\{N,\lfloor(s\Delta)^{-1}\rfloor\}.
\]
The floor is at least two, and
$\frac12\min\{1,Ns\Delta\}\le x=ms\Delta\le1$.
The sine term is at least $2x/\pi$ and the loss at most $x/(2\pi)$.
In particular the separation is at least
$\min\{1,Ns\Delta\}/(2\pi)$.
The readout may be fixed to the ideal logical pair, so this lower is
uniform over all the permitted $F$ as well.  Its code uses an external
reference of dimension at most $2d$ per call.  Finally
\eqref{eq:tubehybrid86} completes the proof.
\end{proof}

\begin{corollary}[All nonzero one-sided first jets]\label{cor:onesidedcurves86}
Let $E(t)$, $0\le t<a$, be a fixed $C^2$ measurement curve, including its
right derivatives at zero, with $H=E'(0)\ne0$.  Then
\eqref{eq:tubedichotomy86} holds with $s=t$ and curve-dependent
constants on a fixed right interval.  The branch is decided by
\eqref{eq:fullrange86}.  Thus first-order support openings, as well
as tangential second-order rank changes, are included.
\end{corollary}
\begin{proof}
One-sided positivity gives \eqref{eq:onesidedcone86}.
On $[0,a_0]$ set
$L_+=\sup_{0\le u\le a_0}\sum_j\|E_j''(u)\|_{\mathrm{op}}$.
Taylor's formula gives \eqref{eq:tangenttube86} with
$\Lambda=L_+/2$.  Apply Theorem~\ref{thm:tubedichotomy86}.
This does not extend the two-sided tangent lemma by dropping one of
its hypotheses; it uses the larger cone of Lemma~\ref{lem:onesidedcone86}.
\end{proof}

\begin{corollary}[Certified controls on the whole neighborhood]
\label{cor:tubecontrols86}
In the tangential linear branch, Theorem~\ref{thm:controlstability85}
holds uniformly over \eqref{eq:tangenttube86} with $t=s$ and
$\kappa=\kappa_{\mathrm{tube}}$ from \eqref{eq:tubelogical86}.
In particular the same initial, cycle and final error allowances
with $z=\min\{1,Ns\Delta\}$ preserve separation at least $z/(2\pi)$.
\end{corollary}
\begin{proof}
Equation~\eqref{eq:tubelogical86} supplies the only curvature input
used in that theorem's telescoping proof.  Pay it once, and pay the
certified implementation errors under both hypotheses exactly as
in \eqref{eq:controlledger85}.  The ideal readout depends on
$E,H,s,m$, not on the unknown remainder within the neighborhood.
The precision allowances remain angle dependent, not fixed-noise
robustness or a control-synthesis guarantee.
\end{proof}

\section{Independent probes and higher-order approaches}
\label{sec:independent86}
A linear finite-use scale need not signal coherent accumulation.
The support-opening case can already be distinguished by independent
classical observations.  We make the acquisition restriction precise.
Let $D_N^{\mathrm{prod}}(E,F)$ be the supremum of the final trace
separation when the input is a product of at most $N$ probe--reference
states, one pair per call.  The pairs may differ and may depend on the
known hypotheses.  All completed outputs and references may undergo
an arbitrary joint final readout.  There is no quantum correlation
between different input pairs, and no intervening feedback into later
inputs.  Public randomization is allowed.  This is narrower than
nonadaptive acquisition with an entangled many-call input.

\begin{theorem}[Three local acquisition regimes]\label{thm:producttrichotomy86}
Fix $E,H,\Lambda$ as in Definition~\ref{def:tangenttube86}.  Uniformly
in $N$, all sufficiently small $s>0$, and every $F$ in
\eqref{eq:tangenttube86}, the two distances have the following orders:
\begin{center}
\begin{tabular}{@{}lll@{}}
\toprule
Direction & $D_N^{\mathrm{prod}}(E,F)$ & $D_N(E,F)$\\
\midrule
$J_j=0$ for all $j$, $\Gamma_E(H)\in\mathcal W_E$
 & $\min\{1,\sqrt N s\}$ & $\min\{1,\sqrt N s\}$\\
$J_j=0$ for all $j$, $\Gamma_E(H)\notin\mathcal W_E$
 & $\min\{1,\sqrt N s\}$ & $\min\{1,Ns\}$\\
Some $J_j\ne0$
 & $\min\{1,Ns\}$ & $\min\{1,Ns\}$\\
\bottomrule
\end{tabular}
\end{center}
Every comparison constant and the common small interval depend only
on the fixed $E,H,\Lambda$.  In particular, collective processing of
independently acquired outputs does not recover the linear scale of
the second row.  The assertion does not bound arbitrary entangled
parallel inputs or all classical-feedback acquisition models.
\end{theorem}
\begin{proof}
Theorem~\ref{thm:tubedichotomy86} gives the adaptive column.  Its
repeated-input lower gives the product lower in the first two rows;
\eqref{eq:leakagelower86} and \eqref{eq:tubehybrid86} give both
product bounds in the last row.  It remains to prove a
reference-uniform product upper when all $J_j=0$.

Use the canonical first-order factors
$A_j=\sqrt{E_j}$ and $B_j=(\sqrt{E_j})^+H_j(I-P_j/2)$ from
Lemma~\ref{lem:curvelogical85}.  They need not be horizontal.
Put $b=\|B\|$, let $G(s)$ be their normalized measurement in
Lemma~\ref{lem:factorremainder86}, and set
\[
 a=\tfrac32 b+\sqrt{\Lambda+K_b}.
\]
For every input state with any reference, denote the corresponding
output states under $E,G(s),F$ by $\rho_0,\rho_g,\rho_f$.
Their Bures distance
$\beta(\rho,\sigma)=\sqrt{2-2f(\rho,\sigma)}$, where
$f(\rho,\sigma)=\|\sqrt\rho\sqrt\sigma\|_1$ is root fidelity,
satisfies
\[
 \beta(\rho_0,\rho_g)\le\tfrac32bs,
 \qquad
 \beta(\rho_g,\rho_f)^2\le\|\rho_g-\rho_f\|_1
                        \le(\Lambda+K_b)s^2.
\]
The first inequality uses purifications from the two stacked factors
and their operator-norm difference; it holds uniformly over the input.
The second uses \eqref{eq:channelnorm86},
\eqref{eq:factorremainder86}, and the fidelity--trace inequalities.
The triangle inequality for $\beta$ follows by choosing common
purifications at the middle state and applying the Hilbert-space
triangle inequality.  Thus $1-f(\rho_0,\rho_f)\le a^2s^2/2$.
These standard fidelity properties and their unhalved trace convention
are as in \cite[Proposition~3.16 and Theorems~3.22, 3.33]{Watrous65}.

For different independent probe--reference pairs the output states
are still products.  Root fidelity is multiplicative, and
$1-\prod_r f_r\le\sum_r(1-f_r)$.  Therefore
\[
 \left\|\bigotimes_r\rho_{0,r}-\bigotimes_r\rho_{f,r}\right\|_1
 \le2\sqrt{1-\prod_r f_r^2}
 \le2a\sqrt N s.
\]
A joint final readout contracts trace norm.  Convexity gives the same
bound for public randomization, including a retained record of that
randomization.  Capping the result at two proves the required upper.
\end{proof}

\begin{corollary}[Higher-order approaches with a quadratic remainder]
\label{cor:powerjets86}
Let $q\ge1$ be an integer, and let a legal measurement family $F(t)$,
$0<t<a$, satisfy for fixed $E,H\ne0,\Lambda$
\begin{equation}\label{eq:powerjet86}
 \|F(t)-E-t^qH\|_{\Sigma}\le\Lambda t^{2q}.
\end{equation}
Then $H$ satisfies \eqref{eq:onesidedcone86}, and the conclusions of
Theorems~\ref{thm:tubedichotomy86} and
\ref{thm:producttrichotomy86} hold with $s=t^q$.
This includes zero first derivatives when $q>1$.
\end{corollary}
\begin{proof}
Compress \eqref{eq:powerjet86} to each missing support, divide by
$t^q$, and let $t\downarrow0$; positivity gives $J_j\succeq0$.
Substitute $s=t^q$ in the two finite-pair theorems.  No inverse
reparametrization is differentiated.  In particular, the conclusion
holds for $F(t)=E_*(t^q)$ whenever $E_*$ is a fixed one-sided $C^2$
curve with nonzero right derivative.  It also holds for a $C^{2q}$
curve whose first nonzero derivative has order $q$ and whose
intermediate derivatives of orders $q+1,\ldots,2q-1$ vanish, by Taylor's remainder after the leading nonzero coefficient.
\end{proof}

Condition \eqref{eq:powerjet86} is substantive.  A first nonzero
coefficient alone does not imply it.  The following exact family
displays what an intervening support-opening term can contribute.

\begin{proposition}[Mixed higher-order rates]\label{prop:mixedrates86}
Let $1\le a<b<2a$ be integers.  For scalar input dimension and three
outcomes put
\[
 E=(\tfrac12,\tfrac12,0),\qquad
 F(t)=\bigl((1-t^b)(\tfrac12+t^a),
            (1-t^b)(\tfrac12-t^a),t^b\bigr),
 \quad 0<t\le\tfrac14.
\]
There are absolute positive constants $c,C$ such that for all $N\ge1$
\begin{equation}\label{eq:mixedrates86}
 c\min\{1,\sqrt N t^a+Nt^b\}
 \le D_N(E,F(t))=D_N^{\mathrm{prod}}(E,F(t))
 \le C\min\{1,\sqrt N t^a+Nt^b\}.
\end{equation}
The leading coefficient $(1,-1,0)$ is in the covariance range.
Its order-$a$ contribution and the order-$b$ opening must both be
retained to obtain the finite-use law.
\end{proposition}
\begin{proof}
A scalar-input device only produces independent classical draws.
Every tester output, including one with bounded stopping, is a common
postprocessing of the full $N$-draw record.  Thus its optimal distance
is exactly the $\ell^1$ distance of the two product laws.
The event that label three occurs gives the lower
$2[1-(1-t^b)^N]\ge2(1-e^{-1})\min\{1,Nt^b\}$.
Send labels one and two to their corresponding bits and label three
to a fair bit.  This gives Bernoulli parameters $1/2$ and
$1/2+(1-t^b)t^a$; their gap is at least $t^a/2$.
Lemma~\ref{lem:bernoulli85} gives a second lower of order
$\min\{1,\sqrt N t^a\}$.  The larger of these lower bounds controls
one half of the truncated sum.

Under $F(t)$ the probability of no third label is $(1-t^b)^N$;
conditioned on that event the record is Bernoulli with parameter
$1/2+t^a$.  Splitting off the complementary event gives the upper
\[
 2[1-(1-t^b)^N]
 +\|\operatorname{Ber}(1/2)^{\otimes N}
       -\operatorname{Ber}(1/2+t^a)^{\otimes N}\|_1.
\]
For $u=t^a$, the one-copy root fidelity obeys
$f^2=(1+\sqrt{1-4u^2})/2$, so $1-f^2\le2u^2$.
Multiplicativity and the trace upper give a bound $2\sqrt{2N}\,u$
for the Bernoulli term.  Use $1-(1-t^b)^N\le Nt^b$ and cap at two.
For the range assertion, the scalar support generator is zero and
the missing-support component of $(1,-1,0)$ vanishes; apply
\eqref{eq:fullrange86}.  Since $b<2a$, the opening remainder is not
$O(t^{2a})$, so this example does not satisfy
\eqref{eq:powerjet86} with that leading coefficient.
\end{proof}

For example, $a=2,b=3$ gives
$\min\{1,\sqrt N t^2+Nt^3\}$ despite a vanishing first derivative.
The scalar family $(t^q,1-t^q)$, in contrast, has the exact distance
$2[1-(1-t^q)^N]$ and satisfies \eqref{eq:powerjet86} in the opening
branch.  A zero tangent supplied without higher-order data is still
an undetermined higher-order problem, not a statement of stationarity.

\begin{proposition}[Exact cone and finite-neighborhood certificates]
\label{prop:coneexact86}
For Gaussian-rational $E,H$, one-sided first-order realizability,
lineality, covariance-range membership and the three mechanisms of
Theorem~\ref{thm:producttrichotomy86} are decidable in polynomial
represented-input bit complexity.  For rational $s>0$, $\Lambda\ge0$
and $\lambda_j\ge0$ with $\sum_j\lambda_j\le\Lambda$, legality of
$F$ and the sufficient finite-neighborhood inequalities
\begin{equation}\label{eq:pairpredicate86}
 -\lambda_js^2I\preceq F_j-E_j-sH_j\preceq\lambda_js^2I
 \quad(1\le j\le k)
\end{equation}
are also exactly decidable with that complexity.
A nonzero opening supplies a rationally represented witness and its
exact output probability at a supplied pair.  No local interval,
curvature bound, recovery circuit or exact adaptive distance is
inferred from a first jet alone.
\end{proposition}
\begin{proof}
The rational support projectors of
Proposition~\ref{prop:supportexact84} give every $J_j$.
Exact positive-semidefinite tests establish \eqref{eq:onesidedcone86};
zero tests establish lineality.  When the $J_j$ vanish, the same
real-rank and Hilbert--Schmidt Gram projection decides whether
$\Gamma_E(H)$ belongs to $\mathcal W_E$.
Lemma~\ref{lem:onesidedcone86} gives necessity and sufficiency,
not only a formal positivity prefilter.

If $J_j\succeq0$ is nonzero, some diagonal element $(J_j)_{rr}$ is
positive.  Put $v=Q_je_r$.  Then $v\ne0$,
$v^*H_jv=(J_j)_{rr}>0$, and $\|v\|^2=(Q_j)_{rr}>0$.
Both the rate $v^*H_jv/\|v\|^2$ and, for a supplied $F$,
the probability $v^*F_jv/\|v\|^2$ are rational.
The unit vector itself need not have rational coordinates: its
rational vector and squared normalization specify the state exactly.
Finally \eqref{eq:pairpredicate86} is a finite list of rational PSD
tests.  Rational elimination, support projection and these tests
have the polynomial bit bounds of the retained exact kernels.
Verifying one supplied pair does not verify a neighborhood for an
unspecified curve, or that a supplied $s$ lies in the theorem's
small interval.  The exact impossible-event probability, when
provided, has its finite $N$-draw interpretation without that caveat.
\end{proof}

\section{Nonempty tangent neighborhoods}
\label{sec:nonempty87}
A prescribed remainder allowance need not admit any legal perturbation.
The following quantitative realization makes the nonvacuous range
explicit.  It is a sufficient allowance, not a minimal one.  The
underlying positive-semidefinite tangent condition is standard
semidefinite variational geometry; see \cite{BCS87}.  What is needed
here is one common normalization and a finite remainder bound for the
whole measurement tuple.

\begin{proposition}[A sufficient allowance]\label{prop:nonempty87}
Let $E$ be an ordered $k$-outcome measurement and let $H$ be a
Hermitian zero-sum tuple with $Q_jH_jQ_j\succeq0$ for every $j$.
Choose $a>0$ and $\lambda>0$ such that
\[
 \|H_j\|_{\mathrm{op}}\le a,\qquad
 E_j\succeq\lambda P_j\quad(1\le j\le k),
 \qquad P_j=\operatorname{supp}E_j.
\]
Set
\begin{equation}\label{eq:nonemptyconstants87}
 c=1+2a^2/\lambda,\qquad
 s_* =\min\{1,\lambda/(2a),1/a\},\qquad
 \Lambda_* =ck(1+2k).
\end{equation}
Then, for every $0\le s\le s_*$,
\begin{equation}\label{eq:nonemptycurve87}
 F_j^{\mathrm c}(s)=\frac{E_j+sH_j+cs^2I_d}{1+kcs^2}
\end{equation}
is a legal measurement and satisfies
\begin{equation}\label{eq:nonemptybound87}
 \|F^{\mathrm c}(s)-E-sH\|_{\Sigma}\le\Lambda_*s^2.
\end{equation}
Thus every allowance $\Lambda\ge\Lambda_*$ gives a nonempty
neighborhood at every $s\in(0,s_*]$.  For arbitrary real input,
\eqref{eq:nonemptycurve87} is a rational function of $s$ with matrix
coefficients.  For Gaussian-rational $E,H$ one may choose rational
$a,\lambda,c,s_*,\Lambda_*$, giving Gaussian-rational coefficients.
\end{proposition}
\begin{proof}
For $P_j\ne0$, the support block of the numerator in
\eqref{eq:nonemptycurve87} is at least $\lambda P_j/2$.
Its cross block is $sP_jH_jQ_j$.  The Schur complement on the
missing support is at least
\[
 sQ_jH_jQ_j+s^2(c-2a^2/\lambda)Q_j\succeq0.
\]
If $P_j=0$, the hypothesis says $H_j\succeq0$, so the numerator
is positive directly.  This also covers zero outcomes without
assigning them a smallest positive eigenvalue.  The common denominator
is positive for all real $s$, and summing the numerators gives exactly
$(1+kcs^2)I_d$.

Direct subtraction, with no series expansion, gives
\[
 F_j^{\mathrm c}(s)-E_j-sH_j
 =\frac{cs^2(I_d-kE_j-ksH_j)}{1+kcs^2}.
\]
Every effect has norm at most one, and $sa\le1$ on the given
interval.  The component remainder is therefore at most
$c(1+2k)s^2$.  Summing proves \eqref{eq:nonemptybound87}.
There is a positive common support floor since there are finitely
many effects and at least one is nonzero.  Rational minorants and
majorants give the last assertion.
\end{proof}

\begin{proposition}[Exact nonemptiness certificates]\label{prop:nonemptyexact87}
On Gaussian-rational represented input, a sufficient certificate
\eqref{eq:nonemptyconstants87} can be constructed and verified in
polynomial bit complexity.  It certifies the whole explicit curve
\eqref{eq:nonemptycurve87}, not only sampled values.  Its interval
$s_*$ is an interval of legality and remainder control; it is not a
computed small interval for the discrimination laws.
\end{proposition}
\begin{proof}
Use the rational support projectors of
Proposition~\ref{prop:supportexact84}.  The maximum absolute row sum,
with $|\Re z|+|\Im z|$ in place of $|z|$, gives a rational upper
bound $a\ge1$ for every Hermitian $H_j$.
Starting at $\lambda=1$, halve until the exact positive-semidefinite
tests $E_j-\lambda P_j\succeq0$ all hold.  Nonzero eigenvalues of a
rational positive matrix have a lower bound $2^{-\operatorname{poly}(d,L)}$
in terms of its represented bit length $L$: clear denominators and
use the nonzero coefficient of lowest degree in the characteristic
polynomial, together with the row-sum bound on the other eigenvalues.
Thus only polynomially many halvings are required.  All remaining
operations in \eqref{eq:nonemptyconstants87} are rational arithmetic.
The displayed PSD predicates and algebraic proof above verify the
certificate on the continuum.  A resource-limited implementation
must reject an exhausted cap, rather than return a partial certificate.
No claim that $\Lambda<\Lambda_*$ is infeasible follows from this
sufficient construction.
\end{proof}

\section{Entanglement width and parallel discrimination}
\label{sec:width87}
We now allow correlations between inputs to several calls, but keep
independent groups of bounded size.  The bound is imposed on the
complete probe--reference groups; a reference shared between groups
would be an additional correlation resource.

\begin{definition}[Independent entangled blocks]\label{def:blocks87}
For integers $1\le b\le N$, a block-$b$ tester chooses a partition
$n_1+\cdots+n_\ell\le N$, $1\le n_r\le b$, and an input state
\[
 \sigma=\bigotimes_{r=1}^{\ell}\sigma_r,
 \qquad
 \sigma_r\in\mathcal S\bigl((\C^d)^{\otimes n_r}\otimes R_r\bigr).
\]
The reference spaces are arbitrary.  All inputs are prepared before
any call; there is no feedback into an input.  All classical outputs
and retained references may be processed jointly after the calls.
Public randomization may choose both the partition and the states,
with its classical record retained.  Write $D_N^{[b]}$ for the
supremum of the unhalved final trace separation in this class.
\end{definition}

Thus $D_N^{[1]}=D_N^{\mathrm{prod}}$.  At $b=N$ there is no
restriction on the joint input and reference, and
\begin{equation}\label{eq:parallelidentity87}
 D_N^{[N]}(E,F)=D_N^{\mathrm{par}}(E,F)
 =\|\mathcal M_E^{\otimes N}-\mathcal M_F^{\otimes N}\|_\diamond.
\end{equation}
Unused calls can be padded and discarded.  Always
$D_N^{[1]}\le D_N^{[b]}\le D_N^{\mathrm{par}}\le D_N$.
Separable mixtures over complete probe--reference pairs are included
by public randomization.  Separability only on the probes, while an
arbitrary shared reference remains, is not the same restriction.

\begin{lemma}[A finite block upper]\label{lem:blockbures87}
Fix $E,H,\Lambda$ as in Definition~\ref{def:tangenttube86}, with
$Q_jH_jQ_j=0$ for every $j$.  Let $B$ be the canonical first-order
factor solution in the proof of Theorem~\ref{thm:producttrichotomy86},
and put
\[
 \beta_0=\|B\|_{\mathrm{op}},\quad
 K=k\beta_0^2(2+\max_j\|H_j\|_{\mathrm{op}}),\quad r=\Lambda+K.
\]
For $0<s\le1$ and $s\beta_0\le1$, every legal $F$ in the
quadratic neighborhood satisfies, for a tester with block sizes $n_r$,
\begin{align}
 \|\rho_E-\rho_F\|_1
 &\le \min\left\{2,
  2s\left[\sum_r
       (\tfrac32\beta_0n_r+\sqrt{rn_r})^2\right]^{1/2}\right\}
                                                        \label{eq:blockfinite87}\\
 &\le\min\{2,\,2(\tfrac32\beta_0+\sqrt r)s\sqrt{Nb}\}.
                                                        \label{eq:blockupper87}
\end{align}
Both inequalities are uniform over entangled states within a block,
all reference dimensions, and final collective processing.
\end{lemma}
\begin{proof}
Let $G(s)$ be the normalized surrogate from
Lemma~\ref{lem:factorremainder86}, and let $W_s,W_0$ be its and the
base measurement's common Stinespring factors, with the duplicated
label in a discarded environment.  That lemma gives
$\|W_s-W_0\|\le3\beta_0s/2$ and
$\|\mathcal M_F-\mathcal M_{G(s)}\|_\diamond\le rs^2$.
For one block of $n$ calls, telescoping tensor products of isometries
gives $\|W_s^{\otimes n}-W_0^{\otimes n}\|\le3\beta_0ns/2$.
Purify any block input, including its reference.  The corresponding
output states have Bures distance at most $3\beta_0ns/2$.
Channel telescoping also gives
\[
 \|\mathcal M_F^{\otimes n}-\mathcal M_{G(s)}^{\otimes n}\|_\diamond
 \le nrs^2.
\]
The Bures distance from the surrogate output to the $F$ output is
therefore at most $s\sqrt{rn}$.  Notice the retained $nrs^2$
channel remainder: it is not set to zero when the block grows.
The Bures triangle inequality bounds the block distance by
$s(3\beta_0n/2+\sqrt{rn})$.

Before final processing the outputs of different blocks are products.
Writing $f_r$ for their root fidelities, their squared Bures distance
is $2(1-\prod_rf_r)\le\sum_r2(1-f_r)$.  The unhalved trace
norm is at most twice the Bures distance.  This proves
\eqref{eq:blockfinite87}.  Since $n_r\ge1$ and
$\sum_rn_r^2\le b\sum_rn_r\le Nb$, it also proves
\eqref{eq:blockupper87}.  A final channel contracts trace norm.
For a retained randomization record, the states are direct sums with
the same weights under both hypotheses; their trace norm difference
is the weighted sum of the conditional differences.  Hence the same
uniform bound holds for randomized preparations and partitions.
The fidelity facts used here are those stated in the proof of
Theorem~\ref{thm:producttrichotomy86}, not an assumption of fidelity
multiplicativity inside an entangled block.
\end{proof}

\begin{lemma}[Parallel corrected phase accumulation]\label{lem:parallelcode87}
Suppose $Q_jH_jQ_j=0$ for all $j$ and
$\Gamma=\Gamma_E(H)\notin\mathcal W_E$.  There are a fixed logical
encoding $\mathcal J$ and a fixed recovery $\mathcal R$ such that
\[
 \Phi_E=\operatorname{Id}_2,\qquad
 \|\Phi_F-\mathcal U_s\|_\diamond\le\kappa s^2,
 \qquad
 \kappa=\Lambda+2\|\Gamma\|_{\mathrm{op}}^2,
\]
where $\Phi_F=\mathcal R(\mathcal M_F\otimes\operatorname{Id})
\mathcal J$ and $\mathcal U_s=\operatorname{Ad}(e^{-is\Gamma_L})$.
The two eigenvalues of $\Gamma_L$ have a gap $\Delta>0$.
For every integer $m\ge1$, a parallel block of $m$ calls has
separation at least
\begin{equation}\label{eq:parallelphase87}
 \bigl(2|\sin(ms\Delta/2)|-m\kappa s^2\bigr)_+.
\end{equation}
It also has a fixed binary readout whose base probability is $1/2$
and whose alternative probability differs from
$1/2+\sin(ms\Delta)/2$ by at most $m\kappa s^2/2$.
The encoding, recovery and readout are independent of the unknown
remainder $F-E-sH$.
\end{lemma}
\begin{proof}
The one-call assertion, including complete trace-preserving recovery
on the actual label and reference, is exactly the corrected-channel
calculation in Theorem~\ref{thm:tubedichotomy86}.  We change the
acquisition, not that calculation.  In the logical eigenbasis prepare
\[
 |\mathrm{GHZ}_m\rangle
   =2^{-1/2}(|0\rangle^{\otimes m}+|1\rangle^{\otimes m}),
\]
and apply the tensor encoding $\mathcal J^{\otimes m}$ before any
call.  All $m$ device inputs are now fixed.  After all calls apply
$\mathcal R^{\otimes m}$ to the actual labels and retained references.
The output is $\Phi_F^{\otimes m}$ applied to this GHZ state, and
channel tensor telescoping gives
\[
 \|\Phi_F^{\otimes m}-\mathcal U_s^{\otimes m}\|_\diamond
 \le m\kappa s^2.
\]
Under the base channel it is the original GHZ state.  The ideal
alternative multiplies its two components by phases whose difference
is $ms\Delta$, giving the pure-state trace separation
$2|\sin(ms\Delta/2)|$ and \eqref{eq:parallelphase87}.

On their two-dimensional span choose the sign of the Pauli $Y$
observable so that its ideal expectation is $\sin(ms\Delta)$;
extend it by zero on the orthogonal complement.  The effect
$(I+Y)/2$ is defined on the entire output, and has base probability
$1/2$.  The probability of one event changes by at most half the
trace norm between normalized states.  This proves the stated
binary readout bound.  The reference dimension in the encoding is
at most $2d$ per call, by Lemma~\ref{lem:supportcode84}.
No recovered logical output is sent into a later device call; all
recoveries are final operations.  In particular this construction
is parallel, not a deferred-description of a sequential input choice.
\end{proof}

\begin{theorem}[Entanglement-width law]\label{thm:widthlaw87}
Fix $E,H\ne0,\Lambda$ as in Definition~\ref{def:tangenttube86}.
There exist $c,C,s_0>0$, depending only on these fixed data, such that
for every pair of integers $1\le b\le N$, every $0<s\le s_0$,
and every legal $F$ with $\|F-E-sH\|_{\Sigma}\le\Lambda s^2$,
\begin{equation}\label{eq:widthlaw87}
 c\,\omega_{N,b}(s)\le D_N^{[b]}(E,F)\le C\,\omega_{N,b}(s),
\end{equation}
where
\[
 \omega_{N,b}(s)=
 \begin{cases}
 \min\{1,\sqrt N s\},&H\in\operatorname{Ran}\mathcal C_E,\\
 \min\{1,\sqrt{Nb}\,s\},&Q_jH_jQ_j=0\ \text{for all }j,
                                  \quad\Gamma_E(H)\notin\mathcal W_E,\\
 \min\{1,Ns\},&\text{some }Q_jH_jQ_j\ne0.
 \end{cases}
\]
The same $c,C,s_0$ work for all $b,N$ and all permitted remainders.
The lower testers may be chosen without knowing that remainder.
Nonemptiness is guaranteed, for instance, by
Proposition~\ref{prop:nonempty87} when $\Lambda\ge\Lambda_*$;
for a smaller prescribed allowance the theorem quantifies only over
its legal members.  The controls in the coherent branch are ideal
known-direction controls, not a common learner or a gate-synthesis
claim.
\end{theorem}
\begin{proof}
The first row follows from the product lower and adaptive upper of
Theorem~\ref{thm:tubedichotomy86}, by the inclusions of strategy
classes.  Its finite upper remains
$2\beta_h\sqrt N s+(\Lambda+K_h)Ns^2$, capped at two, for the
\emph{horizontal} factor solution.  The third row similarly follows
from its impossible-event product lower and hybrid adaptive upper.
Neither row requires entanglement.

For the second row, Lemma~\ref{lem:blockbures87} proves the upper
using the canonical, generally nonhorizontal, solution.  Use
$\Delta,\kappa$ from Lemma~\ref{lem:parallelcode87}, decrease $s_0$
so that
\[
 s_0\le\frac1{4\Delta},\qquad
 s_0\le\frac{\Delta}{\pi\kappa},
\]
and set
\begin{equation}\label{eq:blockchoice87}
 m=\min\{b,\lfloor(2\Delta s)^{-1}\rfloor\},\qquad
 \ell=\lfloor N/m\rfloor,\qquad x=m\Delta s.
\end{equation}
The floor before taking the minimum is at least two; hence
$1\le m\le b$ and $\ell\ge N/(2m)$.  Prepare $\ell$ independent
GHZ blocks of size $m$, discard unused calls, and use the binary
readout from the preceding lemma on each block.  The outputs are
independent Bernoulli variables with common base parameter $1/2$.
Since $0<x\le1/2$, their alternative parameter $q_F$ satisfies
\[
 q_F-\tfrac12
 \ge\tfrac12\sin x-\tfrac12m\kappa s^2
 \ge\frac{x}{\pi}-\frac{\kappa s}{2\Delta}x
 \ge\frac{x}{2\pi}.
\]
This holds for every allowed $F$ and pays its complete finite
remainder.  Lemma~\ref{lem:bernoulli85} gives the lower
\begin{equation}\label{eq:blockbernoulli87}
 D_N^{[b]}(E,F)\ge\frac1{128}
             \min\{1,\sqrt\ell\,x/(2\pi)\}.
\end{equation}
If $m=b$, then
$\sqrt\ell\,x\ge\Delta s\sqrt{Nb/2}$, as required.
If $m<b$, the floor bound implies $x\ge1/4$, so
\eqref{eq:blockbernoulli87} is at least $1/(1024\pi)$, which
also bounds a fixed positive multiple of the truncated target rate.
For example the coherent lower constant may be taken as
$128^{-1}\min\{1,\Delta/(2\sqrt2\pi),1/(8\pi)\}$.
This completes both cases without losing uniformity when $b$ grows.
\end{proof}

\begin{corollary}[Parallel and adaptive local orders]\label{cor:parallel87}
In every fixed quadratic tangent neighborhood with $H\ne0$,
\[
 D_N^{\mathrm{par}}(E,F)\asymp_{E,H,\Lambda}D_N(E,F)
\]
uniformly for all $N\ge1$ and sufficiently small $s>0$.
Both have square-root order in the covariance range and linear
order outside it.  In the coherent tangential branch the gain over
independent acquisitions is realized by pre-existing inter-call
entanglement; feedback is not necessary for that local exponent.
\end{corollary}
\begin{proof}
Set $b=N$ in Theorem~\ref{thm:widthlaw87} and compare with
Theorem~\ref{thm:tubedichotomy86}.
\end{proof}

This compares orders, not exact distances, perfect distinguishability,
or asymptotic error exponents at a fixed nonlocal pair.  It therefore
makes no general assertion that adaptation is useless in channel
discrimination.  The product converse in
Theorem~\ref{thm:producttrichotomy86} remains restricted to its stated
class; the parallel lower is supplied by a different proof here.

\paragraph{A complete qubit probability law.}
Let $E=(|0\rangle\langle0|,|1\rangle\langle1|)$,
$H=(\sigma_y,-\sigma_y)$, and
\[
 F_0(s)=\frac1{1+s^2}
       \begin{pmatrix}1&-is\\is&s^2\end{pmatrix},\qquad
 F_1(s)=I-F_0(s).
\]
These are legal for every real $s$, and
$\|F(s)-E-sH\|_\Sigma=2s^2/\sqrt{1+s^2}\le2s^2$.
On $m$ parallel inputs prepare
$2^{-1/2}(|+\rangle^{\otimes m}-i|-\rangle^{\otimes m})$,
where $|\pm\rangle$ are the $\sigma_x$ eigenvectors.
Direct multiplication of the rank-one effects gives, for every
ordered binary string $z$,
\begin{equation}\label{eq:parallelqubit87}
 p_s(z)=2^{-m}\bigl(1+(-1)^{|z|}\operatorname{Im}(w_s^m)\bigr),
 \qquad w_s=\frac{(1+is)^2}{1+s^2}.
\end{equation}
The two diagonal GHZ terms each give $2^{-m}$ before their
mixture weights; their cross term gives the displayed signed
imaginary part.  The base law is uniform, and summing over all
strings yields
$\|p_s-p_0\|_1=|\operatorname{Im}(w_s^m)|
=|\sin(2m\arctan s)|$.  Parity is a sufficient readout for this
pair of laws.  This is a concrete instance of standard GHZ phase
accumulation, without references or feedback, rather than a
claim of a new metrological mechanism.  Here the actual allowance
$\Lambda=2$ suffices even though the general construction with
$a=\lambda=1$ gives the coarse sufficient value $\Lambda_*=30$.

\begin{corollary}[Controlled power approaches]\label{cor:widthpower87}
If a legal family satisfies
$\|F(t)-E-t^qH\|_{\Sigma}\le\Lambda t^{2q}$ with $H\ne0$
and an integer $q\ge1$, then Theorem~\ref{thm:widthlaw87} and
Corollary~\ref{cor:parallel87} hold with $s=t^q$.
In the coherent branch the block rate is
$\min\{1,\sqrt{Nb}\,t^q\}$.  No such conclusion is asserted
from a leading coefficient without the stated remainder bound.
\end{corollary}
\begin{proof}
The cone condition follows by missing-support compression and
$t^{-q}$ rescaling, as in Corollary~\ref{cor:powerjets86}.
Substitute $s=t^q$; no inverse parametrization is differentiated.
\end{proof}

\section{Classical feedback and coherent reset width}
\label{sec:reset88}
We use normal quantum instruments on trace-class operators.  Countably
many recorded outcomes form an $\ell^1$ direct sum.  All identities below
follow from finite truncations retaining a remainder outcome and passage
to the trace-norm limit; positivity gives monotone convergence of trace
sums.  The hard finite depth makes the conditional induction applicable.
The independent blocks of Definition~\ref{def:blocks87} do not
exhaust the strategies with bounded coherent acquisition.  We now
allow both classical feedback between blocks and adaptive quantum
control within a block.  The receiver may retain and process quantum
memory throughout.  The restriction is on importing that memory
into the next acquisition, not on its dimension or final readout.

\begin{definition}[Reset protocols]\label{def:reset88}
Fix integers $1\le b\le N$.  At a classical history $h$, a protocol
may stop, or choose a positive integer $n(h)\le b$ and prepare a new
complete probe--reference system in a state $\sigma_h$.  Conditional
on $h$, this state is the same under both hypotheses and is in a
tensor product with the receiver's entire existing memory.  A common
$n(h)$-call adaptive experiment acts on the new system alone; its
internal inputs, quantum memory and controls are unrestricted.
After acquisition, its whole accessible output is delivered to the
receiver.  A common quantum instrument may act jointly on this
output and all previous receiver memory, retaining quantum memory
and producing the next classical history.  No quantum subsystem of
that memory enters a subsequent acquisition.  On every history,
\[
 \sum_r n(h_r)\le N.
\]
The integers are reserved call budgets: unused internal slots can
be padded and are still charged.  Public stopping at block boundaries
is allowed.  All classical histories are finite strings with finite
or countable outcome sets; conditional state preparations and
instruments are defined also at histories null under one hypothesis.
Let $D_N^{\mathrm{reset},b}(E,F)$ be the supremum of the unhalved
trace separation of the final accessible states, including the
classical record and arbitrary final joint processing.
\end{definition}

The complete new system, including every reference and purification
used in acquisition, must be independent of the receiver conditional
on $h$.  Separability of the probe marginal alone is not sufficient.
The receiver memory need not be classical, and can be entangled with
previous block outputs.  Measuring that memory to choose a fresh
input is allowed; transmitting it coherently into that input is not.
In particular,
\begin{equation}\label{eq:resetinclusions88}
 D_N^{[b]}\le D_N^{\mathrm{reset},b}\le D_N,
 \qquad D_N^{\mathrm{reset},N}=D_N.
\end{equation}
The first inclusion uses no feedback and postpones receiver processing;
the equality uses one unrestricted adaptive block.  At $b=1$ the
receiver can still adapt future preparations through classical
messages.  Equality with the independent-product distance is not
asserted.

For positive, possibly subnormalized operators, write
$f(R,S)=\|\sqrt R\sqrt S\|_1$.  For normalized states put
$d_B(R,S)^2=2(1-f(R,S))$.  We use
\begin{equation}\label{eq:rootrules88}
\begin{split}
 f(R\otimes\rho,S\otimes\sigma)&=f(R,S)f(\rho,\sigma),\\
 f\left(\bigoplus_zR_z,\bigoplus_zS_z\right)&=\sum_z f(R_z,S_z),\\
 f(\mathcal A(R),\mathcal A(S))&\ge f(R,S)
\end{split}
\end{equation}
for a channel $\mathcal A$.  The last assertion for unnormalized
operators follows by separate normalization and homogeneity.
For normalized states,
\begin{equation}\label{eq:tracebures88}
 d_B(R,S)^2\le\|R-S\|_1
 \le2\sqrt{1-f(R,S)^2}\le2d_B(R,S).
\end{equation}
These are standard fidelity identities and inequalities
\cite[Chapter~3]{Watrous65}.

\begin{lemma}[Conditional fidelity accumulation]\label{lem:resetfidelity88}
Consider a reset protocol in which the fresh output states at node
$h$, before receiver processing, satisfy
\[
 f(\omega_{0,h},\omega_{1,h})\ge(1-a_h)_+,
 \qquad a_h\ge0.
\]
Suppose every root-to-terminal history has $\sum_h a_h\le A$.
For any initial receiver operators $R_0,R_1\succeq0$, the final
classical--quantum operators satisfy
\begin{equation}\label{eq:fidelityaccumulation88}
 f(R_{0,\mathrm{out}},R_{1,\mathrm{out}})
 \ge (1-A)_+f(R_0,R_1).
\end{equation}
Thus the receiver's quantum processing and classical feedback need
not preserve a product state across blocks.
\end{lemma}
\begin{proof}
The assertion is immediate for $A\ge1$.  For $A<1$, prove the stronger
continuation statement at each node, for arbitrary incoming positive
operators and a bound $A_h<1$ on the remaining path sum.  At a terminal
node a final common channel only increases fidelity.  At a nonterminal
node let $a=a_h$ and bound every child's remaining sum by $A_h-a$.
Appending the fresh output multiplies the incoming fidelity by at
least $1-a$.  Apply the receiver instrument, retaining its classical
outcome.  If its two subnormalized outputs at outcome $z$ are
$R_{0,z},R_{1,z}$, channel monotonicity and direct-sum additivity give
\[
 \sum_z f(R_{0,z},R_{1,z})
 \ge (1-a)f(R_0,R_1).
\]
The induction hypothesis applies separately at every $z$, even
though the two conditional receiver states differ.  The continued
output fidelity is at least
\[
 [1-(A_h-a)](1-a)f(R_0,R_1)
 \ge(1-A_h)f(R_0,R_1).
\]
Here $0\le a\le A_h<1$, so the product estimate has the stated
sign.  Each acquisition costs at least one call, hence the depth is
at most $N$ and backwards induction terminates.  Countably many
outcomes are handled by the same nonnegative sum; no normalization
at a zero-probability history is required.  Conditional public
randomization is another retained instrument outcome.  This proves
the statement, including stopping and final joint processing.
\end{proof}

\begin{lemma}[A fresh adaptive block]\label{lem:adaptiveblock88}
Fix $E,H,\Lambda$ in Definition~\ref{def:tangenttube86}, and assume
$Q_jH_jQ_j=0$ for every $j$.  Let $A_j=\sqrt{E_j}$ and choose the
canonical factors $B_j$ of Lemma~\ref{lem:curvelogical85} with
$A_j^*B_j+B_j^*A_j=H_j$.  Set
\[
 \beta_0=\|B\|_{\mathrm{op}},\quad
 K=k\beta_0^2(2+\max_j\|H_j\|_{\mathrm{op}}),\quad
 r=\Lambda+K,\quad \alpha=3\beta_0/2.
\]
If $0<s\le1$ and $s\beta_0\le1$, the outputs of every common fresh
$n$-call adaptive experiment obey
\begin{equation}\label{eq:adaptiveblock88}
 d_B(\omega_E,\omega_F)\le s(\alpha n+\sqrt{rn}),\qquad
 1-f(\omega_E,\omega_F)
 \le \tfrac12(\alpha+\sqrt r)^2 n^2s^2,
\end{equation}
for every legal $F$ with $\|F-E-sH\|_\Sigma\le\Lambda s^2$.
The constants are independent of the input, reference, internal
controls and permitted remainder.
\end{lemma}
\begin{proof}
The normalized surrogate in Lemma~\ref{lem:factorremainder86} has
isometries $W_0,W_s$ with $\|W_s-W_0\|\le\alpha s$ and
\[
 \|\mathcal M_F-\mathcal M_{G(s)}\|_\diamond\le rs^2.
\]
Purify the fresh input and all internal common controls.  Inserting
the difference $W_s-W_0$ successively in the $n$ slots bounds the
final purified vector difference by $n\alpha s$: every other factor
is an isometry.  This argument is valid for adaptive controls, not
only tensor-power acquisition.  The discarded dilation environments
are used in the proof but are unavailable to the experiment.
Purification and contraction give the same Bures upper for the
accessible $E$ and surrogate outputs.  Replacing the surrogate by
$F$ costs at most $nrs^2$ in unhalved trace norm by adaptive channel
telescoping.  The first inequality of \eqref{eq:tracebures88} and the
Bures triangle inequality give the first bound in
\eqref{eq:adaptiveblock88}.  The second follows from
$1-f=d_B^2/2$ and $\sqrt n\le n$.  Internal early stopping is padded
within the already reserved $n$ slots and then discarded.  This
retains the full $nrs^2$ error before taking its square root.
\end{proof}

The canonical factor above is generally not horizontal.  Its
linear-in-$n$ purified displacement is appropriate for a coherent
block.  The separate horizontal construction remains essential to
the stronger regular-direction bound for unrestricted adaptive
strategies.

\begin{theorem}[Reset-width interpolation]\label{thm:resetwidth88}
Fix a measurement $E$, a nonzero Hermitian zero-sum tuple $H$ with
$J_j=Q_jH_jQ_j\succeq0$, and $\Lambda\ge0$.  There are
$c,C,s_0>0$, depending only on $E,H,\Lambda$, such that for every
$N\ge1$, $1\le b\le N$, $0<s\le s_0$, and every legal measurement
$F$ satisfying
\[
 \|F-E-sH\|_\Sigma:=\sum_j\|F_j-E_j-sH_j\|_{\mathrm{op}}
 \le\Lambda s^2,
\]
one has $c\mathfrak r\le D_N^{\mathrm{reset},b}(E,F)\le C\mathfrak r$,
where
\begin{equation}\label{eq:resetwidth88}
\mathfrak r=
\begin{cases}
 \min\{1,\sqrt N s\},&H\in\operatorname{Ran}\mathcal C_E,\\
 \min\{1,\sqrt{Nb}\,s\},&J_j=0\ (\forall j),\quad
                         \Gamma_E(H)\notin\mathcal W_E,\\
 \min\{1,Ns\},&J_j\ne0\ (\exists j).
\end{cases}
\end{equation}
Here $\Gamma_E(H)=\frac i2\sum_j[P_j,H_j]$ and
$\mathcal W_E=\sum_jP_j\mathcal H_dP_j$.
All constants are uniform in $N,b,F$, not in $E,H,\Lambda$.
For each $N,b,s$, the lower experiment can be chosen independently
of the unknown remainder.  It can in fact be chosen without
inter-block feedback and with parallel acquisition inside each block.
\end{theorem}
\begin{proof}
In the tangential cases put $\gamma=(\alpha+\sqrt r)^2/2$ as in
Lemma~\ref{lem:adaptiveblock88}.  For a particular public policy let
\[
 Q=\sup_{\text{terminal histories}}\sum_r n(h_r)^2\le bN.
\]
At a node, the fresh experiment has normalized outputs independent
of the incoming receiver memory, conditional on that history.
Lemma~\ref{lem:adaptiveblock88} bounds its fidelity deficit by
$a_h=\gamma n(h)^2s^2$.  Lemma~\ref{lem:resetfidelity88}, starting
from the same normalized initial memory, therefore gives
\[
 f(\rho_{E,\mathrm{out}},\rho_{F,\mathrm{out}})
 \ge(1-\gamma Qs^2)_+.
\]
If $\gamma Qs^2\le1$, use \eqref{eq:tracebures88} and
$1-(1-x)^2\le2x$; otherwise the trace bound two suffices.  In both
cases we obtain the finite policy-sensitive estimate
\begin{equation}\label{eq:resetfinite88}
 D(\text{this protocol};E,F)
 \le\min\{2,\,2(\alpha+\sqrt r)s\sqrt Q\}
 \le\min\{2,\,2(\alpha+\sqrt r)s\sqrt{bN}\}.
\end{equation}
The uniform per-block bound permits every history-dependent choice
of inputs and internal controls, and arbitrary common receiver
instruments.  The two output states are not assumed to be products
across feedback branches.

For the coherent row, the lower construction and the two integer
allocation cases of Theorem~\ref{thm:widthlaw87} belong to the reset
class by \eqref{eq:resetinclusions88}.  Its fixed binary event pays
$m\kappa s^2$ before the finite Bernoulli amplification.  Thus the
same constants and remainder-independent lower apply here.
For the regular row, the product lower and the fully adaptive
horizontal bound \eqref{eq:tubefiniteupper86} sandwich this class.
That upper retains $(\Lambda+K_h)Ns^2$, which is controlled by
$\sqrt N s$ when $\sqrt N s\le1$ and by the trace cap otherwise.
For the opening row, the fixed impossible event
\eqref{eq:leakagelower86} and the adaptive hybrid bound
\eqref{eq:tubehybrid86} give the result already without entanglement.
Take the minimum of the finitely many fixed scale restrictions.
The full range criterion \eqref{eq:fullrange86} makes the cases
disjoint and exhaustive.
\end{proof}

\begin{corollary}[Classical feedback at unit reset width]
\label{cor:classicalfeedback88}
In a fixed coherent tangent neighborhood,
\[
 D_N^{\mathrm{reset},1}(E,F)\asymp\min\{1,\sqrt N s\},\qquad
 D_N^{\mathrm{par}}(E,F)\asymp D_N(E,F)\asymp\min\{1,Ns\}.
\]
The first class permits arbitrary receiver quantum memory and
classical feedback, but no coherent import of that memory into a
fresh probe--reference preparation.  These are local comparison
orders, not equality of exact distances or fixed-pair error exponents.
\end{corollary}
\begin{proof}
Apply Theorem~\ref{thm:resetwidth88} at $b=1$ and combine the
parallel endpoint in Corollary~\ref{cor:parallel87} with the adaptive
tube theorem.  All intervals and constants refer to the same fixed
$E,H,\Lambda$.
\end{proof}

The role of feedback in channel discrimination has extensive prior
study \cite[Theorem~9 and Remark~10]{SHW87}.  Our assertion is a
finite, remainder-uniform consequence of a conditional fidelity
estimate, with arbitrary receiver memory and variable reserved block
sizes.  It does not introduce the classical-feedback resource model
or the standard-versus-coherent metrological mechanism.  The
isometric-extension estimates likewise have direct antecedents in
channel Bures bounds \cite[Theorems~17 and~19]{HMNW88}.

\subsection{Finite policy certificates}
For a finite public decision tree, its pathwise call maximum and
quadratic acquisition cost are computed recursively.  At a terminal
node set $N_h=Q_h=0$; at an acquisition node set
\[
 N_h=n(h)+\max_zN_{hz},\qquad
 Q_h=n(h)^2+\max_zQ_{hz}.
\]
These are maxima over complete histories, not expected costs under
either hypothesis.  Exact integer arithmetic verifies $N_h\le N$
and $n(h)\le b$ on the whole policy.  Given rational $s$ and a
separately established rational majorant $\bar\gamma\ge\gamma$,
\[
 f_{\rm out}\ge(1-\bar\gamma Qs^2)_+,
 \qquad \|\rho_{E,\rm out}-\rho_{F,\rm out}\|_1^2
       \le\min\{4,8\bar\gamma Qs^2\}.
\]
The same induction also gives the stronger recursion
$F_h=(1-\bar\gamma n(h)^2s^2)_+\min_z F_{hz}$ with terminal
value one.  A policy certificate checks these budget consequences.  It does not
verify the physical reset condition or establish the supplied
one-block fidelity majorant.  Those remain hypotheses.  In
particular a numerical or declared reference dimension is not a
certificate of conditional independence.

The explicit nonempty interval $s_*$ of
Proposition~\ref{prop:nonempty87} continues to certify legality and
the sufficient allowance only.  It is not the discrimination
interval $s_0$ above; that interval also depends on the lower witness
and logical correction gap.  The sufficient $\Lambda_*$ is neither
minimal nor canonical.  Zero tangents, intervening higher-order
terms, arbitrary-pair boundary equivalence and growing-output
learning retain their separate hypotheses and questions.

\section{Arbitrary allocation profiles}
\label{sec:allocation89}
The largest permitted block does not determine how that resource is
actually used.  We retain the entire list of block sizes.  Throughout
this section the groups are complete probe--reference systems, not
merely probe marginals.  We first specify independent acquisition,
then allow the reset-constrained feedback of Definition~\ref{def:reset88}
while keeping the complete reservation profile fixed.

\begin{definition}[A prescribed allocation]\label{def:allocation89}
Let $\boldsymbol n=(n_1,\ldots,n_\ell)$ be a nonempty finite list of
positive integers.  Write
\[
 N_{\boldsymbol n}=\sum_{r=1}^{\ell}n_r,
 \qquad V_{\boldsymbol n}=\sum_{r=1}^{\ell}n_r^2.
\]
An allocation tester has an input
$\bigotimes_r\sigma_r$, where $\sigma_r$ is an arbitrary state on
$(\C^d)^{\otimes n_r}\otimes R_r$.  Reference dimensions and final
joint processing are unrestricted.  Public mixtures, with the
mixture label retained, are allowed for this fixed allocation.
A group may use fewer than its assigned slots by padding with fixed
inputs and discarding their outputs.  Denote the supremum of the
unhalved final trace separation by $D^{\boldsymbol n}(E,F)$.
For the same list, a prescribed-profile reset tester reserves $n_r$
calls in its $r$th fresh block, independent of the classical history.
The preparation and all internal adaptive controls may depend on that
history, but the complete fresh system is conditionally independent
of all old receiver memory.  This is one common conditional rule on
every formal history, not an equality of the two conditional receiver
states.  Early termination can be padded by unused reserved blocks.
Our convention records a deterministic terminal symbol at each remaining
round, executes a fixed fresh preparation and discards its output.  All
positive reservations in the prescribed profile are still charged.
Zero-length terminal symbols are not additional profile components.
Write $D^{\mathrm{reset},\boldsymbol n}$ for this larger supremum.
Both $N_{\boldsymbol n}$ and $V_{\boldsymbol n}$ are hard budgets,
not expectations.  Then
\[
 D^{\boldsymbol n}\le D^{\mathrm{reset},\boldsymbol n}\le D_{N_{\boldsymbol n}}.
\]
\end{definition}

The first moment $N_{\boldsymbol n}$ governs independent statistical
accumulation; the second moment $V_{\boldsymbol n}$ governs coherent
accumulation within these groups.  The second-moment principle has
established metrological antecedents \cite{Hyllus78,Toth78}.  The point
here is a finite trace-distance statement uniform over a whole
measurement neighborhood, including its unknown quadratic remainder.
We first record the finite classical estimate that permits unequal
blocks without grouping them into dyadic size classes.

\begin{lemma}[A common readout for unequal signals]\label{lem:unequalsignals89}
Let $0<\alpha\le1$, $0\le x_r\le1/2$, and
$\alpha x_r\le u_r\le x_r$.  Let $P_0$ be the law of independent
fair signs $X_r\in\{-1,1\}$, and let $P_u$ be the independent law
with $P_u\{X_r=1\}=1/2+u_r$.  There is a binary randomized readout
depending only on $(x_r)$, and not on $(u_r)$, for which the unhalved
separation is at least
\begin{equation}\label{eq:unequalsignals89}
 \frac{\alpha}{64}\min\left\{1,
                         \left(\sum_r x_r^2\right)^{1/2}\right\}.
\end{equation}
Consequently the full product laws have at least this separation.
In the application, $(x_r)$ is fixed by $E,H,\Lambda,s$ and the public
profile.  The readout has precisely this public dependence and no
dependence on the unknown tuple $F-E-sH$ or actual biases $(u_r)$.
\end{lemma}
\begin{proof}
Put $S=\sum_r x_r^2$.  For $S=0$ the assertion is immediate.
Otherwise set
\[
 A=\max\{\sqrt S,S\},\qquad
 t_r=\frac{x_r}{16A},\qquad
 T(X)=\sin\left(\sum_r t_rX_r\right).
\]
Then $0\le t_r\le1/16$ and $\sum_r t_r^2\le1/256$.
Under $P_0$ the expectation of $T$ is zero.  Independence under $P_u$
gives
\[
 \mathbb E_u\exp\left(i\sum_r t_rX_r\right)
 =\prod_r(\cos t_r+2iu_r\sin t_r)=R e^{i\phi},
 \quad \phi=\sum_r\arctan(2u_r\tan t_r).
\]
Each factor has positive real part.  Also
\[
 R\ge\prod_r\cos t_r
 \ge1-\tfrac12\sum_r t_r^2\ge\tfrac12.
\]
For $0\le t\le1/16$, $t\le\tan t\le2t$.  The arguments of
$\arctan$ above lie in $[0,1]$, where $v/2\le\arctan v\le v$.
It follows that
\[
 \frac{\alpha S}{16A}\le\sum_r u_rt_r\le\phi
 \le4\sum_r u_rt_r\le\frac{S}{4A}\le\frac14.
\]
Since $\sin\phi\ge\phi/2$ on this interval,
$\mathbb E_uT=R\sin\phi\ge\alpha S/(64A)$.
Accept with probability $(1+T(X))/2$.  The resulting binary laws
have unhalved separation $|\mathbb E_uT-\mathbb E_0T|$.
Finally $S/A=\min\{1,\sqrt S\}$.  This proves the assertion with
one readout for every permitted vector of biases.
\end{proof}

\begin{theorem}[Allocation-profile law]\label{thm:allocation89}
Fix an ordered measurement $E$, a nonzero Hermitian tuple $H$ with
$\sum_jH_j=0$ and $Q_jH_jQ_j\succeq0$, and a number $\Lambda\ge0$.
Here $P_j=\operatorname{supp}E_j$, $Q_j=I-P_j$, and
$\|R\|_\Sigma=\sum_j\|R_j\|_{\mathrm{op}}$.
There are $c,C,s_0>0$, depending on $E,H,\Lambda$ alone, such that
for every allocation $\boldsymbol n$, every $0<s\le s_0$, and every
legal measurement $F$ satisfying
\begin{equation}\label{eq:allocationtube89}
 \|F-E-sH\|_\Sigma\le\Lambda s^2,
\end{equation}
one has
\begin{equation}\label{eq:allocationlaw89}
 c\,\omega_{\boldsymbol n}(s)
 \le D^{\boldsymbol n}(E,F)
 \le D^{\mathrm{reset},\boldsymbol n}(E,F)\le C\,\omega_{\boldsymbol n}(s),
\end{equation}
where
\[
 \omega_{\boldsymbol n}(s)=
 \begin{cases}
 \min\{1,\sqrt{N_{\boldsymbol n}}s\},
                    &H\in\operatorname{Ran}\mathcal C_E,\\
 \min\{1,\sqrt{V_{\boldsymbol n}}s\},
                    &Q_jH_jQ_j=0\ (1\le j\le k),\quad
                           \Gamma_E(H)\notin\mathcal W_E,\\
 \min\{1,N_{\boldsymbol n}s\},&\text{some }Q_jH_jQ_j\ne0.
 \end{cases}
\]
These constants are uniform in the number and sizes of groups, and
in the entire permitted remainder.  In every row a single lower
tester can be chosen from $E,H,\Lambda,s,\boldsymbol n$, independently
of that remainder.  The coherent construction uses ideal
known-direction controls.  It is neither a common learner nor a
statement about separable probe marginals with a shared reference.
Nonemptiness and a sufficient allowance are as in
Proposition~\ref{prop:nonempty87}; its legality interval $s_*$ does
not replace the discrimination interval $s_0$.
\end{theorem}
\begin{proof}
Write $N=N_{\boldsymbol n}$ and $V=V_{\boldsymbol n}$.
Product-reference inputs are admissible in every allocation, and every prescribed-profile reset tester is an adaptive tester.  The regular and opening
rows therefore follow from Theorem~\ref{thm:tubedichotomy86}.
In the regular row its finite upper is
\[
 2\beta_h\sqrt N s+(\Lambda+K_h)Ns^2,
\]
capped at two, for the horizontal factor solution.  This is of the
asserted order since $x^2\le x$ for $x=\sqrt N s\le1$.
The opening lower uses the fixed impossible-output event, not a
coherent code.

First consider the independent coherent upper.  Use the canonical, generally nonhorizontal,
factor and the constants $\beta_0,r_0$ of
Lemma~\ref{lem:blockbures87}, writing $r_0$ for its remainder constant $r$.
With root fidelity
$f(\rho,\sigma)=\|\sqrt\rho\sqrt\sigma\|_1$ and Bures distance
$\beta^2=2(1-f)$, the normalization is
\[
 \beta(\rho,\sigma)^2\le\|\rho-\sigma\|_1
                         \le2\beta(\rho,\sigma).
\]
The actual channel remainder in a group of $n_r$ calls is
$n_r r_0s^2$; its Bures contribution is $s\sqrt{r_0n_r}$.
The block estimate consequently gives the finite upper
\begin{align}
 D^{\boldsymbol n}(E,F)
 &\le \min\left\{2,\,
  2s\left[\sum_r
      (\tfrac32\beta_0n_r+\sqrt{r_0n_r})^2\right]^{1/2}\right\}
                                                   \label{eq:profilefinite89}\\
 &\le\min\{2,\,2s(\tfrac32\beta_0\sqrt V+\sqrt{r_0N})\}.
                                                   \label{eq:profileminkowski89}
\end{align}
The last line is the Euclidean triangle inequality.  Since $N\le V$,
it implies the required upper.  Fidelity multiplies only between
independent group outputs.  A retained public mixture gives a direct
sum with common weights, whose trace separation is the weighted sum
of the conditional separations; the same bound follows.  Final
joint processing is unrestricted and contracts trace norm.

The same finite upper holds for prescribed-profile reset testers,
for a different reason.  Lemma~\ref{lem:adaptiveblock88} gives the
normalized fresh output deficit
\[
 a_r=\tfrac12s^2(\alpha n_r+\sqrt{r_0n_r})^2,
 \qquad \alpha=\tfrac32\beta_0,
\]
uniformly at every history in round $r$.  These deficits have the
same deterministic sum along every full reservation path.  Apply
Lemma~\ref{lem:resetfidelity88} to positive subnormalized receiver
operators.  If $A=\sum_ra_r\le1$, then
$\|\rho_E-\rho_F\|_1\le2\sqrt{2A}$; if $A>1$, use the trace cap two.
This proves \eqref{eq:profilefinite89} and
\eqref{eq:profileminkowski89} with $D^{\mathrm{reset},\boldsymbol n}$
in place of $D^{\boldsymbol n}$.  It does not multiply unconditional
output fidelities after feedback.  Countable histories are interpreted
as trace-class direct sums; all fidelity summands are nonnegative.

For the lower let $\Delta,\kappa>0$ be as in
Lemma~\ref{lem:parallelcode87}, decreasing $s_0$ if necessary so that
\[
 s_0\le\frac1{4\Delta},\qquad
 s_0\le\frac{\Delta}{\pi\kappa}.
\]
In group $r$ use
\[
 m_r=\min\{n_r,\lfloor(2\Delta s)^{-1}\rfloor\},\qquad
 x_r=m_r\Delta s,
\]
logical GHZ inputs, with the remaining slots padded and discarded.
All encodings precede acquisition and all recoveries follow it.
The reference dimension is at most $2d$ per used call; the reference
of a whole group may grow with its size.
The fixed binary event of that lemma has base probability $1/2$ and
alternative probability $1/2+u_r$, with
\[
 \left|u_r-\tfrac12\sin x_r\right|
       \le\tfrac12m_r\kappa s^2.
\]
This is the full finite tensor-channel error, paid before statistical
amplification.  Since $0<x_r\le1/2$, the displayed choice of $s_0$
implies
\[
 \frac{x_r}{2\pi}\le u_r
 \le\frac{x_r}{2}+\frac{\kappa s}{2\Delta}x_r\le x_r.
\]
Independent group outputs need not have equal biases.
Lemma~\ref{lem:unequalsignals89}, with $\alpha=1/(2\pi)$, gives
\begin{equation}\label{eq:profilelowerfinite89}
 D^{\boldsymbol n}(E,F)\ge
 \frac1{128\pi}\min\left\{1,
                      \Delta s\left(\sum_r m_r^2\right)^{1/2}\right\}.
\end{equation}
Both the individual events and the final randomized classical
readout are independent of $F-E-sH$.

If every $m_r=n_r$, the last minimum is
$\min\{1,\Delta s\sqrt V\}$, at least
$\min\{1,\Delta\}\min\{1,s\sqrt V\}$.
If some $m_r<n_r$, then $(2\Delta s)^{-1}\ge2$ and the floor bound
gives $x_r\ge1/4$.  The last minimum is consequently at least $1/4$.
The target $\min\{1,s\sqrt V\}$ is at most one, so this also controls
it by an absolute positive factor.  In particular the coherent lower
constant may be taken as $\min\{\Delta,1/4\}/(128\pi)$.
This proves the uniform lower for both classes, including profiles with very unequal
or very large groups.  No output is used to prepare a later input.
\end{proof}

\begin{corollary}[Width budgets and effective width]\label{cor:profilebudget89}
For $1\le b\le N$, put $q=\lfloor N/b\rfloor$ and $r=N-qb$.
Among allocations with $\sum n_j\le N$ and $n_j\le b$,
\begin{equation}\label{eq:profilemax89}
 \max V_{\boldsymbol n}=qb^2+r^2,
 \qquad \tfrac12Nb\le qb^2+r^2\le Nb.
\end{equation}
Thus the coherent width law can equivalently be written with
$\sqrt{qb^2+r^2}\,s$, with constants independent of $N,b$.
For a prescribed allocation the effective width is
$b_{\rm eff}=V_{\boldsymbol n}/N_{\boldsymbol n}$, not necessarily
its largest part.  Coarsening an allocation increases its coherent
scale and leaves the other two scales unchanged.
\end{corollary}
\begin{proof}
Adding unused slots increases the objective.  If two parts are below
$b$, transfer slots from the smaller to the larger until one is zero
or the other is $b$.  Their sum of squares does not decrease.
Iteration gives $q$ parts equal to $b$ and, if $r>0$, one part $r$.
The upper bound follows from $n_j^2\le bn_j$.
Since $r<b\le qb$, one has $qb\ge N/2$, which proves the lower.
Merging parts $u,v$ adds $2uv$ to their squared sum.  For independent allocation testers, the operational
inclusion also holds because product inputs within a merged group
remain admissible.  No analogous set inclusion is asserted for
arbitrary reset-feedback testers when a receiver-interaction boundary
is removed; their local orders are the preceding theorem's statement.
\end{proof}

For example, with $N=L^2$ calls and largest block $L$, the allocation
$(L,1,\ldots,1)$ has $V=2L^2-L$, whereas $L$ groups each of size $L$
have $V=L^3$.  Before saturation their coherent scales differ by a
factor of order $\sqrt L$, despite equal call counts and equal largest
blocks.  At the single-group endpoint, the preceding theorem recovers
matching parallel and adaptive \emph{local orders}; it makes no claim
of equality of their exact distances or global error exponents.

The comparison concerns the best tester for a prescribed reservation
profile.  It does not state that each feedback policy using that
profile is informative.  For example, a common final channel that
forgets every output has zero separation.  Likewise, an arbitrary
policy's worst-path cost is not itself a matching lower certificate:
a costly branch can be uninformative or rarely visited.

\section{Quadratic acquisition budgets and receiver memory}
\label{sec:quadraticbudget89}
A width bound can conceal unused coherent acquisition.  The natural
policy quantity is the worst-path sum of squared reservations.  We
first retain the separate finite remainder, then identify the sharp
local order under this resource bound.  None of these quantities
measures the dimension of a quantum memory.

\begin{proposition}[Finite policy profile]\label{prop:policyprofile89}
Fix a legal quadratic tube as in Definition~\ref{def:tangenttube86},
and suppose $Q_jH_jQ_j=0$ for every $j$.  Use the constants
$\alpha,r$ and scale restrictions of Lemma~\ref{lem:adaptiveblock88}.
For a reset policy $\pi$, define the hard path bounds
\[
 N_\pi=\sup_{\zeta}\sum_{h\in\zeta}n(h),\qquad
 Q_\pi=\sup_{\zeta}\sum_{h\in\zeta}n(h)^2,
\]
where $\zeta$ ranges over formal terminal histories, and put
\[
 B_\pi=\sup_{\zeta}\sum_{h\in\zeta}
                  (\alpha n(h)+\sqrt{r n(h)})^2.
\]
For every permitted remainder, the final separation of this policy
satisfies
\begin{equation}\label{eq:policyprofile89}
\begin{split}
 D(\pi;E,F)&\le\min\{2,2s\sqrt{B_\pi}\}\\
 &\le\min\{2,2s(\alpha\sqrt{Q_\pi}+\sqrt{rN_\pi})\}.
\end{split}
\end{equation}
The receiver dimension and all receiver instruments are unrestricted.
The complete fresh acquisition remains conditionally independent of
old receiver memory, including all within-block references and memory.
\end{proposition}
\begin{proof}
At history $h$ the fresh adaptive-block lemma gives a deficit at most
$a_h=s^2(\alpha n(h)+\sqrt{rn(h)})^2/2$.
The term $n(h)rs^2$ is the full tube-to-surrogate channel error before
Bures conversion.  Along every terminal history their sum is at most
$s^2B_\pi/2$.  Conditional fidelity accumulation and
$\|\rho-\sigma\|_1\le2\sqrt{1-f(\rho,\sigma)^2}$ prove the first
bound, with the universal trace cap covering large deficits.
On each path the Euclidean triangle inequality gives
\[
 \left[\sum_h(\alpha n(h)+\sqrt{rn(h)})^2\right]^{1/2}
 \le\alpha\left(\sum_hn(h)^2\right)^{1/2}
                      +\sqrt{r\sum_hn(h)}.
\]
Taking suprema proves the second bound.  No expectation of a path cost
or probability of a history has been substituted for a supremum.
The induction uses positive subnormalized operators and hence applies
also at null histories and to trace-class countable direct sums.
\end{proof}

\begin{definition}[A quadratic reservation budget]\label{def:quadraticbudget89}
For integers $N\ge1$ and $N\le Q\le N^2$, let
$D^{\mathrm{reset}}_{N,Q}(E,F)$ be the supremum over reset protocols
with $N_\pi\le N$ and $Q_\pi\le Q$.  The budgets include every
reserved slot on every formal terminal history, whether or not it is
used internally.  There is no additional width constraint.  The
conditional reset and common-control conventions are those of
Definition~\ref{def:reset88}.
\end{definition}

\begin{theorem}[Sharp quadratic budget law]\label{thm:quadraticbudget89}
Fix $E$, a nonzero Hermitian zero-sum tuple $H$ in the one-sided cone,
and $\Lambda\ge0$.  There exist $c,C,s_0>0$, depending only on
$E,H,\Lambda$, such that for all integers $N\ge1$, $N\le Q\le N^2$,
$0<s\le s_0$, and every legal $F$ with
\[
 \|F-E-sH\|_\Sigma
 :=\sum_j\|F_j-E_j-sH_j\|_{\mathrm{op}}\le\Lambda s^2,
\]
one has $c\omega\le D^{\mathrm{reset}}_{N,Q}(E,F)\le C\omega$, where
\begin{equation}\label{eq:quadraticbudget89}
\omega=
\begin{cases}
 \min\{1,\sqrt N s\},&H\in\operatorname{Ran}\mathcal C_E,\\
 \min\{1,\sqrt Q s\},&Q_jH_jQ_j=0\ (\forall j),\quad
                           \Gamma_E(H)\notin\mathcal W_E,\\
 \min\{1,Ns\},&Q_jH_jQ_j\ne0\ (\exists j).
\end{cases}
\end{equation}
The lower tester is independent of the unknown remainder and can
be chosen with a deterministic profile, no inter-block feedback,
and parallel acquisition within its blocks.  Constants are uniform
in $N,Q,F$, not in the fixed tube.  The theorem asserts local orders,
not exact-distance identities or a classification of memory dimensions.
\end{theorem}
\begin{proof}
For coherent tangents, Proposition~\ref{prop:policyprofile89} gives
\[
 D(\pi;E,F)\le\min\{2,2s(\alpha\sqrt Q+\sqrt{rN})\}
             \le\min\{2,2(\alpha+\sqrt r)s\sqrt Q\}.
\]
For the lower put $b=\lfloor Q/N\rfloor$.  Then $1\le b\le N$,
$Nb\le Q$, and $Nb\ge Q/2$.  Take as many blocks of size $b$ as
possible, followed by the positive remainder when necessary.
Their sum is $N$ and, by Corollary~\ref{cor:profilebudget89},
\[
 Q/4\le Nb/2\le V_{\boldsymbol n}\le Nb\le Q.
\]
The lower construction of Theorem~\ref{thm:allocation89} belongs to
the prescribed budget class and gives the coherent lower.  In
particular, the possible quadratic-budget integrality gaps cost only
an absolute factor; the theorem does not claim to attain every
integer $Q$ exactly.

For regular tangents use the unrestricted adaptive horizontal upper
and the independent-product lower.  The all-singleton profile has
square cost $N\le Q$.  The finite upper still contains
$2\beta_h\sqrt Ns+(\Lambda+K_h)Ns^2$, capped at two; putting
$x=\sqrt Ns$ gives the claimed order.  For support openings the same
singleton profile uses the fixed impossible-label event, and the
ordinary adaptive hybrid bound gives the upper.  This lower is a
classical leakage experiment rather than a coherent phase code.
All lower controls are chosen from the fixed tube and public budgets;
the complete finite $m_r\kappa s^2$ errors were paid in
Theorem~\ref{thm:allocation89} before statistical amplification.
\end{proof}

\begin{proposition}[Certified deviations from reset]\label{prop:approxreset89}
Fix a finite or countably branching reset policy with a hard finite
call bound.  At each acquisition history $h$, replace the ideal
preparation channel $R\mapsto R\otimes\sigma_h$ by a channel
$\widetilde{\mathcal P}_h$ with the same input and output systems.
Assume, including auxiliary references,
\[
 \|\widetilde{\mathcal P}_h-(\operatorname{Id}\otimes\sigma_h)
                      \|_\diamond\le\delta_h,\qquad
 \sup_\zeta\sum_{h\in\zeta}\delta_h\le\varepsilon.
\]
All within-block testers, receiver instruments, and reservation rules
are kept fixed and common.  The implemented final separation obeys
\begin{equation}\label{eq:approxreset89}
 \widetilde D(\pi;E,F)
 \le\min\{2,\ D(\pi;E,F)+2\varepsilon\}.
\end{equation}
For tangential directions one may substitute the right side of
\eqref{eq:policyprofile89} for $D(\pi;E,F)$.
For the same pair of ideal and implemented channels the proof gives the
two-sided estimate $|\widetilde D(\pi;E,F)-D(\pi;E,F)|\le2\varepsilon$.
In particular a lower $c\omega$ remains at least $c\omega/2$ if
$\varepsilon\le c\omega/4$; the condition $\varepsilon=o(\omega)$
suffices along a shrinking-signal sequence.
\end{proposition}
\begin{proof}
For either fixed device, let $\varepsilon_h$ bound the remaining
path sum from a history $h$.  Induct backwards on remaining depth
with arbitrary positive incoming operator $R$.  Replacing the
preparation changes the continued output in trace norm by at most
$\delta_h\tr R$, since every continuation is a channel.  For the
ideal preparation followed by the common block and receiver
instrument, the positive child operators have traces summing to
$\tr R$.  The induction bounds the remaining difference by
$(\varepsilon_h-\delta_h)\tr R$.  Their sum is
$\varepsilon_h\tr R$.  At terminals it is zero; countable sums
converge monotonically in trace.  Thus each hypothesis's normalized
output differs by at most $\varepsilon$.  Applying the triangle
inequality to the two hypotheses proves the result.
\end{proof}

The hypothesis is a uniform channel certificate for a specified
common ideal reset preparation.  Marginal separability, finite
policy arithmetic, or an average calibration error does not supply
it.  Fixed nonzero $\varepsilon$ can dominate a shrinking local
signal; no fixed-error asymptotic robustness is inferred.  The same
argument applies to the particular lower construction when its
preparations satisfy the displayed hypothesis, giving the lower
separation minus $2\varepsilon$.

\subsection{A strict tester-level comparison with classical adaptation}
We now give an explicit comparison with the classically adaptive
model of Ohst--Zhang--Nguyen--Pl\'avala--Quintino
\cite[Definition~23 and (64)]{OZNPQ89}.  For two calls their tester
operators have the form
\begin{equation}\label{eq:catester89}
 T_i=\sum_z R_z\otimes S_{i|z},
\end{equation}
where $R$ and each $S_{\cdot|z}$ are single-call testers.  This is
an entire first-call measurement followed by conditional new
preparation, not retention of an unmeasured first-call quantum
register for a later joint readout.

\begin{proposition}[Retained receiver memory]\label{prop:memorycomparison89}
Every two-call tester of the form \eqref{eq:catester89} is a unit-width
reset tester.  The inclusion is strict at the level of tester
operators, already for two-dimensional inputs and two classical
output labels.  This is not an assertion of a strict optimal-score
gap for every ensemble or for every pair of identical measurement
channels.
\end{proposition}
\begin{proof}
Realize $R$ on the first call, retain only its classical outcome
$z$, then prepare and realize $S_{\cdot|z}$ freshly.  This obeys the
unit-width reset condition and proves inclusion.

For strictness, write $I_1,I_2$ for the two qubit inputs and $Y_1,Y_2$
for their classical output registers.  Prepare independent maximally
entangled states on $I_rR_r$, one per call.  Keep the first reference
and label in the receiver, append the second complete pair freshly,
and make a final joint measurement.  One effect of that measurement
is $\Pi_{R_1R_2}\otimes|00\rangle\langle00|_{Y_1Y_2}$, where
$\Pi=|\Phi^+\rangle\langle\Phi^+|$ and
$|\Phi^+\rangle=(|00\rangle+|11\rangle)/\sqrt2$.
For precision, order factors as $I_1,Y_1,I_2,Y_2$ and take the fully
transposed unnormalized Choi representative
$J^\sharp(E)=\sum_yE_y\otimes|y\rangle\langle y|$, obtained from
$(\operatorname{Id}\otimes\mathcal M_E)(|\Omega\rangle\langle\Omega|)$
by full transpose, with $|\Omega\rangle=|00\rangle+|11\rangle$.
Thus $\tr_YJ^\sharp=I$.  The two normalized pairs produce the receiver
operator $(J^\sharp(E)\otimes J^\sharp(F))^{\mathsf T}/4$.
For any receiver effect $M$, the Born probability is
$\tr[M^{\mathsf T}(J^\sharp(E)\otimes J^\sharp(F))]/4$.
This proves the transpose and factor $1/4$ for arbitrary complex effects,
not merely for the real Bell projector.  Grouping the input factors below
means the canonical permutation of the stated factor order.
In the Choi convention with input marginal $I$, direct contraction
of the two normalized maximally entangled preparations gives the
tester effect
\begin{equation}\label{eq:memorywitness89}
 T_0=\tfrac14\Pi^{\mathsf T}_{I_1I_2}
                      \otimes|00\rangle\langle00|_{Y_1Y_2},
 \qquad T_1=\tfrac14 I_{I_1I_2}\otimes I_{Y_1Y_2}-T_0.
\end{equation}
Both are positive and their sum is the deterministic parallel tester
with input marginal $I_{I_1I_2}/4$.  They are block diagonal in the
classical outputs.  The described realization has reset width one:
no old reference enters the second acquisition.
Indeed $0\preceq\Pi^{\mathsf T}\otimes|00\rangle\langle00|\preceq I$,
so the complementary effect is positive.  The deterministic sum is
$\Xi\otimes I_{Y_2}$ with
$\Xi=I_{I_1}\otimes I_{Y_1}\otimes I_{I_2}/4$ and
$\tr_{I_2}\Xi=(I_{I_1}/2)\otimes I_{Y_1}$; the initial marginal has
trace one.  This verifies both causal tester trace constraints.
The cone of positive separable operators across the complete-call cut
$I_1Y_1\mid I_2Y_2$ is closed: its trace-one slice is the compact convex
hull of the compact set of product states, and trace controls the scaling
in its cone.  Hence the limiting sums mentioned below lie in this same
closed cone.  Partial transpose is used only as a necessary condition
for separability.

Partial transpose on the complete first-call factor $I_1Y_1$ leaves
the output projector unchanged.  Since $\Pi$ is real and its partial
transpose has eigenvalues $1/2,1/2,1/2,-1/2$, the partial transpose of
$T_0$ has an eigenvalue $-1/8$.  Every positive sum
\eqref{eq:catester89} has positive partial transpose across
$I_1Y_1|I_2Y_2$, including limits of such sums.  Hence this effect
cannot occur in a tester of that form.  This proves strictness.
\end{proof}

Consequently the constrained-separability condition on complete
first- and second-call tester factors in \cite[(65)--(68)]{OZNPQ89}
cannot be imposed on the whole unit-reset class.  Reset independence
is a factorization at a specified conditional preparation cut; it
does not imply separability of the eventual tester across call
factors.  The witness is a calibration of model inclusions, not a
new memory-discrimination mechanism or a physical experiment.

\section{A variational principle for two-call reset experiments}
\label{sec:resetvariational91}
The fresh-preparation cut admits an exact finite-dimensional description
of optimal decision values.  No symmetry, second-moment identity, or
special choice of prior is needed in this section.  The description is
by a concave envelope over density matrices, not over pure states alone.
Its finite realization bounds the receiver needed to attain a value;
it does not identify reset width with a prescribed memory dimension.

\subsection{Experiments and strategy classes}
Let $I_\ell=\C^{d_\ell}$ and let $Y_\ell$ be a finite classical alphabet,
$\ell=1,2$.  A finite experiment consists of a probability law
$(p_\theta)$, two ordered measurements $E^{\theta,\ell}$ on $I_\ell$,
and rewards $r_{j\theta}\in[0,1]$ for decisions $1\le j\le m$.
The index $\theta$ is sampled once.  Conditional on it the two slots are
the memoryless measurement channels $\mathcal M_{E^{\theta,\ell}}$.
Repeating one unknown measurement is the case
$E^{\theta,1}=E^{\theta,2}$.  Identification and binary decisions on
subfamilies are obtained by taking indicator rewards.

\begin{definition}[Two-call classes]\label{def:classes91}
An unrestricted strategy prepares a state on $I_1\otimes R$, observes
the first label $y$, applies a common trace-preserving map from $R$ to
$I_2\otimes R'$, and, after observing $z$, measures $R'$ to decide $j$.
A unit-reset strategy instead applies a normal instrument
$(\mathcal K_{yh})_h$ to $R$, retains its old quantum output $O_{yh}$,
and prepares a state $\sigma_{yh}$ on the complete fresh system
$I_2\otimes R_{2,yh}$.  Conditional on $(y,h)$ the preparation is
exactly the tensor product of the old output with $\sigma_{yh}$.
No old quantum register enters $I_2$ or its fresh reference.  After the
second label, the two receivers may be measured jointly.  A classically
adaptive strategy is a unit-reset strategy that measures the entire old
receiver and retains only a classical record before fresh preparation.

All rules are independent of $\theta$ and defined on every formal
history, including histories null under a hypothesis.  Receiver spaces
have no dimension bound.  Instruments may have countably many outcomes;
all classical records are retained.  Finite or countable public
randomization and the
finite-dimensional trace-norm closure of the resulting output-dephased
Choi tester tuples are included in each class.  General public seed laws
are included through this closure: their finitely many tester coordinates
are limits of finite convex combinations.  Early stopping is
implemented by a fixed, discarded fresh call and a decision independent
of its label; it is not a zero-length component of a prescribed profile.
Write $P_{\rm all},P_{\rm reset},P_{\rm ca}$ for the suprema of expected
reward in these classes.
\end{definition}
This reset cut is the two-call instance of Definition~\ref{def:reset88}.
The first class in the ordering
$P_{\rm ca}\le P_{\rm reset}\le P_{\rm all}$ agrees with complete
measure-and-reprepare adaptation, not with a constraint of separability
on all final reset effects.  In the order $I_1,Y_1,I_2,Y_2$, a
classically adaptive effect is a sum of positive products across
$I_1Y_1\mid I_2Y_2$ \cite[Definition~23]{OZNPQ89}.
The positive separable cone is closed in these finite dimensions:
its trace-one slice is the compact convex hull of product states, and
trace controls passage to the cone.  Only this necessary condition,
not its sufficiency for classical adaptive normalization, is used below.

\begin{lemma}[Countable records]\label{lem:countable91}
Let $(X_h)$ be positive trace-class operators with
$\sum_h\tr X_h<\infty$.  Their partial sums converge in trace norm,
and their recorded direct sum lies in the trace-class $\ell^1$ sum.
For effects $0\preceq F_h\preceq I$, the corresponding probability
series converges absolutely.  Replacing all but finitely many branches
of a normal instrument by one remainder branch preserves normalization;
changing the continuation on that branch changes any reward in $[0,1]$
by at most its probability.
\end{lemma}
\begin{proof}
The norm of a positive tail equals its trace.  This proves the first
assertions and bounds the probability tail by $\sum_{h>n}\tr X_h$.
For an instrument, the sum of the omitted completely positive maps is
the remainder map; its trace is precisely the omitted outcome
probability.  Apply the bound to each member of the finite experiment
and then average.  The maximum of its finitely many tail probabilities
tends to zero.  No conditional state is divided by a history probability.
\end{proof}

\subsection{A complete branch normal form}
Fix bases on $I_1,I_2$.  A coefficient map $C:I_1\to R$ represents
$|C\rangle=\sum_i|i\rangle\otimes C|i\rangle$; its squared norm is
$\tr C^*C$.  Measuring its input gives the unnormalized receiver state
$CE_y^{\mathsf T}C^*$.

\begin{lemma}[Unit-reset normal form and converse]\label{lem:resetnormal91}
Every implemented unit-reset strategy has a purification, still respecting the reset
cut, from which its original decision statistics are recoverable by
common receiver channels.  For this purification there are rectangular
maps
\[
 C_0:I_1\longrightarrow R_0,\quad
 C_{yh}:I_1\longrightarrow O_{yh},\quad
 D_{yh}:I_2\longrightarrow R_{2,yh}
\]
with, on the input copies,
\begin{equation}\label{eq:resetgrams91}
 \rho=C_0^*C_0,\quad \tr\rho=1,\quad
 A_{yh}=C_{yh}^*C_{yh},\quad \sum_hA_{yh}=\rho,\quad
 b_{yh}=D_{yh}^*D_{yh},\quad\tr b_{yh}=1.
\end{equation}
The receiver block on the formal record $(y,h,z)$ is
\begin{equation}\label{eq:resetblock91}
 (C_{yh}E_y^{\theta,1\,\mathsf T}C_{yh}^*)
 \otimes(D_{yh}E_z^{\theta,2\,\mathsf T}D_{yh}^*).
\end{equation}
The sum in \eqref{eq:resetgrams91} converges in trace norm.
Conversely, every finite collection satisfying \eqref{eq:resetgrams91}
is realizable, up to common receiver isometries, by a unit-reset strategy
with old receiver dimension at most $d_1$ and fresh reference dimension
at most $d_2$.  Zero-probability histories require no extra assumption.
The tester closure in Definition~\ref{def:classes91} is used by continuity
of decision values; no physical dilation is assumed for an abstract
limit before its value has been realized by the finite optimization.
\end{lemma}
\begin{proof}
Purify the first preparation, retaining the purifying system in $R_0$.
The Schmidt support has dimension at most $d_1$.  For each $y$, dilate
each branch map of the intervening instrument to a single map
$V_{yh}:R_0\to O_{yh}$, retaining its environment as old receiver
memory.  On the relevant support,
$\sum_hV_{yh}^*V_{yh}=I$ in the strong sense.  Thus
$C_{yh}=V_{yh}C_0$ has
$\sum_h C_{yh}^*C_{yh}=C_0^*C_0$.
The positive tails converge in trace norm by Lemma~\ref{lem:countable91}.
This argument uses only the finite Schmidt support even when the
original receiver or a dilation environment is infinite-dimensional.
Tracing those environments recovers the original instrument.

Purify the entire fresh preparation separately.  Its coefficient map
is $D_{yh}$ with squared Hilbert--Schmidt norm one.  The tensor-product
preparation rule, not a statement about a probe marginal, gives
\eqref{eq:resetblock91}.  Common pre-query controls are absorbed into
these preparations and common post-query controls into the final
receiver measurement.  Fresh mixing may be purified inside the fresh
reference.  A retained public seed may be put in the classical record;
randomization before the first call is a mixture of such strategies and
cannot increase the supremum of a linear reward.  More explicitly,
for seed law $q_s$ and first coefficient maps $C_{0,s}$, take
$C_0=\bigoplus_s\sqrt{q_s}C_{0,s}$ and measure the seed register
as part of the first receiver instrument.  The branch maps are
$C_{y,(s,h)}=\sqrt{q_s}V_{ysh}C_{0,s}$, and their Gram sum is
$\sum_sq_sC_{0,s}^*C_{0,s}=C_0^*C_0$ for every $y$.
The fresh rule retains its dependence on $(s,y,h)$, entirely on the
fresh side of the cut.  Thus the original seeded decision statistics
are recovered.  Countable records and tester limits follow from the
preceding lemma and continuity.

For the converse let $r=\operatorname{rank}\rho$ and choose
$C_0:I_1\to\C^r$ with $C_0^*C_0=\rho$.  Its restriction to
$\operatorname{supp}\rho$ is invertible onto $\C^r$.  Let $C_0^+$
be the inverse on this support, and set
\[
 C_{yh}=\sqrt{A_{yh}},\qquad
 V_{yh}=\sqrt{A_{yh}}C_0^+,\qquad D_{yh}=\sqrt{b_{yh}}.
\]
Positivity and $\sum_hA_{yh}=\rho$ imply
$\operatorname{supp}A_{yh}\subseteq\operatorname{supp}\rho$.
Hence $V_{yh}C_0=\sqrt{A_{yh}}$ and
$\sum_hV_{yh}^*V_{yh}=I_r$.  These are Kraus operators of a
trace-preserving recorded instrument with outputs in $\C^{d_1}$.
The fresh vector $|\sqrt{b_{yh}}\rangle$ is normalized and independent
of that output.  Polar decomposition identifies any other rectangular
maps of the same Gram matrices with these maps on their supports;
receiver isometries may be absorbed into the final effects.
Zero Gram matrices give zero branch maps.  All other null histories are
handled by the same unnormalized formulas, including projective devices.
\end{proof}

\subsection{The compact primal problem}
Let $\mathcal D_d$ denote the density matrices on $\C^d$, and put
\begin{equation}\label{eq:gameblocks91}
 G_{j,yz}=\sum_\theta p_\theta r_{j\theta}
     (E_y^{\theta,1})^{\mathsf T}\otimes(E_z^{\theta,2})^{\mathsf T}.
\end{equation}
These are fixed positive matrices on the two input copies.  For
$A\succeq0$ and $b\in\mathcal D_{d_2}$ define
\begin{equation}\label{eq:leafvalue91}
 \Phi_y(A,b)=\sum_z\max_{\substack{F_j\succeq0\\\sum_jF_j=I}}
 \sum_j\tr\bigl[F_j(\sqrt A\otimes\sqrt b)
            G_{j,yz}(\sqrt A\otimes\sqrt b)\bigr].
\end{equation}
The POVM is chosen separately at every recorded $z$.  It acts on
$\C^{d_1}\otimes\C^{d_2}$.  For $a\in\mathcal D_{d_1}$ let
$g_y(a)=\max_{b\in\mathcal D_{d_2}}\Phi_y(a,b)$ and let
\begin{equation}\label{eq:roof91}
 c_y(\rho)=\sup\left\{\sum_h\lambda_h g_y(a_h):
 \lambda_h\ge0,\ \sum_h\lambda_h=1,\
 \sum_h\lambda_ha_h=\rho\right\}.
\end{equation}
The envelope is taken over \emph{all} density matrices $a_h$.  Replacing
them by rank-one states would change the resource problem.  The
construction of concave envelopes is standard \cite{Uhlmann91}; the
operational content here is the reset realization and the common
barycenter constraint across all first labels.

\begin{theorem}[Exact reset variational principle]\label{thm:resetvariational91}
For every finite experiment of Definition~\ref{def:classes91},
\begin{equation}\label{eq:resetvariational91}
 P_{\rm reset}=\max_{\rho\in\mathcal D_{d_1}}\sum_yc_y(\rho)
 =\max_{\substack{A_{yh}\succeq0,\ \sum_hA_{yh}=\rho\\
                    \rho\in\mathcal D_{d_1},\ b_{yh}\in\mathcal D_{d_2}}}
                    \sum_{y,h}\Phi_y(A_{yh},b_{yh}).
\end{equation}
All maxima are attained.  For a fixed $\rho$ of rank $r$, at most $r^2$
nonzero instrument outcomes are required after each first label.
An optimal strategy exists with old receiver dimension at most $d_1$,
fresh reference dimension at most $d_2$, and a final quantum receiver
of dimension at most $d_1d_2$, separately from the finite classical
record.  The dimension bounds concern the retained registers at the recording
cuts, not temporary workspaces for implementing an instrument.
These assertions include arbitrary priors and arbitrary fixed
pairs as particular decision problems; they assert an exact optimum,
not a strict advantage for every such problem.
\end{theorem}
\begin{proof}
For fixed $y$, $\Phi_y$ is continuous, and is positively homogeneous
in $A$.  Continuity follows from continuity of matrix square roots and
maximization of a continuous function over a compact POVM set.
Maximizing also over compact $\mathcal D_{d_2}$ makes $g_y$ continuous.
The graph $\{(a,g_y(a)):a\in\mathcal D_{d_1}\}$ is compact in a
real space of affine dimension at most $d_1^2$.  Its convex hull is
compact.  The largest height above $\rho$ is therefore attained and
is exactly $c_y(\rho)$, also allowing countable decompositions.
It is an upper semicontinuous concave function of $\rho$.
Consequently $\sum_yc_y$ attains a maximum.

Carath\'eodory first supplies at most $r^2+1$ atoms at a maximizing
height over a rank-$r$ matrix: each atom is supported on
$\operatorname{supp}\rho$, since positive matrices cannot cancel on its
orthogonal complement.  If more than $r^2$ weights are nonzero, the
Hermitian atoms are linearly dependent.  Choose nonzero real $v_h$ with
$\sum_hv_ha_h=0$; tracing gives $\sum_hv_h=0$.
Small perturbations $\lambda_h\pm\epsilon v_h$ preserve the barycenter
and positivity.  Optimality forces $\sum_hv_hg_y(a_h)=0$, since otherwise
one sign raises the height.  Move in this direction until a weight
vanishes.  The value is unchanged.  Iteration gives at most $r^2$ atoms.
For each atom the maximum over $b$ and the leaf POVMs is attained.

By Lemma~\ref{lem:resetnormal91}, the value of any purified competitor
is bounded by the right side of \eqref{eq:resetvariational91}, with
$A_{yh}=\lambda_{yh}a_{yh}$.  Countable records follow by the positive
tail bound; public mixtures and norm limits preserve the upper.
Conversely, the lemma realizes every finite maximizing collection,
and its optimal leaf POVMs realize every maximum in
\eqref{eq:leafvalue91}.  This proves equality and the dimension bounds.
The final quantum dimension excludes the separately recorded labels.
\end{proof}
For a binary equal-prior fixed pair, its optimal unhalved final trace
separation is $4P_{\rm reset}-2$.  Thus the theorem is also a general
variational formula for that operational distance.  The primal has
finitely many bounded matrix variables, but is not asserted to be a
convex semidefinite program or an efficient synthesis algorithm.

\subsection{Dual certificates and the role of receiver feedback}
\begin{theorem}[Duality and contact conditions]\label{thm:resetdual91}
For the same experiment,
\begin{equation}\label{eq:resetdual91}
 P_{\rm reset}=\min_{(H_y)}\lambda_{\max}\left(\sum_yH_y\right),
 \qquad \tr(H_ya)\ge g_y(a)\quad(a\in\mathcal D_{d_1}),
\end{equation}
where $H_y$ are Hermitian $d_1\times d_1$ matrices.  The minimum is
attained.  It is equivalent to minimizing $\lambda$ subject to the
same majorizations and $\sum_yH_y=\lambda I$.
For any optimal primal and dual, every atom of positive weight is a
contact point, $\tr(H_ya_{yh})=g_y(a_{yh})$, and $\rho$ is supported
on the maximal eigenspace of $\sum_yH_y$.  Conversely these conditions,
with maximizing fresh states and leaf POVMs, certify optimality.
\end{theorem}
\begin{proof}
Every majorization of $g_y$ majorizes its concave envelope.  Therefore
$\sum_yc_y(\rho)\le\tr(\rho\sum_yH_y)\le
\lambda_{\max}(\sum_yH_y)$, which proves weak duality.
For completeness, the strong-duality step is finite-dimensional.
Extend $c_y$ homogeneously to positive matrices by
$\widehat c_y(A)=\tr A\,c_y(A/\tr A)$ and $\widehat c_y(0)=0$.
Its hypograph
\[
 \mathcal K_y=\{(A,u):A\succeq0,\ u\le\widehat c_y(A)\}
\]
is a closed convex cone.  Closedness at $A=0$ follows from boundedness
of $g_y$; elsewhere it follows from the compact-graph construction.
The conic problem maximizing $\sum_yu_y$, with $(A_y,u_y)\in\mathcal K_y$,
$A_y=\rho$ and $\tr\rho=1$, is precisely the primal problem.
The point $A_y=\rho=I/d_1$, $u_y=-1$ is strictly interior to all
hypographs, since rewards are nonnegative and $A_y$ is positive definite.
Its value is finite, in fact at most one by the operational realization.
The finite-dimensional separating-hyperplane form of conic Slater
duality thus gives an attained dual with no gap
\cite[Sections~5.2.3 and~5.9.1]{BV91}.

Its multipliers can be read directly: the supremum of
$u-\tr(H_yA)$ on $\mathcal K_y$ is zero exactly when
$\tr(H_ya)\ge g_y(a)$ for every density matrix.  Taking $\rho$ as a
free Hermitian variable under its trace constraint imposes
$\sum_yH_y=\lambda I$.  This gives the stated dual.  Conversely, for
majorants with sum $S$, add the positive matrix
$\lambda_{\max}(S)I-S$ to one $H_y$; all majorizations remain valid
and the sum becomes scalar.  The two forms have the same minimum.
Finally, the weak-duality difference is a sum of nonnegative atom
majorization defects and a nonnegative spectral defect.  It vanishes
exactly under the stated contact and support conditions.
\end{proof}
The majorizations in \eqref{eq:resetdual91} are universal inequalities
on density matrices.  A finite sampling of them is not a dual certificate.

\begin{corollary}[Quantitative contact certificates]\label{cor:contactdefect91}
Let $(H_y)$ be any feasible all-density majorants in
Theorem~\ref{thm:resetdual91}, put $S_H=\sum_yH_y$ and
$U=\lambda_{\max}(S_H)$, and consider a finite or countable reset
normal form with achieved reward $v$.  For every nonzero branch write
$w_{yh}=\tr A_{yh}$, $a_{yh}=A_{yh}/w_{yh}$, and let $v_{yh}$
be its contribution to the reward, summed over the second labels and
decisions.  Then
\begin{equation}\label{eq:contactdefect91}
\begin{split}
 U-v={}&\tr[\rho(UI-S_H)]\\
 &+\sum_{y,h}w_{yh}\bigl[\tr(H_ya_{yh})-g_y(a_{yh})\bigr]\\
 &+\sum_{y,h}\bigl[w_{yh}g_y(a_{yh})-v_{yh}\bigr].
\end{split}
\end{equation}
Every term is nonnegative.  Consequently, if $U-v\le\varepsilon$,
the total Gram weight of branches with majorization defect at least
$u>0$ is at most $\varepsilon/u$.  If the maximal eigenspace projection
$P$ of $S_H$ has a spectral gap $\gamma>0$ to its complement, then
$\tr[(I-P)\rho]\le\varepsilon/\gamma$.
\end{corollary}
\begin{proof}
The common-barycenter constraints give
$\sum_{y,h}w_{yh}\tr(H_ya_{yh})=\tr(\rho S_H)$.
Add and subtract this quantity and $\sum_{y,h}w_{yh}g_y(a_{yh})$.
The leaf optimization bounds $v_{yh}$ by $w_{yh}g_y(a_{yh})$.
Feasible majorization and $S_H\preceq UI$ prove nonnegativity.
All series converge absolutely: the first-label alphabet is finite,
$\sum_hw_{yh}=1$ for every $y$, and the functions and matrices are
bounded.  Restricting a nonnegative sum proves the first estimate;
$UI-S_H\succeq\gamma(I-P)$ proves the second.
\end{proof}
The weights in this certificate are input Gram weights, not the
hypothesis-dependent probabilities of recorded histories.  With an
attained optimal dual, $U-v$ is precisely the loss in optimal reward.
The formula requires a genuine universal majorant, not sampled tests
of its inequality.

\begin{corollary}[Label feedback versus receiver feedback]\label{cor:feedbackroof91}
If no receiver-instrument outcome is used to choose the second
preparation, but the first device label may be used, the optimum is
\[
 P_{\rm label}=\max_{\rho\in\mathcal D_{d_1}}\sum_yg_y(\rho).
\]
Allowing receiver feedback replaces $g_y$ by $c_y$ in this formula.
In particular, concavity of every $g_y$ is sufficient for
$P_{\rm label}=P_{\rm reset}$.  More generally, a label-only strategy
with a common density $\rho$ is optimal for the full reset class if
it attains the contact and spectral conditions of
Theorem~\ref{thm:resetdual91}.
\end{corollary}
\begin{proof}
With one preparation choice after each $y$, the normal form has
$A_y=\rho$; a trace-preserving receiver operation cannot improve a
subsequent unrestricted joint measurement and can be postponed to it.
Conversely these choices are admissible.  The assertion follows from
the primal formula and the definition of concave envelope.
\end{proof}

\subsection{A geometric criterion for retained receiver advantage}
Let $\mathcal P_{d_1}$ be the rank-one density matrices.  In parallel
with \eqref{eq:roof91}, define
\begin{equation}\label{eq:pureroot91}
 c_y^{\rm cl}(\rho)=\sup\left\{\sum_h\lambda_h g_y(a_h):
 a_h\in\mathcal P_{d_1},\quad
 \lambda_h\ge0,\quad\sum_h\lambda_h a_h=\rho\right\}.
\end{equation}
The weights sum to one by taking traces.  Every density matrix admits
such a decomposition, but the value of this pure-atom envelope need
not equal the all-density envelope $c_y$.

\begin{theorem}[Classical and quantum receiver envelopes]
\label{thm:classicalroof91}
For every finite experiment of Definition~\ref{def:classes91},
\begin{equation}\label{eq:classicalroof91}
 P_{\rm ca}=\max_{\rho\in\mathcal D_{d_1}}\sum_yc_y^{\rm cl}(\rho)
 =\min_{(H_y)}\lambda_{\max}\left(\sum_yH_y\right),
 \qquad \tr(H_ya)\ge g_y(a)\quad(a\in\mathcal P_{d_1}).
\end{equation}
Both optima are attained.  At most
$(\operatorname{rank}\rho)^2$ nonzero first-receiver outcomes per
first device label suffice at an optimal $\rho$.
The passage from complete classical readout to retained quantum
receiver memory replaces the pure-atom envelopes $c_y^{\rm cl}$ by
the all-density envelopes $c_y$.  In particular, the invariant
\begin{equation}\label{eq:receivergap91}
 \Gamma=\max_\rho\sum_yc_y(\rho)
          -\max_\rho\sum_yc_y^{\rm cl}(\rho)
\end{equation}
is exactly the optimal retained-receiver advantage in that experiment.
\end{theorem}
\begin{proof}
After complete first readout, a classical branch has a positive input
Gram effect $A$ and retains no quantum system depending on $\theta$.
Its probability under a first effect is $\tr(AE_y^{\mathsf T})$.
Refine a spectral decomposition $A=\sum_i\lambda_i a_i$ into the
classical record and keep the former second preparation and decision
on each refined branch.  This does not change the score and gives a
rank-one decomposition with the same barycenter.  Optimizing the
second action separately after the refined record can only increase it.

Conversely, for a rank-one atom $a=|u\rangle\langle u|$, the old
receiver in the canonical normal form has
\[
 \sqrt a E_y^{\mathsf T}\sqrt a
 =\langle u,E_y^{\mathsf T}u\rangle a.
\]
Its normalized state is independent of the unknown device.  A joint
POVM with this known factor is equivalent to a POVM on the fresh
reference alone, by taking its expectation against $a$.
Thus $g_y(a)$ is attainable with only a classical first record.
The converse part of Lemma~\ref{lem:resetnormal91}, followed by discarding
this known old state, realizes every finite pure-atom decomposition.
These two observations give the first equality in
\eqref{eq:classicalroof91}.

The graph of $g_y$ on the compact pure-state set is compact.  The
compact-convex-hull and rank-face atom-elimination arguments of
Theorem~\ref{thm:resetvariational91} therefore apply verbatim, giving
attainment and the stated branch count.  Its homogeneous hypograph is
again a closed convex cone, with a strictly feasible point at a
positive definite common barycenter and negative heights.  The proof
of Theorem~\ref{thm:resetdual91} applies with majorization required only
on pure atoms and gives the attained dual.  Majorization on pure atoms
is equivalent to majorization of $c_y^{\rm cl}$ on every density matrix,
not to majorization of $g_y$ on every density matrix.  The last assertion
now follows from the two exact primal formulas.
\end{proof}

\begin{corollary}[Complete separation certificates]\label{cor:gapcertificate91}
There is a strict retained-receiver advantage if and only if there exist
a finite feasible reset collection from \eqref{eq:resetvariational91},
with achieved reward $v$, and Hermitian matrices $H_y$ such that
\[
 \tr(H_ya)\ge g_y(a)\quad(a\in\mathcal P_{d_1}),\qquad
 v>\lambda_{\max}\left(\sum_yH_y\right).
\]
There is no advantage if and only if an attaining classical strategy
has contact and spectral equality with an all-density dual majorant
of Theorem~\ref{thm:resetdual91}.
\end{corollary}
\begin{proof}
The displayed inequalities compare a physically attained reset lower
with a valid classical upper.  If the gap is positive, take an optimal
finite reset collection and an optimal classical dual, whose existence
was proved above.  This proves both directions of the first assertion.
If the gap is zero, an optimal classical strategy also attains the reset
value; complementarity with any attained reset dual gives the second
condition.  Conversely, that condition certifies a common upper and
lower for the two classes, so their values agree.
\end{proof}
These are certificates with finitely many matrix variables and universal
majorization requirements, not finite-sampling tests.  The criterion
applies to any fixed-pair problem as well as to ensemble decisions,
but does not require every experiment to have positive $\Gamma$.
For example, the equal-prior qubit calculation below gives pure-atom
majorants $H_0=H_1=(7/24)I$, whereas a reset strategy attains $5/8$.
The gap is $1/24$.  Its universal pure-atom majorization follows from
the classical upper in Theorem~\ref{thm:equalprior91}, not from the
finite arithmetic regression.

\subsection{The unrestricted causal normalization}
\begin{proposition}[Physical two-call tester blocks]\label{prop:causalnormal91}
In the unnormalized convention
$J^\sharp(E)=\sum_yE_y\otimes|y\rangle\langle y|$, every unrestricted
strategy has positive blocks $T_{j,yz}$ and positive matrices $\Xi_y$ with
\begin{equation}\label{eq:causalnormal91}
 \sum_jT_{j,yz}=\Xi_y,\qquad
 \tr_{I_2}\Xi_y=\rho_1,\qquad\tr\rho_1=1.
\end{equation}
Its probability of decision $j$ on the output block $(y,z)$ is
$\tr[T_{j,yz}(E_y^{\theta,1}\otimes E_z^{\theta,2})]$.
In particular $0\preceq T_{j,yz}\preceq\Xi_y$ and $\tr\Xi_y=1$.
\end{proposition}
\begin{proof}
Purify the first preparation to $|\psi\rangle\in I_1\otimes R$.
For its common intervening channel $\mathcal A_y:R\to I_2\otimes R'$
put
\[
 X_y=(\operatorname{Id}\otimes\mathcal A_y)(|\psi\rangle\langle\psi|),
 \quad\Xi_y=\tr_{R'}X_y,\quad\rho_1=\tr_R|\psi\rangle\langle\psi|.
\]
For the final POVM $(F_{j,yz})_j$ on $R'$ define
\[
 T_{j,yz}=\tr_{R'}[(I\otimes\sqrt{F_{j,yz}})X_y
                              (I\otimes\sqrt{F_{j,yz}})].
\]
Positivity is immediate.  Cyclicity in the traced factor and
$\sum_jF_{j,yz}=I$ give the sum; trace preservation of $\mathcal A_y$
gives the marginal.  Expanding the Born rule gives the contraction.
Here $\rho_1$ is the physical input marginal, whereas the coefficient
Gram matrix $\rho$ in \eqref{eq:resetgrams91} is its transpose.
These two conventions must not be interchanged.  Every complementary
effect is positive by the same construction.  No factor $d_1d_2$
occurs, because $J^\sharp$ uses the unnormalized maximally entangled
vector.  Countable records and limits follow by continuity.
\end{proof}

\section{An exact operational separation of acquisition and receiver memory}
\label{sec:operationalmemory90}
The local profile law does not identify exact operational distances.
We now separate three strategy classes on one finite ensemble of ordered
measurement channels.  The separation holds in every input dimension,
throughout a depolarizing family, with exact optimal success probabilities.
The unknown interface still has only classical output.  The experiment
has a binary hypothesis decision, not a request to identify every member
of the ensemble.

\subsection{A finite family and its second moments}
Let $d\ge2$, let $S$ interchange the two factors of $\C^d\otimes\C^d$,
and put $\Pi_\pm=(I\pm S)/2$.  A weighted family of ordered orthonormal
bases consists of projections $P_y^b$, $1\le y\le d$, and positive
weights $w_b$ summing to one.  We require
The weights are real; rational or algebraic coordinates are not required
by the theorem.  We call this a weighted exact second-moment family,
not necessarily an unweighted projective design.
\begin{equation}\label{eq:basismoments90}
 \sum_b w_b P_y^b\otimes P_z^b=
 \begin{cases}
 (I+S)/(d(d+1)),&y=z,\\
 (dI-S)/(d(d^2-1)),&y\ne z.
 \end{cases}
\end{equation}
\begin{lemma}[Finite simultaneous second moments]\label{lem:finitedesign90}
In every dimension $d\ge2$ there is such a family with at most $d^6+1$
bases.  For $d=2$, the six ordered eigenbases of the three Pauli matrices,
with both orders and equal weights, satisfy \eqref{eq:basismoments90}.
\end{lemma}
\begin{proof}
Average the ordered coordinate basis over Haar measure on $U(d)$.
The second moment of a single column is supported on the symmetric
subspace, invariant under $U\otimes U$, and has trace one.  It is
therefore $2\Pi_+/[d(d+1)]$.  For distinct columns the invariant moment
has the form $aI+bS$; its trace is one and its contraction with $S$ is
zero, since the columns are orthogonal.  These two equations give the
second line of \eqref{eq:basismoments90}.  Equivalently, this line follows
by summing over the other $d-1$ columns and using the first moment $I/d$.
The vector of all $d^2$ Hermitian second moments lies in a real space of
dimension at most $d^6$.  Its Haar average belongs to the convex hull
of the compact set of these vectors.  Carath\'eodory's theorem gives a
convex combination of at most $d^6+1$ vectors; discard zero weights.
For the stated qubit family expand $P=(I\pm\sigma_a)/2$ and use
$S=(I\otimes I+\sum_{a=1}^3\sigma_a\otimes\sigma_a)/2$.
\end{proof}
This is a finite existence argument, not an efficient design construction.
Second-moment designs themselves are established objects \cite{GAE90}.

Fix $0\le t\le1$, and consider the ordered measurements
\begin{equation}\label{eq:gamefamily90}
 C_y=I/d,\qquad E_y^b(t)=(1-t)I/d+tP_y^b.
\end{equation}
Set
\begin{equation}\label{eq:gameprior90}
 L=2(d+1)-t^2,\qquad
 p_C=\frac{d+1-t^2}{L},\qquad p_E=\frac{d+1}{L}.
\end{equation}
With probability $p_C$ the device is $C$; otherwise a basis $b$ is chosen
with law $w_b$ and the device is $E^b(t)$.  This choice is made once.
Both calls use the same chosen, memoryless, classical-output device.
The task is to decide $C$ versus the alternative family.  In particular,
the second moments cannot be replaced by two independent first moments.
Every effect is strictly positive for $t<1$.
Taking the partial trace in \eqref{eq:basismoments90} gives
$\sum_b w_bE_y^b(t)=C_y$.  Thus the two hypotheses have identical
one-call channels, including on entangled inputs, and the one-call
optimum is $p_E$.  The two-call signal concerns a device chosen once;
it is absent if the alternative is independently redrawn at each call.

We compare: (i) complete first-call measurement followed by
measure-and-reprepare classical adaptation, as in
\cite[Definition~23]{OZNPQ89}; (ii) unit-reset protocols of
Definition~\ref{def:reset88}, with profile $(1,1)$ and arbitrary receiver
quantum memory and classical feedback; and (iii) all two-call adaptive
protocols.  Let their optimal Bayes success probabilities be
$P_{\rm ca}$, $P_{\rm reset}$ and $P_{\rm all}$, respectively.
Early stopping is allowed and can be implemented by a discarded fresh call.
The three closed optimization classes are precisely those of
Definition~\ref{def:classes91}.  Lemma~\ref{lem:resetnormal91} gives
the full branch normal form, and Proposition~\ref{prop:causalnormal91}
gives the physical causal normalization independently of this example.

\begin{theorem}[Exact two-call hierarchy]\label{thm:operationalhierarchy90}
The alternative basis index in this experiment is drawn once and
reused at both calls.  Independently redrawing it instead gives the
scalar product channel and no improvement over the larger prior.
For every finite family satisfying \eqref{eq:basismoments90}, every
$d\ge2$, and every $0\le t\le1$, the game \eqref{eq:gamefamily90}--\eqref{eq:gameprior90}
has the exact optimal values
\begin{equation}\label{eq:operationalhierarchy90}
\begin{split}
 P_{\rm ca}&=\frac{d+1}{L},\\
 P_{\rm reset}&=\frac{d+1}{L}+\frac{t^2(d-1)}{2dL},\\
 P_{\rm all}&=\frac{d+1+t^2}{L}.
\end{split}
\end{equation}
For $t>0$ both inequalities are strict.  At $t=0$ all three
values are $1/2$.  The priors sum to one and
$p_E-p_C=t^2/L$, so $p_E$ is the constant-decision baseline.
The middle value is attained
by independent maximally entangled input--reference pairs and a final
joint reference readout, without feedback.  Its upper bound allows all
unit-reset feedback and unrestricted receiver dimension.  The last
value is attained by a parallel antisymmetric two-input state, and its
upper bound allows every adaptive two-call tester.
\end{theorem}
The theorem separates a finite ensemble decision task, not an arbitrary
pair of fixed channels and not a complete memory-dimension hierarchy.
The middle class has hard resources $N=Q=2$; the last construction is
one acquisition group of size two, with square cost four.

\subsection{A swap-filter identity}
For a Hermitian matrix $X$, write $X_+$ for its positive part.
\begin{lemma}[Negative mass of a filtered swap]\label{lem:swapfilter90}
Let $A,B\succeq0$ be $d\times d$ matrices and
$K=(\sqrt A\otimes\sqrt B)S(\sqrt A\otimes\sqrt B)$.  Then
\begin{equation}\label{eq:swapfilter90}
 \tr(-K)_+=\frac12\left[
  \bigl(\tr\sqrt{\sqrt A B\sqrt A}\bigr)^2-\tr(AB)\right]
 \le\frac{d-1}{2d}\tr A\tr B.
\end{equation}
The same negative mass is obtained from $(U\otimes V)S(U\otimes V)^*$
when $U^*U=A$ and $V^*V=B$, even if the receiving spaces are larger.
If $\tr A\tr B>0$, equality in the upper bound holds if and only if
$A=(\tr A)I/d$ and $B=(\tr B)I/d$.
\end{lemma}
\begin{proof}
First suppose $A,B$ are positive definite.  The matrix $K$ is similar
to $(A\otimes B)S$.  Conjugation of the latter by $I\otimes B$ gives
$(AB\otimes I)S$.  The matrix $AB$ is similar to $\sqrt A B\sqrt A$;
let its positive eigenvalues be $\lambda_1,\ldots,\lambda_d$.
Conjugating by the tensor square of an eigenbasis matrix leaves $S$
unchanged.  The resulting matrix has eigenvalues $\lambda_i$ on the
one-dimensional spaces indexed by $(i,i)$, and
$\pm\sqrt{\lambda_i\lambda_j}$ on the two-dimensional spaces indexed
by $(i,j),(j,i)$, $i<j$.  Thus its negative mass is
$\sum_{i<j}\sqrt{\lambda_i\lambda_j}$, which gives the equality.
By Cauchy--Schwarz,
$\sum_i\lambda_i\ge d^{-1}(\sum_i\sqrt{\lambda_i})^2$.
Moreover, the polar decomposition and Hilbert--Schmidt Cauchy--Schwarz
give $\|\sqrt A\sqrt B\|_1\le\sqrt{\tr A\tr B}$.
These inequalities give the bound.  Approximation by positive definite
matrices proves the singular case.  The polar decompositions of $U,V$
identify their filtered swap with $K$ on its support, adding only zero
eigenvalues in the orthogonal complement.
For the equality assertion, put $f=\|\sqrt A\sqrt B\|_1$ and
$a=\tr A$, $b=\tr B$.  Equality in the stated bound, with $ab>0$,
forces equality in both $\tr(AB)\ge f^2/d$ and $f^2\le ab$.
The variational formula for the trace norm and equality in
Hilbert--Schmidt Cauchy--Schwarz imply $B=(b/a)A$: indeed a maximizing
unitary $U$ has $\sqrt B U=c\sqrt A$, and multiplication by its
adjoint gives this proportionality.  Equality in the first inequality
makes all eigenvalues of $\sqrt A B\sqrt A$ equal and positive.
Consequently $A=(a/d)I$ and $B=(b/d)I$.  These scalar matrices plainly
attain the bound.  This also excludes equality at nonzero singular
factors.
\end{proof}
The normalized equality statement and its singular case are expanded
in the proof of Lemma~\ref{lem:centeredswap91}: normalized fidelity one
forces equality of the two density matrices, and eigenvalue
Cauchy--Schwarz then forces all $d$ eigenvalues to be $1/d$.

\subsection{Payoff and complete normalization}
We use the fully transposed, unnormalized Choi representative
\begin{equation}\label{eq:choiconvention90}
 J^\sharp(E)=\left[(\operatorname{Id}\otimes\mathcal M_E)
          (|\Omega\rangle\langle\Omega|)\right]^{\mathsf T}
       =\sum_y E_y\otimes|y\rangle\langle y|,
 \qquad |\Omega\rangle=\sum_i|ii\rangle.
\end{equation}
Its output trace is $I$.  A simultaneous full transpose of the usual
tester representative makes its contraction with $J^\sharp$ the Born
probability.  Factor order is $I_1,Y_1,I_2,Y_2$; expressions grouping
$I_1,I_2$ use the canonical permutation of these factors.
Dephasing the output factors does not change any score.  Write $T_{0,yz}$
for the input block of the tester decision $C$.  The gain over always
guessing the alternative is
\begin{equation}\label{eq:payoff90}
 P-p_E=\sum_{y,z}\tr(T_{0,yz}W_{yz}),\qquad
 W_{yz}=p_C C_y\otimes C_z-p_E\sum_b w_b E_y^b(t)\otimes E_z^b(t).
\end{equation}
The averaged tensor in the second term is $(1-t^2)I/d^2$ plus $t^2$
times \eqref{eq:basismoments90}.  With $k=t^2/(dL)$, substitution gives
\begin{equation}\label{eq:payoffblocks90}
 W_{yy}=-kS,\qquad W_{yz}=-\frac{2k}{d-1}\Pi_-\quad(y\ne z).
\end{equation}
These operators are real in the coordinate basis; a full transpose
therefore does not alter the following congruences.

\begin{proof}[Proof of Theorem~\ref{thm:operationalhierarchy90}]
\emph{Classical adaptation.}
Each block of a measure-and-reprepare tester is a sum of positive
products $A\otimes B$.  Since
$\tr[(A\otimes B)S]=\tr(AB)\ge0$, equal labels contribute no positive
gain in \eqref{eq:payoff90}; different labels cannot contribute positively
either.  Hence $P_{\rm ca}\le p_E$, with equality by the constant decision.
This proof also covers limits of such testers.

\emph{All unit-reset protocols.}
Purify the first acquisition and write its state as
$\sum_i|i\rangle C_0|i\rangle$, with $\tr C_0^*C_0=1$.
After the first device label $y$, combine all intervening receiver
operations into an instrument with recorded outcome $h$.  For an upper
bound, retain its Stinespring environment in the receiver.  The first
conditional reference operator then has the form
$C_{yh}E_y^{\mathsf T}C_{yh}^*$, where
\[
 A_{yh}=C_{yh}^*C_{yh},\qquad
 \sum_h A_{yh}=C_0^*C_0\quad\hbox{for each }y.
\]
Retaining this environment only enlarges receiver processing and does
not import it into the second acquisition.  Conditional reset permits
a purification of that fresh acquisition with a map $D_{yh}$ satisfying
$B_{yh}=D_{yh}^*D_{yh}$ and $\tr B_{yh}=1$.
Its complete system is independent of the old receiver conditional on
$(y,h)$.  The two-reference payoff operator at record $(y,h,z)$ is
therefore
\[
 (C_{yh}\otimes D_{yh})W_{yz}^{\mathsf T}
                         (C_{yh}\otimes D_{yh})^*.
\]
For different labels it is nonpositive.  For equal labels, maximizing
any final effect bounded by the identity gives at most its positive
mass, which by Lemma~\ref{lem:swapfilter90} is
$k(d-1)\tr A_{yh}/(2d)$.  Summing over $y,h$ gives
\[
 P_{\rm reset}-p_E\le k(d-1)/2=\frac{t^2(d-1)}{2dL}.
\]
Public randomization can be included in the recorded history.  A normal
instrument with countably many outcomes is treated by trace-class
sums, or by finite truncations retaining a remainder outcome and then
taking the trace-norm limit.  The bound is independent of receiver
size.  Thus no finite-dimensional restriction on receiver storage has
entered the upper bound.
For $t>0$, the equality case of Lemma~\ref{lem:swapfilter90} also shows
that a nonrandomized protocol attaining this bound in the displayed
purified normal form must satisfy
$A_{yh}=(\tr A_{yh})I/d$ and $B_{yh}=I/d$ on every nonzero branch.
All deficits in the summed upper bound are nonnegative, so equality
forces their branchwise vanishing.  Summing over $h$ then gives
$C_0^*C_0=I/d$.  Maximal mixing is thus a consequence of equality,
not a restriction imposed on the competing protocols.  This statement
concerns the purified normal form, not uniqueness of its dilations
or of the final readout.

For attainment, prepare $|\Omega\rangle/\sqrt d$ independently on each
input and reference.  Declare $C$ precisely when the two device labels
agree and the references lie in $\Pi_-$.  The event probability under
$C$ is $a=(d-1)/(2d^2)$.  Under any $E^b(t)$ it is $(1-t^2)a$:
indeed $\tr E_y^b(t)=1$, $\tr E_y^b(t)^2=t^2+(1-t^2)/d$, and
$\tr[\Pi_-(E\otimes E)]=[(\tr E)^2-\tr E^2]/2$.
The factor $d^{-2}$ comes from the two normalized input--reference
pairs.  The resulting gain is
$a[p_C-p_E(1-t^2)]=t^2(d-1)/(2dL)$.

\emph{All adaptive protocols.}
The deterministic normalization of any two-call tester is
$T_0+T_1=\Xi\otimes I_{Y_2}$, where, after dephasing $Y_1$,
$\Xi$ has input blocks $\Xi_y\succeq0$ satisfying
\[
 \tr_{I_2}\Xi_y=\rho\quad(1\le y\le d),\qquad \tr\rho=1.
\]
Here is a direct realization of these constraints in the present
classical-output interface.  Purify the initial state as
$|\psi\rangle\in I_1\otimes R$, and let $\mathcal A_y$ be the
trace-preserving intervening map from $R$ to $I_2\otimes R'$ after
observing $y$.  Put
\[
 X_y=(\operatorname{Id}_{I_1}\otimes\mathcal A_y)
                 (|\psi\rangle\langle\psi|),\qquad
 \Xi_y=\tr_{R'}X_y,\qquad
 \rho=\tr_R|\psi\rangle\langle\psi|.
\]
Positivity and trace preservation give the displayed marginal equations.
If $F_{yz}$ is the final receiver effect for decision $C$, then
\[
 T_{0,yz}=\tr_{R'}\bigl[(I_{I_1I_2}\otimes\sqrt{F_{yz}})
           X_y(I_{I_1I_2}\otimes\sqrt{F_{yz}})\bigr].
\]
The Choi contraction \eqref{eq:choiconvention90} gives its Born
probability, and $0\preceq F_{yz}\preceq I$ gives
$0\preceq T_{0,yz}\preceq\Xi_y$.  This also proves positivity of the
complement, without an input-dimension factor from a normalized Choi
state.  In particular $\tr\Xi_y=1$.
Since $-kS\preceq kI$ and the other payoff blocks are nonpositive,
\[
 P_{\rm all}-p_E\le k\sum_y\tr\Xi_y=kd=t^2/L.
\]
Prepare any normalized antisymmetric pure state on the two inputs and
declare $C$ when their labels agree.  The event probability is $1/d$
under $C$, and $(1-t^2)/d$ under every $E^b(t)$, because
$(P_y^b\otimes P_y^b)\Pi_-=0$.  The gain is $t^2/L$.
This is a legitimate parallel two-call strategy and attains the
adaptive upper.  The case $t=0$ gives equal values $1/2$ directly.
\end{proof}

\begin{corollary}[Seven qubit devices]\label{cor:qubitgame90}
Take the scalar qubit measurement with prior $2/5$ and each of the six
ordered Pauli measurements with prior $1/10$.  With the device fixed
for both calls and the binary decision specified above, the exact
success probabilities are
\[
 P_{\rm ca}=3/5,\qquad P_{\rm reset}=13/20,\qquad P_{\rm all}=4/5.
\]
\end{corollary}
\begin{proof}
Use Lemma~\ref{lem:finitedesign90} and substitute $d=2,t=1$.
\end{proof}
The antisymmetric readout and comparison mechanism have antecedents in
measurement comparison \cite{ZHS90}; reductions between measurement and
state discrimination also have a substantial history \cite{SZ74}.
Here the assertion is the joint exact optimization over the three
specified classes, including unrestricted reset feedback, on one finite
classical-output ensemble and its full-rank deformation.

\subsection{Quantitative stability of the operational gap}
\begin{corollary}[A stable score separation]\label{cor:robustgame90}
Suppose every device in the finite game is replaced by a classical-output
measurement channel at diamond distance at most $\delta$ from its
specified channel.  Keep the priors and the two-call resource classes
fixed.  Each of their optimal success probabilities changes by at most
$\delta$.  If the displayed reset implementation also has total hard
pathwise preparation defect at most $\varepsilon$ in the sense of
Proposition~\ref{prop:approxreset89}, its score exceeds the optimal
perturbed classical-adaptation score whenever
\begin{equation}\label{eq:robustgap90}
 2\delta+\varepsilon/2<\frac{t^2(d-1)}{2dL}.
\end{equation}
\end{corollary}
\begin{proof}
A two-call hybrid changes each normalized output in unhalved trace norm
by at most $2\delta$.  A fixed decision event changes by at most
$\delta$, and averaging over the unchanged priors preserves this bound.
Taking suprema over a fixed strategy class gives the same two-sided
bound for its optimum.  The preparation certificate changes each
implemented output by at most $\varepsilon$, hence its decision
probability by at most $\varepsilon/2$.  Subtract the classical upper
from the reset lower in \eqref{eq:operationalhierarchy90}.
\end{proof}
These are conditional norm estimates, not physical calibration claims.
In the preceding corollary the same $\delta$ bounds the scalar channel
and each entire basis measurement channel individually.  Priors and the
once-drawn-index law stay fixed; index-dependent channel perturbations
are permitted.  Corollary~\ref{cor:gamestability91} additionally treats
prior and law perturbations, with the unhalved norm factors explicit.
They make the relative scale needed to preserve a small score gap explicit.

\subsection{Hard budget conventions beyond the normalized domain}
\begin{proposition}[Normalization of arbitrary hard thresholds]\label{prop:budgetdomain90}
For positive integer reservations, $N_\pi\le Q_\pi\le N_\pi^2$.
Given real hard bounds $n,q\ge1$, put
\[
 M=\min\{\lfloor n\rfloor,\lfloor q\rfloor\},\qquad
 R=\min\{\lfloor q\rfloor,M^2\}.
\]
The class $\{\pi:N_\pi\le n,\ Q_\pi\le q\}$ equals
$\{\pi:N_\pi\le M,\ Q_\pi\le R\}$, and $M\le R\le M^2$.
Consequently Theorem~\ref{thm:quadraticbudget89} applies with $(M,R)$.
\end{proposition}
\begin{proof}
On each path $\sum n(h)\le\sum n(h)^2\le(\sum n(h))^2$.
Taking hard suprema proves the first assertion.  Bounded integer costs
have integer suprema, so both constraints can be floored.  The first
inequality reduces the call bound to $M$ and the last reduces the
quadratic bound to $R$.  Conversely the reduced bounds imply the original
ones.  Finally $M\le\lfloor q\rfloor$ and $M\le M^2$.
\end{proof}
The normalization $N\le Q$ allows the $N$-singleton lower construction;
it does not say that every admissible individual policy has $Q_\pi\ge N$.
If a hard threshold is below one, only the zero-call class is available.
Fractional-duration reservations and expected-cost constraints are
different models.  None is substituted for a hard integer reservation here.

\section{An exact hierarchy with equal hypothesis priors}
\label{sec:equalprior91}
The strict separation does not require the prior in
\eqref{eq:gameprior90}.  Equal hypothesis priors lead to a different
optimal receiver readout and to another exact hierarchy.  The finite
weighted second-moment family is unchanged.

\subsection{A trace-norm bound}
Let $S$ be the swap on $\C^d\otimes\C^d$ and
$\Pi_\pm=(I\pm S)/2$.
\begin{lemma}[The centered filtered swap]\label{lem:centeredswap91}
For $d\ge2$, $A,B\succeq0$, and $L=\sqrt A\otimes\sqrt B$,
\begin{equation}\label{eq:centeredswap91}
 \|L(I-dS)L\|_1\le\left(d-\frac1d\right)\tr A\tr B.
\end{equation}
For nonzero factors equality holds precisely when $A$ and $B$ are
positive scalar multiples of $I$.  The same norm and bound hold for
rectangular factors having Gram matrices $A,B$.
\end{lemma}
\begin{proof}
Put $K=LSL$ and $f=\|\sqrt A\sqrt B\|_1$.
The spectrum calculation in Lemma~\ref{lem:swapfilter90} gives
$\|K\|_1=f^2$, including singular factors by continuity.
Since $I-dS=2\Pi_--(d-1)S$, the triangle inequality gives
\begin{align*}
 \|L(I-dS)L\|_1
 &\le 2\tr(L\Pi_-L)+(d-1)\|K\|_1\\
 &=\tr A\tr B-\tr(AB)+(d-1)f^2\\
 &\le\tr A\tr B+(d-1-d^{-1})f^2\\
 &\le(d-d^{-1})\tr A\tr B.
\end{align*}
The penultimate step is eigenvalue Cauchy--Schwarz,
$\tr(AB)\ge f^2/d$; the last is $f^2\le\tr A\tr B$.

Here is the equality argument, also supplying a detailed treatment of
singular nonzero factors in Lemma~\ref{lem:swapfilter90}.
Normalize $a=A/\tr A$, $b=B/\tr B$.  The variational formula for the
trace norm selects a unitary $U$ with
$\operatorname{Re}\tr(\sqrt a\sqrt b U)=\|\sqrt a\sqrt b\|_1$.
Both $\sqrt a$ and $\sqrt b U$ have Hilbert--Schmidt norm one.
Equality in Cauchy--Schwarz, with the phase chosen as above, forces
$\sqrt b U=\sqrt a$; multiplication by the adjoint gives $b=a$.
This implication uses no invertibility.  Equality in
$\tr(AB)\ge f^2/d$ then gives $\tr(a^2)=1/d$.  The $d$
nonnegative eigenvalues of $a$, including zeros, sum to one, so equality
in their Cauchy--Schwarz inequality forces all of them to equal $1/d$.
Thus nonzero singular factors cannot attain equality.  Since
$d-1-d^{-1}>0$, equality in \eqref{eq:centeredswap91} forces both
of these equalities.  Scalar factors attain the bound directly from
the eigenvalues $1-d$ on $\Pi_+$ and $1+d$ on $\Pi_-$.
Polar decompositions of rectangular factors preserve the nonzero
spectrum and add only zero eigenvalues, proving the last assertion.
\end{proof}

\subsection{Equal-prior optimization}
\begin{theorem}[Equal-prior operational hierarchy]\label{thm:equalprior91}
Let $(w_b,P_y^b)$ be any weighted exact second-moment family satisfying
\eqref{eq:basismoments90}.  Choose the scalar measurement $C_y=I/d$
with probability $1/2$, or choose the alternative family with probability
$1/2$ and then draw one basis index $b$ with law $w_b$.  In the latter
case the same device $E_y^b(t)=(1-t)I/d+tP_y^b$ is used at both calls.
For $d\ge2$ and $0\le t\le1$, the three classes of
Definition~\ref{def:classes91} have exact optimal success probabilities
\begin{equation}\label{eq:equalprior91}
 \begin{split}
 P_{\rm ca}^{\rm eq}&=\frac12+\frac{t^2(d-1)}{2d(d+1)},\\
 P_{\rm reset}^{\rm eq}&=\frac12+\frac{t^2(d-1)}{2d^2},\\
 P_{\rm all}^{\rm eq}&=\frac12+\frac{t^2}{2d}.
 \end{split}
\end{equation}
Both inequalities are strict for $t>0$, including the full-rank interval
$0<t<1$; all three values equal $1/2$ at $t=0$.
The reset optimum is attained with independent maximally entangled
input--reference pairs and no receiver feedback.  Its upper includes
all unit-reset receiver instruments and arbitrary receiver dimension.
The unrestricted optimum is attained by a normalized antisymmetric
two-input state and bounds every adaptive strategy.
If the alternative basis is instead redrawn independently at each
call, its two-slot averaged channel is the scalar product channel and
all three equal-prior values are $1/2$.
\end{theorem}
\begin{proof}
Put $a=t^2/[2d^2(d+1)]$.  Subtracting the alternative-weighted
second moment from the scalar-weighted product gives
\begin{equation}\label{eq:equalpayoff91}
 W_{yy}=a(I-dS),\qquad
 W_{yz}=-\frac{a}{d-1}(I-dS)\quad(y\ne z).
\end{equation}
Indeed the diagonal moment is
$(1-t^2)I/d^2+t^2(I+S)/[d(d+1)]$, and the off-diagonal
moment is $(1-t^2)I/d^2+t^2(dI-S)/[d(d^2-1)]$.
These substitutions prove \eqref{eq:equalpayoff91}.
Both $S$ and $\Pi_\pm$ are real in the chosen basis; their full
transpose does not alter these payoff blocks.  The gain over the
constant alternative decision is $\sum_{yz}\tr(T_{C,yz}W_{yz})$.

\emph{Classically adaptive upper.}
After a complete first readout, a branch has an input effect
$A\succeq0$, with $\sum_h A_{yh}=\rho_1$ for each $y$ and
$\tr\rho_1=1$.  The second single-call tester has normalization
$B\succeq0$, $\tr B=1$, and decision blocks $0\preceq T_z\preceq B$.
For $\tr A>0$ set $\widehat A=A/\tr A$.
Contracting \eqref{eq:equalpayoff91} with $A$ in the first slot gives
$a\tr A(I-d\widehat A)$ for $z=y$ and its negative divided by
$d-1$ for $z\ne y$.  For a Hermitian $Q$ and $0\preceq T\preceq B$,
$\tr(TQ)\le\tr(BQ_+)$.  Therefore the branch gain is at most
\[
 a\tr A\,\tr[B|I-d\widehat A|]\le a(d-1)\tr A.
\]
The final inequality follows from $0\preceq\widehat A\preceq I$ and
$d\ge2$.  The zero branch contributes zero.
Summing $\sum_{yh}\tr A_{yh}=d$ proves the first upper bound.
It covers arbitrary second-call references and decisions, not only
unassisted classical probes.  Norm closure preserves it.

To attain it, use the same fixed pure input $|u\rangle$ twice and
declare $C$ on unequal labels.  Under $C$ the equal-label probability is
$1/d$; under the averaged alternative it is
\[
 \sum_{by}w_b\langle u,E_y^b(t)u\rangle^2
 =\frac1d+\frac{t^2(d-1)}{d(d+1)}.
\]
This follows directly from the second moments.  The corresponding
Bayes gain is half their difference, namely $ad(d-1)$.

\emph{Unit-reset upper.}
Use Lemma~\ref{lem:resetnormal91} without any restriction on its
rectangular receiver spaces.  On a branch $(y,h)$ put
\[
 X=(C_{yh}\otimes D_{yh})(I-dS)(C_{yh}\otimes D_{yh})^*.
\]
The elementary identity
$\max_{0\preceq F\preceq I}\tr(FX)=\tr(X_+)$ and
\eqref{eq:equalpayoff91} give total branch gain at most
\[
 a\tr(X_+)+a\tr((-X)_+)=a\|X\|_1
 \le a(d-d^{-1})\tr A_{yh},
\]
since $\tr b_{yh}=1$.  There are $d-1$ opposite-sign output blocks,
which cancel their denominator.  Summing over $y,h$ gives
$a(d^2-1)=t^2(d-1)/(2d^2)$.
Countable outcomes and null histories use unnormalized positive tails.

For attainment, prepare two independent normalized maximally entangled
pairs $|\Omega\rangle/\sqrt d$.  On equal device labels declare $C$
when the references lie in $\Pi_-$; on unequal labels declare $C$
when the references lie in $\Pi_+$.  This is one final joint receiver
measurement, with complementary effect $I-F_C$ on each classical
block.  Its event probabilities are
\begin{equation}\label{eq:equalresetevents91}
 q_C=\frac{(d-1)(d+2)}{2d^2},\qquad
 q_b=\frac{(d-1)(d+2-2t^2)}{2d^2}
 \quad\hbox{for every }b.
\end{equation}
To verify them, the receiver block is
$(E_y^b(t))^{\mathsf T}\otimes(E_z^b(t))^{\mathsf T}/d^2$ and
$\tr[\Pi_\pm(A\otimes B)]=[\tr A\tr B\pm\tr(AB)]/2$.
Here $\tr E_y=1$,
$\tr E_y^2=t^2+(1-t^2)/d$, and
$\tr E_yE_z=(1-t^2)/d$ for $y\ne z$.
Counting the $d$ equal-label and $d(d-1)$ unequal-label blocks gives
\eqref{eq:equalresetevents91}.  Half the difference is the reset upper.

\emph{Unrestricted upper.}
By Proposition~\ref{prop:causalnormal91}, $T_{C,yz}\preceq\Xi_y$
and $\tr\Xi_y=1$.  The sum of the positive payoff parts for a fixed
first label is
\[
 (W_{yy})_++\sum_{z\ne y}(W_{yz})_+
 =a(d+1)\Pi_-+a(d-1)\Pi_+\preceq a(d+1)I.
\]
Thus the gain is at most $ad(d+1)=t^2/(2d)$.
Any normalized state supported on the nonzero antisymmetric input
subspace attains it by declaring $C$ on equal labels: their probability
is $1/d$ under $C$ and $(1-t^2)/d$ under every alternative.

The gaps are
\begin{equation}\label{eq:equalgaps91}
 P_{\rm reset}^{\rm eq}-P_{\rm ca}^{\rm eq}
 =\frac{t^2(d-1)}{2d^2(d+1)},\qquad
 P_{\rm all}^{\rm eq}-P_{\rm reset}^{\rm eq}
 =\frac{t^2}{2d^2}.
\end{equation}
Finally $\sum_bw_b\mathcal M_{E^b(t)}=\mathcal M_C$ as a channel
identity.  Two independent draws give its tensor square, hence the
same two-slot statistics even for adaptive testers.  This proves the
last assertion and the endpoint statements.
\end{proof}
The equal priors concern the two hypotheses, not a uniform prior on all
individual devices.  The result removes prior tuning, but still uses
a shared latent device.  It does not assert a gap for every fixed pair.
Together with Theorem~\ref{thm:resetvariational91}, it gives a general
optimization principle and an explicitly solved family, rather than
identifying these two levels of generality.

\subsection{Seven explicit qubit devices}
\label{sec:qubitexplicit91}
For $d=2$ write
\[
 X=\begin{pmatrix}0&1\\1&0\end{pmatrix},\quad
 Y=\begin{pmatrix}0&-i\\i&0\end{pmatrix},\quad
 Z=\begin{pmatrix}1&0\\0&-1\end{pmatrix}.
\]
The devices are $C=(I/2,I/2)$ and, for $A\in\{X,Y,Z\}$ and
$\epsilon\in\{-1,1\}$,
\[
 E^{A,\epsilon}(t)=
 \left(\frac{I+t\epsilon A}{2},\frac{I-t\epsilon A}{2}\right).
\]
The six alternatives have conditional weight $1/6$.
Both outcome orders are included: they cancel the first Pauli moments
and supply the required ordered second moments.  At $t=1$ they are the
six ordered Pauli eigenbasis measurements.  For the prior of
Theorem~\ref{thm:operationalhierarchy90}, $C$ has mass $2/5$ and each
alternative mass $1/10$; for Theorem~\ref{thm:equalprior91}, these masses
are $1/2$ and $1/12$.

For reset acquisition prepare
\[
 |\phi^+\rangle_{I_1R_1}\otimes|\phi^+\rangle_{I_2R_2},\qquad
 |\phi^+\rangle=(|00\rangle+|11\rangle)/\sqrt2.
\]
The inputs are independent.  The two retained references are measured
jointly only after both calls.  Let
$|\psi^-\rangle=(|01\rangle-|10\rangle)/\sqrt2$,
$\Pi_-=|\psi^-\rangle\langle\psi^-|$ and $\Pi_+=I_4-\Pi_-$
on $R_1\otimes R_2$.
The receiver block for any two devices is $E_y^{\mathsf T}\otimes
F_z^{\mathsf T}/4$.  Under $J^\sharp$, the decision tester block is
$F_{C,yz}^{\mathsf T}/4$, where
\[
 F_{C,yz}=\mathbf1_{y=z}\Pi_-
 \quad\hbox{for the prior }2/5,
 \qquad
 F_{C,yz}=\begin{cases}\Pi_-,&y=z,\\\Pi_+,&y\ne z\end{cases}
 \quad\hbox{for equal priors}.
\]
The other block is $(I_4-F_{C,yz})^{\mathsf T}/4\succeq0$.
Their sum is $I_4/4$ for each $(y,z)$, with partial trace
$I_2/2$ over $I_2$ and trace one.  This verifies the causal normalization
in the order $I_1,Y_1,I_2,Y_2$, using the canonical permutation to group
the input copies.  The transpose is required for arbitrary complex
receiver effects, even though the displayed antisymmetric projectors
are real.

At $t=1$, the first reset event has probabilities $1/8$ under $C$ and
zero under every alternative.  Its Bayes score is
$(2/5)(1/8)+(3/5)=13/20$.
The equal-prior reset event has probabilities $1/2$ and $1/4$,
respectively, giving $\frac12(\frac12)+\frac12(\frac34)=5/8$.
For unrestricted acquisition put the normalized state $|\psi^-\rangle$
on $I_1\otimes I_2$, not on the references, and declare $C$ on equal
labels.  The probabilities are $1/2$ and zero, yielding $4/5$ for the
first prior and $3/4$ for equal priors.  For classical adaptation,
the biased-prior optimum is the constant decision $3/5$.
For equal priors, two inputs $|0\rangle$ and the unequal-label event
have probabilities $1/2$ under $C$ and $1/3$ averaged over alternatives,
giving $7/12$.  Consequently the two exact triples on the same devices are
\[
 \frac35<\frac{13}{20}<\frac45,
 \qquad
 \frac7{12}<\frac58<\frac34.
\]
The class upper bounds, not these evaluations alone, prove optimality.

\subsection{An open neighborhood of the equal-prior experiment}
\begin{corollary}[Channel, prior, and ensemble stability]\label{cor:gamestability91}
Fix the family and $t>0$.  Replace the scalar device and each entire
basis measurement channel individually by channels within unhalved
diamond distance $\delta$.  The perturbations may depend on the basis
index and may change all its effects, but must form legal channels on
the same interfaces.  Keep the once-drawn-index rule.  Replace the
scalar prior $1/2$ by $q$ and the conditional basis law $w$ by $\nu$,
and put
\[
 \eta=\tfrac12\sum_b|\nu_b-w_b|,\qquad
 \kappa=\delta+|q-\tfrac12|+\eta.
\]
For each of the three fixed operational classes,
$|\widetilde P_{\rm class}-P_{\rm class}^{\rm eq}|\le\kappa$.
Thus both class gaps remain strict when
$2\kappa<t^2(d-1)/[2d^2(d+1)]$.
If the exhibited reset protocol has, in addition, a uniform
reference-complete diamond certificate of preparation defect whose
hard pathwise sum is at most $\varepsilon$, its implemented score
still exceeds the perturbed classical optimum provided
\[
 2\kappa+\varepsilon/2<\frac{t^2(d-1)}{2d^2(d+1)}.
\]
\end{corollary}
\begin{proof}
For a fixed protocol and device index, the two-call hybrid bounds
output trace distance by $2\delta$.  Normalization makes the change
of any decision probability at most $\delta$, with the factor $1/2$
for unhalved trace distance.  Changing the binary prior changes success
by at most $|q-1/2|$; changing the alternative law changes an average
of $[0,1]$-valued successes by at most $\eta$.  The bound is uniform
in the protocol, so taking suprema preserves it in both directions.
Subtract and use \eqref{eq:equalgaps91}.
The preparation hybrid changes each implemented output by at most
$\varepsilon$, hence its success by at most $\varepsilon/2$.
This last statement concerns a certified implementation of the displayed
reset protocol, not an enlargement of the ideal optimization class.
\end{proof}

\subsection{Quantitative rigidity of near-optimal reset strategies}
The equality condition for the centered swap has a stable form.  It
controls the input Gram matrices in a chosen purified normal form,
rather than specifying unique instruments, dilations, or receiver
readouts.  For $d\ge2$ set
\[
 c_d=d-1-d^{-1}>0,\qquad
 K_d=d\left(1+\frac{9}{2c_d}\right).
\]
\begin{lemma}[Stable centered-swap equality]\label{lem:stableswap91}
For density matrices $a,b$ on $\C^d$, let
\[
 \Delta(a,b)=d-d^{-1}
  -\| (\sqrt a\otimes\sqrt b)(I-dS)
                         (\sqrt a\otimes\sqrt b)\|_1.
\]
Then $\Delta(a,b)\ge0$ and
\begin{equation}\label{eq:stableswap91}
 \|a-I/d\|_1^2\le K_d\Delta(a,b),\qquad
 \|b-I/d\|_1^2\le K_d\Delta(a,b).
\end{equation}
No invertibility is required.
\end{lemma}
\begin{proof}
Put $f=\|\sqrt a\sqrt b\|_1$, $\eta=1-f^2$, and
$e=\tr(ab)-f^2/d$.  The proof of Lemma~\ref{lem:centeredswap91}
gives $\eta,e\ge0$ and
\begin{equation}\label{eq:swapdefectdecomposition91}
 \Delta(a,b)\ge c_d\eta+e.
\end{equation}
Choose a unitary $V$ extending the right polar factor so that
$Z=\sqrt a\sqrt b V=\sqrt{\sqrt a b\sqrt a}\succeq0$.
This is possible also when the matrices are singular.  The
Hilbert--Schmidt norm, denoted by $\|\cdot\|_2$, satisfies
\[
 \|\sqrt b V-\sqrt a\|_2^2=2(1-f),\quad
 \|Z-a\|_2\le\sqrt{2\eta},\quad
 \|Z-fI/d\|_2^2=e.
\]
Here $\|\sqrt a\|_{\rm op}\le1$ and $1-f\le\eta$.
Since also $1-f\le\sqrt\eta$, the triangle inequality yields
\[
 \|a-I/d\|_2\le(\sqrt2+d^{-1/2})\sqrt\eta+\sqrt e.
\]
Weighted Cauchy--Schwarz bounds its square by
$[1+(\sqrt2+d^{-1/2})^2/c_d](c_d\eta+e)$.
For $d\ge2$, $(\sqrt2+d^{-1/2})^2\le9/2$; and
$\|T\|_1^2\le d\|T\|_2^2$ for a $d\times d$ matrix $T$.
These facts and \eqref{eq:swapdefectdecomposition91} prove the first
estimate.  Interchanging $a,b$ proves the second, since swap
conjugation leaves the centered-swap norm unchanged.
\end{proof}

\begin{corollary}[Near-optimal acquisition marginals]\label{cor:resetstability91}
Fix the equal-prior experiment of Theorem~\ref{thm:equalprior91}
with $t>0$.  Suppose a unit-reset strategy achieves reward
$v\ge P_{\rm reset}^{\rm eq}-\varepsilon$.
Choose its purified normal form from Lemma~\ref{lem:resetnormal91}.
On nonzero branches put
\[
 a_{yh}=A_{yh}/\tr A_{yh},\qquad
 \mu_{yh}=\tr A_{yh}/d.
\]
Then $\sum_{yh}\mu_{yh}=1$ and
\begin{equation}\label{eq:resetstability91}
 \sum_{y,h}\mu_{yh}
 \left(\|a_{yh}-I/d\|_1^2+\|b_{yh}-I/d\|_1^2\right)
 \le \frac{4K_d d(d+1)}{t^2}\,\varepsilon.
\end{equation}
For a public mixture the same conclusion holds after adjoining the
seed and multiplying these weights by its law.  Finite and countable
recorded instruments are both allowed.
\end{corollary}
\begin{proof}
Write $\alpha=t^2/[2d^2(d+1)]$ for the coefficient in
\eqref{eq:equalpayoff91}.  The complete reset upper, before applying
the centered-swap bound, gives
\[
 v\le\tfrac12+\alpha\sum_{yh}\tr A_{yh}
          [d-d^{-1}-\Delta(a_{yh},b_{yh})].
\]
Purification does not obstruct recovering the original decision rule,
so this inequality bounds its achieved reward as well.  The common
Gram identity gives $\sum_{yh}\tr A_{yh}=d$.  Subtract from the
attained optimum to obtain
$\sum_{yh}\mu_{yh}\Delta(a_{yh},b_{yh})
\le\varepsilon/(\alpha d)$.
Apply Lemma~\ref{lem:stableswap91} to both factors and substitute
$\alpha$.  Positive trace tails justify countable sums.  For a public
mixture, apply the same upper to each component and average before
subtracting its reward from the common optimum.
\end{proof}
In particular, the total $\mu$-weight on which either normalized Gram
matrix is at trace distance at least $u>0$ from $I/d$ is at most
$4K_d d(d+1)\varepsilon/(t^2u^2)$.  These are Gram weights, not
probabilities under a selected hypothesis.  At zero loss the result
recovers branchwise isotropy.  The estimate is nontrivial at the
relative scale $\varepsilon/t^2$; it does not turn fixed calibration
error into vanishing-signal robustness.  It asserts neither uniqueness
of a dilation nor a device-independent certification of the reset cut.

\section{Rank-resolved receiver resources}
\label{sec:rankmemory92}
The distinction between classical and quantum receivers extends to an
attained hierarchy at each pair of retained-register dimensions.  This
section keeps the reset cut fixed and resolves its two sides separately.
It supplies a general variational characterization and an exactly solved
spectrum, rather than identifying all notions of bounded quantum memory.

\subsection{Two specified recording cuts}
Use the finite experiment of Definition~\ref{def:classes91}.
Fix integers $1\le k\le d_1$ and $1\le \ell\le d_2$.
\begin{definition}[Retained-register dimensions]\label{def:rankclass92}
The class $\mathfrak R_{k,\ell}$ consists of unit-reset strategies for
which, conditional on every recorded first history, the entire retained
old quantum system has dimension at most $k$ and the entire fresh
reference accompanying the second input has dimension at most $\ell$.
The old system and the complete fresh input--reference system are in a
tensor product at that cut.  All other quantum registers are discarded
before the second call; classical records may be retained without a
cardinality bound.  The initial first-call reference and temporary
workspaces before this cut have no dimension restriction.  The final
measurement on the two retained receivers is arbitrary.
Mixed preparations, recorded normal instruments, public randomization,
null histories, discarded terminal calls and finite-dimensional tester
closure have the conventions of Definition~\ref{def:classes91}.
Denote the optimal reward by $P_{k,\ell}$.
\end{definition}
In particular, this is not a bound on the largest workspace throughout
an implementation.  Enlarging the old output by keeping an instrument
environment would violate the stated resource.  The refinement in the
next lemma records that environment's Kraus index \emph{classically}
instead, without retaining it coherently.

\begin{lemma}[Dimension-preserving refinement]\label{lem:ranknormal92}
For the optimization over $\mathfrak R_{k,\ell}$, the normal form
\eqref{eq:resetgrams91}--\eqref{eq:resetblock91} is exact with
\begin{equation}\label{eq:rankconstraints92}
 \operatorname{rank}A_{yh}\le k,\qquad
 \operatorname{rank}b_{yh}\le\ell.
\end{equation}
Every implemented competitor can be refined into such branches while
preserving its score if the added classical indices are ignored in
its decisions.  Conversely every finite Gram collection satisfying
these ranks and the common-barycenter constraints is realizable in
$\mathfrak R_{k,\ell}$.
\end{lemma}
\begin{proof}
Purify the first input, retaining a reference whose relevant Schmidt
support has dimension at most $d_1$.  A completely positive branch of
the subsequent map from that support to $\C^k$ has a finite Kraus
representation $X\mapsto\sum_qV_{yhq}XV_{yhq}^*$, with
$V_{yhq}:\C^{r_0}\to\C^k$.  Refine $h$ to $(h,q)$, record $q$
classically, and use the original fresh preparation and final effects.
Summing over $q$ recovers the old channel and its decision statistics.
Thus $C_{yhq}=V_{yhq}C_0$ has rank at most $k$, and completeness
still gives $\sum_{h,q}C_{yhq}^*C_{yhq}=\rho$ for every $y$.
An original mixed first preparation is recovered by discarding its
purification before this old-output cut; this discarding is part of
that same completely positive branch, not an additional retained
quantum register.

A mixed fresh state on $I_2\otimes\C^\ell$ has a finite spectral
mixture of normalized pure states in that same space.  Sample and
record its mixture index, use the corresponding vector, and ignore
that index in the original decision.  Its coefficient map
$D:I_2\to\C^\ell$ has Gram rank at most $\ell$.  Multiplying the
old coefficient map by the square root of this mixing probability
preserves the common Gram sum.  This is a classical refinement, not a
purification into a larger fresh reference.  It is independent of the
old quantum output conditional on the refined record.  Countable
outcomes are handled by the trace tails of Lemma~\ref{lem:countable91}.
Public mixtures and tester limits preserve the resulting value bound.

Conversely choose $C_0^*C_0=\rho$ on its Schmidt support.  Factor each
$A_{yh}$ as $C_{yh}^*C_{yh}$ with
$C_{yh}:I_1\to\C^k$, and each $b_{yh}$ as $D_{yh}^*D_{yh}$ with
$D_{yh}:I_2\to\C^\ell$.  The maps $V_{yh}=C_{yh}C_0^+$ satisfy
$\sum_hV_{yh}^*V_{yh}=I$ on that support.  They give a recorded
trace-preserving instrument.  The normalized fresh vector $|D_{yh}\rangle$
is prepared independently, and polar isometries transfer the prescribed
leaf POVMs to $\C^k\otimes\C^\ell$.  Zero Gram branches are simply
zero maps.  The original unnormalized formulas define all null histories.
\end{proof}

\subsection{An attained hierarchy and its dual}
Write $\mathcal D_d^{\le q}$ for the compact set of density matrices of
rank at most $q$.  In the notation of \eqref{eq:leafvalue91}, put
\begin{align}
 g_y^\ell(a)&=\max_{b\in\mathcal D_{d_2}^{\le\ell}}\Phi_y(a,b),
 \label{eq:rankleaf92}\\
 c_y^{k,\ell}(\rho)&=\max\left\{\sum_h\lambda_hg_y^\ell(a_h):
 a_h\in\mathcal D_{d_1}^{\le k},\quad
 \lambda_h\ge0,\quad\sum_h\lambda_ha_h=\rho\right\}.
 \label{eq:rankroof92}
\end{align}
Weights sum to one by taking traces.  All density matrices have such a
representation, since rank-one atoms are allowed.

\begin{theorem}[Rank-resolved variational principle]\label{thm:rankroof92}
For every finite two-call experiment,
\begin{equation}\label{eq:rankprimaldual92}
\begin{split}
 P_{k,\ell}&=\max_{\rho\in\mathcal D_{d_1}}
                       \sum_yc_y^{k,\ell}(\rho)\\
 &=\min_{(H_y)}\lambda_{\max}\left(\sum_yH_y\right),\qquad
 \tr(H_ya)\ge g_y^\ell(a)\quad(a\in\mathcal D_{d_1}^{\le k}),
\end{split}
\end{equation}
where the $H_y$ are Hermitian.  Both extrema are attained.  At an
optimal barycenter $\rho$, at most $(\operatorname{rank}\rho)^2$
nonzero receiver outcomes after each first label suffice; a final
quantum receiver of dimension $k\ell$ suffices.
The values increase monotonically in either dimension, with
\[
 P_{1,d_2}=P_{\rm ca},\qquad P_{d_1,d_2}=P_{\rm reset}.
\]
The contact, spectral and quantitative defect criteria of
Theorem~\ref{thm:resetdual91} and Corollary~\ref{cor:contactdefect91}
hold with $g_y^\ell$ and the rank-restricted atoms.
\end{theorem}
\begin{proof}
The function $\Phi_y$ is continuous and homogeneous in its first
positive argument.  The rank-bounded sets are compact, so
$g_y^\ell$ is continuous and its graph on
$\mathcal D_{d_1}^{\le k}$ has compact convex hull.
The largest height over each $\rho$ is an attained, upper
semicontinuous concave envelope.  Lemma~\ref{lem:ranknormal92} gives
both operational directions.  The atom-elimination proof of
Theorem~\ref{thm:resetvariational91} uses only linear dependence of
atoms on the support of $\rho$, and does not change any atom.  It
therefore preserves the rank restriction and gives the same branch
bound.  Compactness supplies all fresh states and final POVMs.

For clarity, dual attainment does not require the rank-bounded atom
set to be convex.  The homogeneous hypograph of its concave envelope
is a closed convex cone, finite on every positive matrix and
nonnegative there.  At the common matrix $I/d_1$ with every height
$-1$, all these cones are strictly feasible.  The conic argument of
Theorem~\ref{thm:resetdual91} applies without alteration: its dual is
precisely majorization of the atom graph in \eqref{eq:rankprimaldual92}.
Bounded reward gives a finite value.  Slater duality yields no gap and
an attained Hermitian majorant.  Subtracting the primal value from the
dual upper gives the same three nonnegative defects, proving the
contact and quantitative assertions.  Inclusion of the rank-bounded
atom sets proves monotonicity.  The two endpoint identifications are
Theorems~\ref{thm:classicalroof91} and~\ref{thm:resetvariational91}.
\end{proof}
The dual inequalities are universal on their rank-bounded domains;
sampling is not a certificate.  In particular, strict improvement
between any two nested dimension pairs is equivalent to a feasible
finite strategy for the larger pair exceeding an attained dual upper
for the smaller pair.  This criterion covers arbitrary fixed pairs,
without asserting that every such pair exhibits a gap.

\subsection{Rank-sensitive swap bounds}
The rank in the next lemma is an upper bound on each of the two Gram
ranks; the active overlap may have a smaller rank.
\begin{lemma}[Swap bounds at finite rank]\label{lem:rankswap92}
Let $1\le r\le d$, $d\ge2$, and let $A,B\succeq0$ have ranks at most
$r$.  With $L=\sqrt A\otimes\sqrt B$,
\begin{align}
 \tr[-LSL]_+&\le\frac{r-1}{2r}\tr A\tr B,
 \label{eq:rankswapnegative92}\\
 \|L(I-dS)L\|_1&\le\left(d-\frac1r\right)\tr A\tr B.
 \label{eq:rankswapcenter92}
\end{align}
The bounds also hold if only the smaller of the two ranks is at most
$r$, and for rectangular factors with the same Grams.  Both bounds
are sharp: take $A=aP/r$, $B=bP/r$, where $P$ is any rank-$r$
projection and $a,b\ge0$.  For nonzero factors, equality in the first
bound with $r>1$ requires these forms with a common $P$; equality in
the second requires them when $(d,r)\ne(2,1)$.
\end{lemma}
\begin{proof}
The nonzero eigenvalues of $\sqrt A B\sqrt A$ number at most $r$.
Consequently, for $f=\|\sqrt A\sqrt B\|_1$,
\[
 \tr(AB)\ge f^2/r,\qquad f^2\le\tr A\tr B.
\]
Insert the first inequality into Lemma~\ref{lem:swapfilter90} to get
\eqref{eq:rankswapnegative92}.  The triangle step of
Lemma~\ref{lem:centeredswap91} gives
\[
 \|L(I-dS)L\|_1
 \le\tr A\tr B-\tr(AB)+(d-1)f^2
 \le\tr A\tr B+(d-1-r^{-1})f^2.
\]
The coefficient is nonnegative, proving \eqref{eq:rankswapcenter92}.
These arguments only use the smaller rank, not a common support
assumption.  The spectrum identities hold for singular factors by the
proved continuity argument, and polar isometries handle rectangular
maps.

For equality, when the coefficient of $f^2$ is positive, normalized
fidelity one gives $A/\tr A=B/\tr B$ without an invertibility
assumption, as in Lemma~\ref{lem:centeredswap91}.  Equality in the
rank-$r$ Cauchy--Schwarz inequality then forces exactly $r$ equal
nonzero eigenvalues, proving the common-projection forms.
For such forms the centered spectrum is $(1-d)ab/r^2$ with
multiplicity $r(r+1)/2$ and $(1+d)ab/r^2$ with multiplicity $r(r-1)/2$;
the negative swap has mass $ab(r-1)/(2r)$.  This proves sharpness.
The excluded centered case has zero coefficient and different equality
cases: for $d=2,r=1$, orthogonal pure factors also attain the bound.
For the first inequality at $r=1$, the negative mass is always zero.
\end{proof}

\subsection{The complete two-cut spectrum of the basis experiment}
The following formula resolves every pair of retained-register
thresholds on the same once-selected finite basis family.  No
additional design is introduced at a new threshold.
\begin{theorem}[Exact dimension spectra]\label{thm:rankspectrum92}
Fix $d\ge2$, $0\le t\le1$, and a weighted exact second-moment family
satisfying \eqref{eq:basismoments90}.  Under the once-drawn device rule,
let $P_{k,\ell}^{\rm b}$ use the priors of
Theorem~\ref{thm:operationalhierarchy90}, and let
$P_{k,\ell}^{\rm eq}$ use the equal hypothesis priors of
Theorem~\ref{thm:equalprior91}.  For $1\le k,\ell\le d$, set
$r=\min\{k,\ell\}$ and $L_0=2(d+1)-t^2$.  Then
\begin{align}
 P_{k,\ell}^{\rm b}
 &=\frac{d+1}{L_0}+\frac{t^2(r-1)}{2rL_0},
 \label{eq:rankbiased92}\\
 P_{k,\ell}^{\rm eq}
 &=\frac12+\frac{t^2(dr-1)}{2dr(d+1)}.
 \label{eq:rankequal92}
\end{align}
These are optima over the complete classes
$\mathfrak R_{k,\ell}$.  For $t>0$ they strictly increase at every
increase of the smaller threshold, and attain the unrestricted reset
values precisely at $k=\ell=d$.  The two-cut resource dependence is
therefore exactly through $\min\{k,\ell\}$ for these experiments.
\end{theorem}
\begin{proof}
Use the refined normal form of Lemma~\ref{lem:ranknormal92}; there is
no retained dilation environment outside the counted receivers.
Every branch has $\tr b_{yh}=1$ and smaller Gram rank at most $r$,
while $\sum_{yh}\tr A_{yh}=d$.
For the biased priors the equal-label payoff is
$-t^2 S/(dL_0)$ and all unequal-label payoffs are nonpositive.
Equation~\eqref{eq:rankswapnegative92} bounds the gain by
$t^2(r-1)/(2rL_0)$.
For equal priors write $\alpha=t^2/[2d^2(d+1)]$.
The $d-1$ opposite-sign payoff blocks cancel their denominator, so the
branch gain is at most
$\alpha\|(C_{yh}\otimes D_{yh})(I-dS)(C_{yh}\otimes D_{yh})^*\|_1$.
Equation~\eqref{eq:rankswapcenter92} and the Gram sum give
$\alpha d(d-1/r)$, which is \eqref{eq:rankequal92} minus $1/2$.
Classical mixtures and trace-norm closure preserve both upper bounds.

Here is a single finite construction attaining them.  Index the input
basis cyclically modulo $d$, and let $J_h:\C^r\to\C^d$ include the
$r$ consecutive coordinates starting at $h$, $0\le h<d$.
Put $P_h=J_hJ_h^*$; then $\sum_hP_h=rI$.
Choose the first coefficient map $C_0=I/\sqrt d$ and, after every
first device label, the receiver instrument and fresh coefficient maps
\begin{equation}\label{eq:cyclicrankprotocol92}
 V_h=J_h^*/\sqrt r,\qquad
 C_{yh}=J_h^*/\sqrt{dr},\qquad D_{yh}=J_h^*/\sqrt r.
\end{equation}
Thus $\sum_hV_h^*V_h=I$,
$A_{yh}=P_h/(dr)$, $b_{yh}=P_h/r$ and $\sum_hA_{yh}=I/d$.
The fresh vector has norm one and is independent of the old receiver.
Both retained quantum registers have dimension $r$; $h$ is classical.
The initial reference has dimension $d$, as explicitly allowed before
the recording cut in Definition~\ref{def:rankclass92}.

On these receivers the filtered swap is $S_r/(dr^2)$ and the centered
one is $(I_{r^2}-dS_r)/(dr^2)$.  For the biased prior, declare $C$
only on equal labels and the antisymmetric receiver projection.
For equal priors, declare $C$ on the antisymmetric projection on
equal labels and on the symmetric projection on unequal labels.
Complementary effects complete a POVM.  Their spectral signs attain
every branch bound above.  At $r=1$ the antisymmetric projection is
zero; the second readout reduces to the unequal-label rule.  At $r=d$
the instrument's redundant classical randomization may be omitted.
All constructions remain normalized at $t=0$ and on projective null
histories at $t=1$.
\end{proof}

\begin{corollary}[Adjacent gains and dimension certificates]
\label{cor:rankwitness92}
For $1\le r<d$, abbreviate the values in
Theorem~\ref{thm:rankspectrum92} by $P_r^{\rm b}$ and $P_r^{\rm eq}$.
Their adjacent gains are
\[
 P_{r+1}^{\rm b}-P_r^{\rm b}=\frac{t^2}{2L_0r(r+1)},\qquad
 P_{r+1}^{\rm eq}-P_r^{\rm eq}
       =\frac{t^2}{2d(d+1)r(r+1)}.
\]
For the equal-prior perturbations of Corollary~\ref{cor:gamestability91},
with its score error $\kappa$, every perturbed class value moves by at
most $\kappa$.  In particular an achieved ideal-reset score strictly
above $P_r^{\rm eq}+\kappa$ requires both counted receiver dimensions
to exceed $r$.  The adjacent gap survives if $2\kappa$ is smaller
than its displayed value.  A certified additional preparation defect
$\varepsilon$ changes the implemented score by at most $\varepsilon/2$
and may be added to these error budgets.
\end{corollary}
\begin{proof}
Subtract the consecutive exact formulas.  The hybrid and probability
law bounds are uniform in a fixed protocol, independently of its
receiver dimensions, so they hold for every rank-restricted supremum.
If either dimension were at most $r$, its nominal value would be at
most $P_r^{\rm eq}$ by Theorem~\ref{thm:rankspectrum92}.  This proves
the asserted certificate and the perturbation statements.
\end{proof}
These dimension certificates assume the specified reset architecture
and calibration bounds.  They are not device-independent conclusions
from an uncalibrated score.  The entire two-cut hierarchy still lies
inside the reset class; the unrestricted coherent-acquisition optimum
of Theorem~\ref{thm:equalprior91} is a different resource.

\subsection{Relation to constrained-memory discrimination}
Memory-constrained tester optimization and its relation to constrained
separability are established in \cite{OZNPQ89}, including the classical
adaptive boundary.  Theorem~\ref{thm:rankroof92} is a rank-envelope
characterization for the specified classical-output, fresh-preparation
cut, and Theorem~\ref{thm:rankspectrum92} solves its two-coordinate
values on one fixed finite family.  It does not replace the general
channel-memory formulations of that work.  The autonomous multi-time
hierarchy of \cite{ZB90} carries coherent information between successive
acquisitions and compiles arbitrary finite-time testers with a counter.
Here increasing the two receiver dimensions never removes the reset
cut.  Consequently saturation means the full reset optimum, not the
unrestricted strategy norm.  These different saturation statements
are compatible and cannot be interchanged.

\section{Initial spectra and rigidity of rank-optimal receivers}
\label{sec:spectralrigidity93}
The dimension spectrum determines an optimal value.  We now determine
which initial spectra can attain it, and how a small loss in value
restricts the acquisition marginals.  Throughout this section the basis
index is drawn once and retained for both calls, the hypothesis priors
are equal, and $0<t\le1$.  The family satisfies
\eqref{eq:basismoments90}.  Put
\[
 r=\min\{k,\ell\},\qquad
 p_r=\frac12+\frac{t^2(dr-1)}{2dr(d+1)},\qquad
 \alpha=\frac{t^2}{2d^2(d+1)}.
\]
All ranks and norms refer to the fixed input copies.  Trace norms are
unhalved.  An initial Gram matrix $\rho$ means a pure first acquisition
$|C_0\rangle$ with $C_0^*C_0=\rho$, in the convention of
Lemma~\ref{lem:resetnormal91}; its input marginal is $\rho^{\mathsf T}$.
The receiver cut is exactly Definition~\ref{def:rankclass92}.

\subsection{The complete set of optimal initial spectra}
\begin{theorem}[Optimal initial Gram matrices]\label{thm:initialspectra93}
A fixed initial Gram matrix $\rho\in\mathcal D_d$ admits a strategy
in $\mathfrak R_{k,\ell}$ attaining $p_r$ if and only if
\begin{equation}\label{eq:spectralcap93}
                  \rho\preceq I/r.
\end{equation}
Every such $\rho$ has a decomposition
\begin{equation}\label{eq:projectionframe93}
 \rho=\sum_{h=1}^{s}w_hP_h/r,\qquad s\le d,\quad
 w_h>0,\quad\sum_hw_h=1,
\end{equation}
where the $P_h$ are rank-$r$ projections commuting with $\rho$.
This decomposition gives an attaining strategy with the same receiver
instrument after every first device label and with both retained
quantum registers of dimension $r$.

Except when $(d,r)=(2,1)$, every attaining strategy in the
dimension-preserving refined normal form has, on each nonzero branch,
\begin{equation}\label{eq:optimalgrams93}
       A_{yh}=w_{yh}P_{yh}/r,\qquad b_{yh}=P_{yh}/r,
       \qquad \operatorname{rank}P_{yh}=r.
\end{equation}
Together with the common Gram sums and maximizing leaf decisions,
these conditions are also sufficient.  They specify Gram data, not
unique dilations or unique final measurements.
\end{theorem}
\begin{proof}
For a normalized branch write $a=A/\tr A$ and $b=b_{yh}$.
The equal-prior upper in Theorem~\ref{thm:rankspectrum92}, before
applying its final inequality, is
\begin{equation}\label{eq:unsaturatedrankupper93}
 v\le\frac12+\alpha\sum_{y,h}w_{yh}
       \| (\sqrt{a_{yh}}\otimes\sqrt{b_{yh}})(I-dS)
                    (\sqrt{a_{yh}}\otimes\sqrt{b_{yh}})\|_1.
\end{equation}
The smaller Gram rank is at most $r$, and $\sum_{y,h}w_{yh}=d$.
All branch deficits from the bound $d-1/r$ are nonnegative.
Attainment therefore forces equality on every branch of positive
weight.  When $(d,r)\ne(2,1)$, Lemma~\ref{lem:rankswap92} gives
\eqref{eq:optimalgrams93}.  For each fixed $y$, summing yields
$\rho=\sum_hw_{yh}P_{yh}/r\preceq I/r$, since
$\sum_hw_{yh}=1$.  In the exceptional case $r=1$, the claimed spectral
condition holds for every density matrix anyway.  The fixed-$\rho$
compact-envelope argument of Theorem~\ref{thm:rankroof92} gives
attainment of its constrained value, so the same conclusion applies
when an optimum was initially defined through tester closure.

We give a finite construction for the converse.  Diagonalize $\rho$
with eigenvalues $\lambda_i$ and set $x_i=r\lambda_i$ and
$s_i=\sum_{j\le i}x_j$, $s_0=0$.  Then $0\le x_i\le1$ and
$s_d=r$.  For almost every $u\in[0,1)$, mark the $r$ points
$u,u+1,\ldots,u+r-1$ in $[0,r)$.  Select coordinate $i$ if its
half-open interval $[s_{i-1},s_i)$ contains a marked point.
Each interval has length at most one, so exactly $r$ distinct
coordinates are selected.  The probability of selecting coordinate
$i$ is $x_i$.  The selected subset is constant between successive
fractional parts of the $s_i$.  There are at most $d$ such intervals
of positive length, because $s_0$ and $s_d$ have the same fractional
part.  Their lengths are the weights $w_h$, and their selected
coordinate projections are $P_h$.  Taking diagonal entries proves
\eqref{eq:projectionframe93}.  Endpoints have zero measure and may
be assigned either adjacent choice.  Rational eigenvalues give rational
weights by this construction; no rational eigenbasis is presumed.

Choose isometries $J_h:\C^r\to\C^d$ with $J_hJ_h^*=P_h$ and put
\[
 C_{yh}=\sqrt{w_h/r}\,J_h^*,\qquad
 D_{yh}=J_h^*/\sqrt r,\qquad V_{yh}=C_{yh}C_0^+.
\]
On the initial Schmidt support, the common Gram sum gives
$\sum_hV_{yh}^*V_{yh}=I$ and $V_{yh}C_0=C_{yh}$.
Thus $V_{yh}$ is a recorded instrument, and $|D_{yh}\rangle$ is an
independently prepared normalized fresh vector.  The same maps may be
used for every $y$.  The receiver centered swap is
$w_h(I_{r^2}-dS_r)/r^2$.  Declare the scalar hypothesis on the
antisymmetric receiver projection for equal labels and on the symmetric
projection for unequal labels.  Complementary effects complete the
measurement.  These decisions attain each positive-part bound, proving
sufficiency, including $r=1$.  Applying the same construction without
a prescribed commuting eigenbasis proves sufficiency of
\eqref{eq:optimalgrams93} with the common Gram sums.  Countable
normal forms are covered by the positive trace tails of
Lemma~\ref{lem:countable91}.
\end{proof}
In particular the initial state need not be maximally mixed when
$r<d$.  Its largest Schmidt weight must be at most $1/r$.
When $r=d$, the condition reduces to $\rho=I/d$.
The finite interval construction concerns a known initial spectrum,
not construction of the second-moment family itself.

\subsection{Counting the first acquisition as well}
Let $\mathfrak R_{q;k,\ell}$ impose the additional condition that the
entire reference accompanying the first input has dimension at most
$q$, conditional on the public preparation record.  Mixed states on
$I_1\otimes\C^q$ are allowed.  The old receiver after its recorded
instrument has dimension at most $k$ and the independent second
reference at most $\ell$, as before.  Workspaces discarded before a
specified cut do not enlarge a counted register.  Classical records
and the final joint receiver measurement have the same conventions as
Definition~\ref{def:rankclass92}.
\begin{corollary}[Three specified dimension cuts]\label{cor:threecuts93}
For $1\le q,k,\ell\le d$, put $s=\min\{q,k,\ell\}$.
On the same finite family the complete classes
$\mathfrak R_{q;k,\ell}$ have exact optimal values
\begin{equation}\label{eq:threecuts93}
 \begin{split}
 P_{q;k,\ell}^{\rm eq}
      &=\frac12+\frac{t^2(ds-1)}{2ds(d+1)},\\
 P_{q;k,\ell}^{\rm b}
      &=\frac{d+1}{2(d+1)-t^2}
          +\frac{t^2(s-1)}{2s[2(d+1)-t^2]}.
 \end{split}
\end{equation}
The second line uses exactly the prior of
Theorem~\ref{thm:operationalhierarchy90}.  Both optima are attained by
two independently prepared maximally entangled vectors on any fixed
rank-$s$ input subspace, with no intervening receiver instrument and
no feedback.  No first-call reference of dimension $d$ is required
unless $s=d$.
\end{corollary}
\begin{proof}
Resolve a mixed first preparation into a finite pure-state mixture in
$I_1\otimes\C^q$ and record its mixture index classically.  This
recovers its old score if the index is ignored, without adding a quantum
purification.  For each component $\operatorname{rank}C_0\le q$.
The Kraus refinement of Lemma~\ref{lem:ranknormal92} then gives
$\operatorname{rank}A_{yh}\le\min\{q,k\}$, while
$\operatorname{rank}b_{yh}\le\ell$.  The two rank-sensitive swap
bounds apply with $s$.  The trace sums, constant-decision baselines
and positive-part optimizations in Theorem~\ref{thm:rankspectrum92}
give both upper bounds.  Averaging over initial mixtures and public
seeds, and then taking norm limits, preserves these bounds.

For attainment fix an isometry $J:\C^s\to\C^d$ and use
$C_0=D=J^*/\sqrt s$.  Keep the first receiver unchanged and prepare
the second vector freshly.  For every first label the old Gram and
the fresh Gram are both $JJ^*/s$.  The two final spectral readouts
used in Theorem~\ref{thm:rankspectrum92} attain the bounds.  All
three counted registers have dimension $s$.  This proves the assertion
also at $s=1$, where the antisymmetric receiver projection is zero.
\end{proof}
The three bounds apply at specified cuts.  They are not a cap on the
total workspace of an arbitrary implementation or a removal of the
reset architecture.  In particular, saturation remains the reset
value, not the unrestricted coherent-acquisition value.

\subsection{A stable common-projection theorem}
For density matrices $a,b$ with smaller rank at most $r$, set
\[
 \Delta_{d,r}(a,b)=d-r^{-1}
  -\|(\sqrt a\otimes\sqrt b)(I-dS)(\sqrt a\otimes\sqrt b)\|_1.
\]
When $(d,r)\ne(2,1)$ define
\begin{equation}\label{eq:rankrigidityconstant93}
 c_{d,r}=d-1-r^{-1}>0,\qquad
 K_{d,r}=r+\frac{(3+\sqrt2)^2}{c_{d,r}}.
\end{equation}
\begin{lemma}[Quantitative common support]\label{lem:projectionstability93}
For $d\ge2$, $1\le r\le d$ and $(d,r)\ne(2,1)$, there exists
a rank-$r$ projection $P$ such that
\begin{equation}\label{eq:projectionstability93}
 \max\{\|a-P/r\|_1^2,\|b-P/r\|_1^2\}
       \le K_{d,r}\Delta_{d,r}(a,b).
\end{equation}
Only the smaller Gram rank need be bounded by $r$; neither
invertibility nor equality of the two supports is assumed.
\end{lemma}
\begin{proof}
Put $f=\|\sqrt a\sqrt b\|_1$, $\eta=1-f^2$ and
$e=\tr(ab)-f^2/r$.  The overlap rank bound gives $e\ge0$, and
$f\le1$ gives $\eta\ge0$.  The proof of
Lemma~\ref{lem:rankswap92} gives
\begin{equation}\label{eq:rankdefects93}
              \Delta_{d,r}(a,b)\ge c_{d,r}\eta+e.
\end{equation}
Extend a right polar factor to a unitary $U$ so that
$Z=\sqrt a\sqrt bU=\sqrt{\sqrt a b\sqrt a}\succeq0$.
The Hilbert--Schmidt identity and the Schatten product inequality give
\[
 \|\sqrt bU-\sqrt a\|_2^2=2(1-f),\qquad
 \|Z-a\|_1\le\sqrt{2(1-f)}\le\sqrt{2\eta}.
\]
Choose any rank-$r$ projection $P$ containing the range of $Z$.
This is possible since $\operatorname{rank}Z\le r$.  Since
$\tr Z=f$, one has
\[
 \|Z-fP/r\|_2^2=e,\qquad
 \|Z-fP/r\|_1\le\sqrt{re}.
\]
These formulas remain valid at $f=0$, without normalizing $Z$.
As $1-f\le\sqrt\eta$, the triangle inequality yields
\[
       \|a-P/r\|_1\le(1+\sqrt2)\sqrt\eta+\sqrt{re}.
\]
The normalized coefficient vectors for $\sqrt a$ and $\sqrt bU$
have inner product $f$ and receiver marginals $a,b$.
Their pure-state trace distance is $2\sqrt{1-f^2}$, by its
rank-two spectral calculation.  Contractivity under partial trace
therefore gives $\|a-b\|_1\le2\sqrt\eta$, the trace-distance/root-fidelity estimate of
Fuchs--van de Graaf~\cite[Theorem~1]{FvdG93}.  Consequently
\[
       \|b-P/r\|_1\le(3+\sqrt2)\sqrt\eta+\sqrt{re}.
\]
Weighted Cauchy--Schwarz bounds the square of either right side by
$[r+(3+\sqrt2)^2/c_{d,r}](c_{d,r}\eta+e)$.
Now use \eqref{eq:rankdefects93}.  Every argument used matrices on
finite supports and a unitary extension of a partial isometry; no
inverse of a singular matrix or limit of an equality statement occurs.
\end{proof}
The constants in this lemma are explicit but not asserted to be optimal.
The projection may vary between branches; the theorem does not force
all receivers to occupy one common subspace throughout the protocol.

\begin{corollary}[Near-optimal projection decompositions]
\label{cor:nearoptimalframes93}
Assume $(d,r)\ne(2,1)$ and a strategy in
$\mathfrak R_{k,\ell}$ has reward $v\ge p_r-\varepsilon$.
Choose a dimension-preserving refined normal form; on its nonzero
branches put
\[
 w_{yh}=\tr A_{yh},\qquad a_{yh}=A_{yh}/w_{yh},\qquad
 \mu_{yh}=w_{yh}/d,\qquad
 B_\varepsilon=\frac{2K_{d,r}d(d+1)}{t^2}\varepsilon.
\]
There are rank-$r$ projections $P_{yh}$ such that
\begin{align}
 \sum_{y,h}\mu_{yh}
  \bigl(\|a_{yh}-P_{yh}/r\|_1^2+
              \|b_{yh}-P_{yh}/r\|_1^2\bigr)&\le2B_\varepsilon,
 \label{eq:nearoptimalframes93}\\
 \frac1d\sum_y\left\|\rho-\sum_hw_{yh}P_{yh}/r\right\|_1^2
          &\le B_\varepsilon.
 \label{eq:framebarycenter93}
\end{align}
Finite and countable recorded instruments are allowed.  The weights
$\mu_{yh}$ sum to one and are Gram weights, not probabilities under a
chosen hypothesis.  For public mixtures the same statements hold with
the public seed adjoined to the refined record, using the common
purified initial Gram matrix.
\end{corollary}
\begin{proof}
Subtract \eqref{eq:unsaturatedrankupper93} from the attained class
value.  Since $\sum_{y,h}w_{yh}=d$, this gives
\[
 \sum_{y,h}\mu_{yh}\Delta_{d,r}(a_{yh},b_{yh})
       \le\frac{\varepsilon}{\alpha d}
       =\frac{2d(d+1)}{t^2}\varepsilon.
\]
Apply Lemma~\ref{lem:projectionstability93} to each branch and sum.
For the second estimate, each fixed $y$ has $\sum_hw_{yh}=1$ and
$\sum_hw_{yh}a_{yh}=\rho$.  The triangle inequality followed by
Jensen's inequality bounds its squared residual by
$\sum_hw_{yh}\|a_{yh}-P_{yh}/r\|_1^2$.
Average over $y$.  The proof uses nonnegative sums with finite total
weight, so countable traces pass to the limit.  The seed construction
in Lemma~\ref{lem:resetnormal91}, followed by the dimension-preserving
refinement, preserves the common Gram sum and the achieved reward.
\end{proof}
The total $\mu$-weight on branches for which either deviation is at
least $u>0$ is at most $2B_\varepsilon/u^2$.  This is a quantitative
statement at the relative loss scale $\varepsilon/t^2$, not a claim
of vanishing error at fixed nonzero calibration defect.

\subsection{A spectral loss certificate}
Let
\[
 \mathcal C_{d,r}=\{\sigma\in\mathcal D_d:\sigma\preceq I/r\},
 \qquad e_r(\rho)=\tr(\rho-I/r)_+.
\]
\begin{proposition}[Distance to the optimal initial spectra]
\label{prop:spectralloss93}
For every density matrix $\rho$,
\begin{equation}\label{eq:capdistance93}
 \inf_{\sigma\in\mathcal C_{d,r}}\|\rho-\sigma\|_1=2e_r(\rho).
\end{equation}
For $(d,r)\ne(2,1)$, any strategy in $\mathfrak R_{k,\ell}$
with initial Gram $\rho$ and achieved equal-prior reward $v$ satisfies
\begin{equation}\label{eq:spectralloss93}
 p_r-v\ge \frac{2t^2}{K_{d,r}d(d+1)}\,e_r(\rho)^2.
\end{equation}
\end{proposition}
\begin{proof}
Let $Q$ be the spectral projection of $\rho$ above $1/r$.
For $\sigma\preceq I/r$, one has
$\tr Q(\rho-\sigma)\ge e_r(\rho)$.
As $\tr(\rho-\sigma)=0$, positive-part optimization gives
$\|\rho-\sigma\|_1\ge2e_r(\rho)$.
Conversely, in an eigenbasis of $\rho$, cap every eigenvalue exceeding
$1/r$ and redistribute the removed mass among eigenvalues below $1/r$.
There is enough capacity since $d/r\ge1$.  This produces a density
matrix in $\mathcal C_{d,r}$ at trace distance exactly $2e_r(\rho)$.

Every matrix $\sum_hw_{yh}P_{yh}/r$ in
\eqref{eq:framebarycenter93} belongs to $\mathcal C_{d,r}$.
Set $\varepsilon=p_r-v\ge0$ there and use
\eqref{eq:capdistance93}; rearranging $4e_r(\rho)^2\le
2K_{d,r}d(d+1)(p_r-v)/t^2$ proves the bound.
\end{proof}
The estimate concerns the actual first-input marginal, up to its fixed
transpose, and a specified ideal reset architecture.  Under the
calibrated perturbations of Corollary~\ref{cor:gamestability91}, the
same strategy's nominal reward differs from its observed reward by at
most the stated channel, prior and latent-law error; a preparation
defect adds its stated half-trace contribution.  That total error
must be added to $p_r-v$ before applying the certificate.  There is no
uncalibrated inference of the reset cut from a score.

\subsection{The exceptional qubit face}
\begin{proposition}[Pinching at the exceptional endpoint]
\label{prop:qubitpinching93}
Let $d=2$, $r=1$, and suppose $a$ is pure.  Put
$\mathcal P_a(b)=aba+(I-a)b(I-a)$ and
$\Delta=\Delta_{2,1}(a,b)$.  Then
\begin{equation}\label{eq:qubitpinching93}
 \|b-\mathcal P_a(b)\|_1^2=2\Delta-\Delta^2\le2\Delta.
\end{equation}
In particular $\Delta=0$ if and only if $[a,b]=0$.
If instead $b$ is the pure factor, interchange $a,b$.
\end{proposition}
\begin{proof}
By simultaneous unitary conjugation take $a=|0\rangle\langle0|$ and
$b=\left(\begin{smallmatrix}x&z\\\bar z&1-x\end{smallmatrix}\right)$.
The nonzero block of the centered filtered swap is
$\sqrt b(I-2a)\sqrt b$.  Its trace is $1-2x$ and its determinant
is $-\det b$.  Its trace norm, also when $b$ is singular, is therefore
\[
 \sqrt{(1-2x)^2+4\det b}=\sqrt{1-4|z|^2}=1-\Delta.
\]
The pinching difference has eigenvalues $\pm|z|$.
Squaring gives \eqref{eq:qubitpinching93}.  Vanishing is exactly
$z=0$, or $[a,b]=0$.
\end{proof}
Thus the exceptional equality set includes a pure factor with any
commuting mixed qubit state, not just identical or orthogonal pure
factors.  A common rank-one projection cannot replace pinching here.
The initial-spectrum theorem remains valid because
$\mathcal C_{2,1}=\mathcal D_2$.  This exception is separate from the
zero-signal case $t=0$, in which every strategy has reward $1/2$.

\section{Relation to channel metrology and measurement information}
\label{sec:comparison84}

\subsection{Support spans, channel estimation and finite pairs}
For the measurement Kraus operators
$L_{j,a}=|j\rangle\langle a|\sqrt{E_j}$, the real Hermitian
Kraus-product span is exactly $\mathcal W_E$.  For the orbit
$L_{j,a}(t)=L_{j,a}e^{-itG}$ the effective generator is $G$.
Zhou--Jiang \cite[Theorems~1--3 and Section~V]{ZJ84} establish the
Kraus-span criterion for the linear or quadratic asymptotic channel
Fisher-information rate and give error-correction attainability.
We do not claim a new metrological dichotomy or a new correction
principle.  Theorems~\ref{thm:supportkernel84} and
\ref{thm:orbittrichotomy84} identify its complete support-kernel
form and give finite-angle trace-distance inequalities with an
explicit $4mt^2\|G\|_{\mathrm{op}}^2$ remainder.
Theorem~\ref{thm:curveclassification85} adds a normalized first-order
realization for general two-sided $C^2$ effect curves; its
$\sqrt N|t|+O(Nt^2)$ bound handles rank changes without a smooth
exact Kraus hypothesis.  Theorem~\ref{thm:controlstability85} separately
budgets certified control errors.  Neither constitutes a new claim
to the underlying metrological exponents.

Sieniawski--Demkowicz-Dobrza\'nski
\cite[Section~IV, equations~(16)--(23)]{SD76} integrate channel
Fisher-information bounds to bound adaptive discrimination and
compare metrological scaling with perfect finite discrimination.
The same horizontal method is a principal antecedent of
\eqref{eq:frontupper84}.  Our support-kernel identity is stated in
measurement coordinates and our lower \eqref{eq:finiteanglelower84}
comes from an actual corrected finite-angle channel.  It does not
claim the exact optimal tester or perfect discrimination.


\subsection{Strategy classes and finite-use comparisons}
Our independent inputs are products of \emph{complete} probe--reference
systems.  Public mixtures preserve the bound; a correlated reference
shared between groups does not satisfy this definition.  Block-$b$
preparations interpolate between those inputs and unrestricted
entangled parallel inputs.  Classical feedback permits earlier
labels to change later preparations and is not a parallel strategy;
fully adaptive control additionally retains quantum memory.  The
independent-block theorem is not a classification of every feedback
model.  Theorem~\ref{thm:resetwidth88} treats the separately defined
classical-feedback reset class with arbitrary receiver memory; it
still excludes coherent import into a fresh block.

Harrow--Hassidim--Leung--Watrous \cite[Sections~3--4]{HHLW87}
construct channels perfectly distinguishable by two adaptive calls
but not by any finite entangled parallel strategy.  Those are
perfect-discrimination statements about a fixed pair of general
channels.  Corollary~\ref{cor:parallel87} concerns comparison orders
in shrinking fixed-tangent neighborhoods of classical-output
measurements; it asserts neither equality of distances nor perfect
discrimination.  There is no contradiction with that separation.

Salek--Hayashi--Winter \cite[Theorem~9 and Remark~10]{SHW87}
compare asymptotic Chernoff and Hoeffding rates for classical
feed-forward with no quantum memory at the input against nonadaptive
product strategies.  Their restrictions and fixed-pair error
exponents differ from the present hard call cap and finite trace
separation.  In particular their product result is not an upper
bound on our entangled parallel inputs.

Li--Hirche--Tomamichel \cite[Theorems~3.1--3.3]{LHT87}
characterize sequential error-exponent regions under expected and
high-probability stopping constraints in their specified adaptive,
block and nonadaptive models, under finite max-relative-entropy
hypotheses.  Our budget instead bounds the number of calls on every
record, and our parameter $s$ varies jointly with that budget.
Their general instrument discussion is not invoked as a theorem of
unrestricted quantum-memory optimality here.

Bounds in terms of sums of squared entangled group sizes, and GHZ
attainability of metrological enhancements, are standard
\cite{Hyllus78,Toth78}.  Theorem~\ref{thm:widthlaw87} does not
claim these mechanisms as new.  It proves the finite-pair trace law
for a normalized measurement tube, controls arbitrary legal quadratic
remainders, and uses a readout independent of that unknown remainder.
The factor $\sqrt{Nb}$ comes with both a block-state Bures upper and
an integer-block Bernoulli lower, not just a Fisher-information scaling.


\subsection{Memory-constrained testers and retained receiver systems}
Ohst--Zhang--Nguyen--Pl\'avala--Quintino \cite{OZNPQ89}
formulate discrimination with prescribed auxiliary dimensions as
optimization over constrained separable operators.  Their
Definition~6 specifies single-call memory-limited testers, while
Theorem~9 characterizes convex combinations of memoryless testers
by a convergent symmetric-extension hierarchy;
Sections~5.2--5.4 treat multi-call quantum and classical memory
and affine constraints on the factors of the lifted decomposition.
These are dimension and causal-transfer restrictions, not a bound on
the squared number of calls between fresh preparations.

Their Definition~23, equation~(64), describes two-call classically
adaptive testers as $T_i=\sum_zR_z\otimes S_{i|z}$, with single-call
testers $R$ and $S_{\cdot|z}$.  The first complete output is measured
and only the classical result chooses the second tester.  Their
Theorem~24 shows that this class and the parallel class are not
nested; both lie in the adaptive class.  Our inclusion and explicit
negative-partial-transpose witness in
Proposition~\ref{prop:memorycomparison89} answer the relevant
comparison without identifying these classical testers with all
unit-reset protocols.  Reset protocols may keep an unmeasured old
receiver register for a final joint measurement.  It is independent
of the complete next acquisition, but need not be separable from
all other receiver systems after receiver processing.

The constrained-separability formulation in their Section~6.2,
equations~(65)--(68), optimizes over the stated classically adaptive
set.  It is not a characterization of our larger retained-receiver
set: its cross-use separability requirement excludes the tester
\eqref{eq:memorywitness89}.  Conditional tensor independence at
preparation is not the same condition as separability of every
final tester effect across uses.  Consequently an upper bound for
that smaller optimization cannot be used as an upper for the reset
class merely by dropping a dimension cap.  Their general adaptive
optimization, when applicable to the same finite channel ensemble,
can of course bound a reset subclass; that inclusion does not
supply a sharp $Q$-dependent shrinking-tube law.  We do not claim
that no other lifted constrained-separability representation of
reset preparation could be developed.

Thus arbitrary receiver memory with zero coherent import has an
explicit operational relation to their complete-measurement
class, but it is not identified with any fixed auxiliary-dimension
slice of their hierarchy.  The contribution here is not a new
memory hierarchy or a new separability method.  It is the
support-coordinate local profile and hard-budget theorem, including
the unknown quadratic remainder and a matching lower in those
resources.  The Bell-reference witness alone establishes a model
boundary; Theorem~\ref{thm:operationalhierarchy90} now supplies a strict
exact score hierarchy for the specified finite ensemble, whereas the general memory-constrained discrimination
problem and independent priority assessment remain distinct.

\subsection{Channel Bures bounds and block testing}
Huang--Meyer--Nuradha--Wilde \cite[Theorem~17, Remark~18]{HMNW88}
restate Yuan--Fung's parallel bound \cite[equation~(23)]{YF76}.
For aligned canonical isometries put
$M_W=V_0^*(W\otimes I)V_1$, where $\|W\|\le1$,
$a_W=\|I-\operatorname{Re}M_W\|$ and $b_W=\|I-M_W\|$.
Their Theorems~17 and~19 give, respectively,
\begin{align*}
 d_B^2(\mathcal N_0^{\otimes n},\mathcal N_1^{\otimes n})
 &\le n\inf_W\{2a_W+(n-1)b_W^2\},\\
 \sup_{\text{adaptive controls}}d_B^2(\mathcal P_0^{(n)},\mathcal P_1^{(n)})
 &\le n\inf_W\{2a_W+(n-1)\sqrt{2a_W}\,b_W\}.
\end{align*}
Channel Bures distance takes the supremum over complete
probe--reference inputs.  For our chosen surrogate dilations,
$2a_W=\|W_s-W_0\|^2\le\alpha^2s^2$ and
$b_W\le\alpha s$; either expression yields
$d_B\le\alpha ns$.  Restoring the actual tube member costs
$nrs^2$ in diamond norm and hence $s\sqrt{rn}$ in Bures distance.
Thus the block upper specializes established isometric-extension
bounds; it is not a new fidelity principle.  The finite-tube
contribution is the uniform remainder, complete support cases and
matching width-dependent lower.  The reset theorem additionally
handles history-dependent allocation and receiver instruments by
conditional fidelity, rather than unconditional tensorization.

Ghosal--Halder--Patra--Sen \cite[Section~III, Theorem~1]{GHPS88}
use block-i.i.d. probes, with arbitrary within-block correlations
and independent identical repetition.  For a closed convex null
channel set $\mathsf R$ and a simple alternative $\mathcal N$,
their memory-assisted fixed-block Stein exponent is
\[
 \sup_{\sigma_{RA^{\otimes\ell}}}\inf_{\Phi\in\mathsf R}
 \frac1\ell D\!\left((\operatorname{id}_R\otimes\Phi^{\otimes\ell})(\sigma)
 \middle\|(\operatorname{id}_R\otimes\mathcal N^{\otimes\ell})(\sigma)\right).
\]
Their memory-free expression restricts $\sigma$ to the probes.
The block hierarchy is a direct antecedent of our resource taxonomy.
Our blocks may have unequal sizes and nonidentical states; the reset
class also has classical feedback and adaptive within-block control.
We consider finite unhalved trace separation in a shrinking fixed tube,
not a composite Stein exponent at fixed block size.  Neither theorem
is substituted for the other, and no first introduction of block
strategies is claimed.

\subsection{Semidefinite tangent geometry and nonemptiness}
The positive-semidefinite tangent condition $QHQ\succeq0$ and
second-order corrections are classical variational geometry;
Bonnans--Cominetti--Shapiro \cite{BCS87} study parabolic tangent
sets and semidefinite second-order regularity.  We use no new
primitive cone theorem.  Proposition~\ref{prop:nonempty87}
provides a common normalization, explicit coefficient and interval,
and a simultaneous operator-norm remainder bound for the measurement
tuple.  The sufficient bound can be large and is not a minimal
feasibility threshold.  Its exact certificate does not determine
optimal discrimination constants or synthesize quantum controls.

\subsection{Fisher frames and covariance embeddings}
Saini--Kiukas--Burgarth--Gilchrist \cite[Theorem~1]{SKBG84} compare
classical and quantum Fisher information for a varying full-rank
\emph{state} measured by one fixed informationally complete POVM,
using a state-dependent frame operator.  Here the varying object
is the measurement, the tangent has coupled normalization, and the
loss permits $N$ adaptive uses with a reference.  In the rank-one
case our support span equals the effect span, so informational
completeness excludes the linear orbit regime by
Corollary~\ref{cor:orbitexamples84}.  That intersection does not
identify the two quadratic forms or their resource models.

Mayer--Yun \cite[Theorem~3.6, Definition~3.7 and Remark~4.4]{MY84}
use a kernel embedding of operator-valued measures, a kernel-dependent
discrepancy, and a tomographic Gram operator whose nullspace records
unobserved state directions.  Our support span contains entire
matrix blocks $P_j\mathcal H_dP_j$, rather than one observable
per effect; this distinction disappears for rank-one effects but
not in general.  Theorem~\ref{thm:supportkernel84} concerns a
coupled measurement tangent, and \eqref{eq:frontupper84} concerns
adaptive finite-use loss.  Neither kernel embedding nor injectivity
alone supplies those statements.

\subsection{Known symmetry and tomography primitives}
The full primary text of Yoshida--Okigami--Posta--Grinko
\cite{YOPG84} distinguishes known-symmetry estimation from known-pair
discrimination.  Theorem~7 gives invariant-state copy complexity
$\Theta((D_{\rm st}-1+\log(1/\eta))/\epsilon^2)$, where
$D_{\rm st}=\sum_\lambda m_\lambda^2$; $D_{\rm st}=1$ is a
known singleton.  Corollary~9 gives parallel covariant-channel bounds
\[
 O((D'_G+K_G\log(1/\eta))/\epsilon_\diamond^2),\qquad
 O((C'_G+\log(1/\eta))/\epsilon_{\rm tr}^2),
\]
in one-use halved diamond distance and normalized-Choi trace distance,
respectively, for $0<\epsilon\le1$, $0<\eta\le1/2$.
Here $D_G=\sum_\lambda m_\lambda M_\lambda$ counts Choi commutant
parameters before trace preservation,
$K_G=\sum_\lambda m_\lambda$,
$D'_G=D_G+K_G\log(4L)$, and
$C'_G=d_{\mathcal I}^{-1}(\sum_\lambda m_\lambda
\sqrt{d_\lambda M_\lambda})^2$; $L$ counts occupied blocks.
Theorem~13 does not give a general matching inverse-square diamond
lower: its bound is $\Omega_\xi(D_G/\epsilon_\diamond^{\xi/2})$
under its affine-dimension hypothesis.

For permutation covariance on $\ell$ systems of fixed local dimension
$q$, Theorems~12 and~15 match the Choi order
$\ell^{q^4-q^2}\epsilon_{\rm tr}^{-2}$ at constant confidence.
Theorem~18 approximates the Hayashi measurement in
$O((s+q)\operatorname{polylog}(s,q,\tau^{-1},\zeta^{-1}))$ gates,
with coupled output error at most $\tau$ except with probability
$\zeta$.  This is not a diamond-error certificate for our recovery.
Known-symmetry estimation, its parameter dimensions, and that
specialized gate construction do not themselves give our support
kernel or finite-pair curve theorem.  Nor do our theorems claim
faster learning on those symmetry classes.

For ordinary tomography the main-text substitutions remain explicit.
Mele--Bittel \cite[Theorem~III.3]{MB67} give the sufficient
constant-confidence order $O(dkr\epsilon^{-2})$ in halved diamond
loss, with $r\le dk$ for measurement channels.
Zambrano--Ramos-Calderer--Kueng \cite[Theorem~2]{ZRK76} give the
sufficient single-copy worst-input outcome-TV bound
\[
 m\ge\frac{8[d^2+\epsilon(d^2+1)/6]}{\epsilon^2}
          \log\frac{2^{k+1}9^{2d}}{\eta}.
\]
Either loss bounds the component norm error by $\epsilon$ after
output dephasing where needed.  A legal estimator need not be
balanced: apply Theorem~\ref{thm:affinerepair83} before using the
interior metric.  The resulting future loss is at most
$8k\sqrt N\epsilon$.  With $\epsilon=\delta/(8k\sqrt N)$ the
respective sufficient orders are
\[
 O(d^2k^4N\delta^{-2}),\qquad
 O\!\left(k^2N\delta^{-2}d^2[d+k+\log(1/\eta)]\right).
\]
These are upper-bound substitutions, not optimality claims about
the cited estimators.  Our component procedure credits the existing
binary operator-norm primitive and uses a different confidence
allocation.  The full calculations and confidence restrictions
are retained in the supplement, Section~\ref{sec:priority83}.

\subsection{What remains beyond the theorems}
The orbit theorem and Theorem~\ref{thm:curveclassification85} provide
matching scales on each fixed two-sided smooth curve with nonzero
first derivative, including second-order rank openings.  Their
constants need not be uniform across curves.  Arbitrary-pair
midpoint equivalence and full-body multi-outcome entropy are not
consequences of these local statements.  Growing-$k$ minimax sharpness,
polynomial entropy-optimal dictionaries, and general readout synthesis
remain separate mathematical questions.  Independent expert priority
assessment is not inferred from the absence of an identical statement
in this targeted comparison.

\subsection{Tangent neighborhoods and independent acquisition}
Theorem~\ref{thm:tubedichotomy86} fixes $E,H$ and an explicit
second-order remainder allowance, rather than assuming a particular
smooth family of Kraus operators.  It treats the full one-sided
positive tangent cone.  A nonzero missing-support compression has a
classical impossible-event witness; only within the cone's lineality
space does the support-generator test select the coherent branch.
The latter still uses the credited correction mechanism.
Theorem~\ref{thm:producttrichotomy86} separates these mechanisms by
an explicit product probe--reference restriction.  Its upper uses
standard root-fidelity inequalities and tensor multiplicativity
\cite[Chapter~3]{Watrous65}, not a new fidelity principle or a bound
on all nonadaptive entangled strategies.  The intervening parallel
class is addressed instead by the new proof in
Section~\ref{sec:width87}.  Its conclusion matches the adaptive
local order and does not promote the product upper to that class.

The higher-order statement requires the finite remainder
\eqref{eq:powerjet86}.  Proposition~\ref{prop:mixedrates86} instead
proves a two-scale law when an intermediate support-opening term is
present.  These results address one-sided openings and selected
higher-order approaches without identifying a first jet with an
arbitrary-pair global geometry.  No general boundary entropy or
unrestricted growing-outcome learning law is inferred.

\subsection{Autonomous coherent-memory hierarchies}
Zonnios--Binder \cite[Proposition~2]{ZB90} prove monotonicity of their
memory-parametrized process distinguishability and its upper bound by
the strategy norm.  Their Theorem~1 compiles any finite-time tester into
repeated use of one instrument with an internal counter and sufficient
coherent memory; Corollary~1 gives finite-time completeness.  Their
unknown object is a multi-time process, and their resource is the coherent
ancilla dimension of an autonomous repeated-instrument machine.
Our reset class instead allows public time-dependent history rules and
unrestricted receiver storage, but forbids coherent import across a
specified fresh-preparation cut.  Its hard square cost $Q$ is not an
ancilla dimension.  Neither their completeness theorem nor our fixed-tube
profile law supplies the other's conclusion.  The exact finite-ensemble
hierarchy of Theorem~\ref{thm:operationalhierarchy90} likewise compares
three acquisition/receiver classes, not every dimension slice of their
hierarchy.  Their work is a direct antecedent for systematic coherent-memory
interpolation; no first such framework is claimed here.

\subsection{Classical outputs and accessible post-measurement states}
Eid--Quintino \cite{EQ90} consider the quantum state returned by the
associated L\"uders instrument in addition to the classical label.
Our unknown channel returns only that label.  The references used in
Theorem~\ref{thm:operationalhierarchy90} were supplied by the experimenter;
they are not residual quantum outputs of the device.  Thus an instrument
discrimination advantage cannot be imported into the present interface.

\subsection{What the exact hierarchy adds}
The strict tester inclusion in Proposition~\ref{prop:memorycomparison89}
alone is not an operational score theorem.  Theorem~\ref{thm:operationalhierarchy90}
provides one finite classical-output ensemble, exact upper bounds for all
three classes, and matching realizations in every dimension and along a
full-rank deformation.  Ohst et al.\ \cite[Theorem~24]{OZNPQ89} already
separate other strategy classes by channel-discrimination examples; the
present assertion does not claim a first memory advantage.  Its middle
optimization permits arbitrary unmeasured receiver storage and all
unit-reset feedback.  Lemma~\ref{lem:swapfilter90} is the step that prevents
a separable-tester upper from being incorrectly assigned to this larger
class.  These exact ensemble values and the local allocation orders
answer different questions; neither is substituted for the other.

\subsection{Finite reset optimization}
The constrained-separability formulations of Ohst et al.\ \cite{OZNPQ89}
and the finite-time compilation theorem of Zonnios--Binder
\cite{ZB90} are direct antecedents for optimization with memory
restrictions.  Theorem~\ref{thm:resetvariational91} instead fixes a
complete fresh-preparation cut and retains an arbitrary old receiver;
it realizes a common-barycenter concave-envelope formula with a finite
branch bound for every finite two-call measurement decision experiment.
It does not classify all auxiliary-dimension slices of those models.
Concave envelopes, Carath\'eodory reduction and conic duality are
established tools \cite{Uhlmann91,BV91}, not separate novelty claims.
The equal-prior values in Theorem~\ref{thm:equalprior91} are an exact
application, with all class uppers proved rather than inferred from
exhibited strategies.  Antisymmetric comparison and finite design
mechanisms remain credited to \cite{ZHS90,GAE90}.  Independent specialist
assessment of the resulting optimization statements remains appropriate.


\paragraph{Rank-resolved recording cuts.}
Theorems~\ref{thm:rankroof92} and~\ref{thm:rankspectrum92} separate
an arbitrary-experiment value characterization from an explicitly
solved finite family.  The resource counts the old receiver after its
recorded instrument and the fresh reference during the second call;
it does not count the unrestricted initial reference.  In contrast,
\cite[Sections~5--6]{OZNPQ89} formulates constrained channel-memory
optimization through separability with affine constraints, including
complete classically adaptive testers in Definition~23.
The present rank envelope gives the exact common-barycenter form for
a reset measurement interface, with explicit two-coordinate values
and adjacent gains on the same ensemble.  Neither a new general
memory concept nor priority over constrained-memory optimization is
asserted.  The autonomous process hierarchy of
\cite[Theorem~1 and Corollary~1]{ZB90} saturates the unrestricted
strategy norm at sufficient coherent transfer memory.  Our hierarchy
preserves fresh reset at every threshold and saturates only the
reset optimum.  Its conditional dimension witnesses require that
architecture and calibrated channels.

\subsection{Spectral equality and its quantitative form}
The initial-state constraint $\rho\preceq I/r$ is the convex hull of
normalized rank-$r$ projections.  We use this elementary spectral
geometry as a tool; neither its convex-hull identity nor the standard
trace-distance/fidelity inequality is a priority claim.
Theorem~\ref{thm:initialspectra93} identifies this entire set with the
attainable initial Grams of the specified reset discrimination optimum.
Lemma~\ref{lem:projectionstability93} and
Proposition~\ref{prop:spectralloss93} convert its score deficit into
quantitative common-support and initial-spectrum constraints.
The proof includes the familiar purification estimate, credited to
Fuchs--van de Graaf~\cite{FvdG93}, but derives its use with singular,
unequal-rank factors directly.  The exceptional qubit equality face
is treated separately by Proposition~\ref{prop:qubitpinching93}.
The dimension cap on the initial reference is independent of the
old-output cut.  The three-cut formula does not replace the stationary
multi-time hierarchy of Zonnios--Binder or the constrained tester
framework of Ohst et al.; its endpoint is still the reset value.

\begin{thebibliography}{99}
\bibitem{Watrous65} J. Watrous, \emph{The Theory of Quantum Information},
Cambridge University Press, 2018. Theorems~2.22 and~2.26
(Choi criteria), Corollary~3.40 (trace-norm contraction),
and Section~3.3, equation~(3.306) (hybrid comparison).
\bibitem{PLS66} T. Surawy-Stepney, J. Kahn, R. Kueng and M. Gu\c{t}\u{a},
Projected least-squares quantum process tomography,
\emph{Quantum} \textbf{6} (2022), 844;
arXiv:2107.01060v2, 18 October 2022.
\bibitem{Gutoski66} G. Gutoski,
On a measure of distance for quantum strategies,
\emph{J. Math. Phys.} \textbf{53} (2012), 032202;
arXiv:1008.4636v4, 14 March 2012.
\bibitem{NZGB67} S. G. Naik, M. Zartab, N. Gisin and M. Banik,
No-Go Theorem for Generic Simulation of Qubit Channels with Finite
Classical Resources, arXiv:2501.15807v2 (16 July 2025),
Theorems~2 and~5.
\bibitem{MB67} A. A. Mele and L. Bittel,
Optimal learning of quantum channels in diamond distance,
arXiv:2512.10214v3 (2026), Theorem~I.1, Lemmas~III.7--III.8,
Corollary~III.9 and Sections~III.1.1, IV.1.1.
\bibitem{OG67} A. Oufkir and F. Girardi,
Improved Lower Bounds for Learning Quantum Channels in Diamond
Distance, arXiv:2601.04180v4 (2026).
\bibitem{DKG67} R. Demkowicz-Dobrza\'nski, J. Ko\l ody\'nski
and M. Gu\c{t}\u{a}, The elusive Heisenberg limit in quantum-enhanced
metrology, \emph{Nature Communications} \textbf{3} (2012), 1063.
DOI: 10.1038/ncomms2067; arXiv:1201.3940v2, equations~(4)--(6)
(classical simulation), equations~(10)--(13) (channel extension),
and the classical-simulation argument in Methods.
\bibitem{YBCH68} Y.~Yang, G.~Bai, G.~Chiribella and M.~Hayashi,
Compression for quantum population coding,
\emph{IEEE Transactions on Information Theory} \textbf{64} (2018), no.~7,
doi:10.1109/TIT.2017.2788407;
arXiv:1701.03372v8.
\bibitem{SHW68} F.~Salek, M.~Hayashi and A.~Winter,
Usefulness of adaptive strategies in asymptotic quantum channel
 discrimination, \emph{Phys. Rev. A} \textbf{105} (2022), 022419;
arXiv:2011.06569v3 (17 March 2024).
\bibitem{CMW70} T.~Cooney, M.~Mosonyi and M.~M.~Wilde,
Strong converse exponents for a quantum channel discrimination problem
and quantum-feedback-assisted communication,
arXiv:1408.3373v3 (26 February 2016), p.~5 and Sections~4.1--4.2.
\bibitem{Acin70} A.~Ac\'in,
Statistical distinguishability between unitary operations,
\emph{Phys. Rev. Lett.} \textbf{87} (2001), 177901.
DOI: 10.1103/PhysRevLett.87.177901.
\bibitem{DFY70} R.~Duan, Y.~Feng and M.~Ying,
Entanglement is not necessary for perfect discrimination between
unitary operations, \emph{Phys. Rev. Lett.} \textbf{98} (2007), 100503.
DOI: 10.1103/PhysRevLett.98.100503.

\bibitem{DW72} S. Das and M. M. Wilde,
Quantum reading capacity: General definition and bounds,
\emph{IEEE Trans. Inform. Theory} \textbf{65} (2019), 7566--7583;
arXiv:1703.03706v3; doi:10.1109/TIT.2019.2929925.
\bibitem{WBHK72} M. M. Wilde, M. Berta, C. Hirche and E. Kaur,
Amortized channel divergence for asymptotic quantum channel discrimination,
\emph{Lett. Math. Phys.} \textbf{110} (2020), 2277--2336;
arXiv:1808.01498v2; doi:10.1007/s11005-020-01297-7.
\bibitem{WW72} X. Wang and M. M. Wilde,
Resource theory of asymmetric distinguishability for quantum channels,
\emph{Phys. Rev. Research} \textbf{1} (2019), 033169;
arXiv:1907.06306v2; doi:10.1103/PhysRevResearch.1.033169.
\bibitem{PPKKO72} Z. Pucha\l a, \L. Pawela, A. Krawiec, R. Kukulski and
M. Oszmaniec, Multiple-shot and unambiguous discrimination of von Neumann
measurements, \emph{Quantum} \textbf{5} (2021), 425;
arXiv:1810.05122v4; doi:10.22331/q-2021-04-06-425.


\bibitem{SZ74} M. Sedl\'ak and M. Ziman,
Optimal single-shot strategies for discrimination of quantum measurements,
\emph{Phys. Rev. A} \textbf{90} (2014), 052312;
arXiv:1408.0934; doi:10.1103/PhysRevA.90.052312.
Section VI, Example 4, equation (37).
\bibitem{PPKK74} Z. Pucha\l a, \L. Pawela, A. Krawiec and R. Kukulski,
Strategies for optimal single-shot discrimination of quantum measurements,
\emph{Phys. Rev. A} \textbf{98} (2018), 042103;
arXiv:1804.05856; doi:10.1103/PhysRevA.98.042103. Theorem 1.
\bibitem{FM75} J.~Fiur\'a\v sek and M.~Mi\v cuda,
Optimal two-copy discrimination of quantum measurements,
\emph{Phys. Rev. A} \textbf{80} (2009), 042312;
arXiv:0909.2940; doi:10.1103/PhysRevA.80.042312.

\bibitem{KPP76} A.~Krawiec, \L.~Pawela and Z.~Pucha\l a,
Discrimination of POVMs with rank-one effects,
\emph{Quantum Information Processing} \textbf{19} (2020), 428;
arXiv:2002.05452v1; doi:10.1007/s11128-020-02883-3.
\bibitem{DBSA76} C.~Datta, T.~Biswas, D.~Saha and R.~Augusiak,
Perfect discrimination of quantum measurements using entangled systems,
\emph{New Journal of Physics} \textbf{23} (2021), 043021;
arXiv:2012.07069v2; doi:10.1088/1367-2630/abecaf.
\bibitem{FI76} A.~Fujiwara and H.~Imai,
A fibre bundle over manifolds of quantum channels and its application
to quantum statistics,
\emph{J. Phys. A: Math. Theor.} \textbf{41} (2008), 255304;
doi:10.1088/1751-8113/41/25/255304. Theorems~4 and~5.
\bibitem{DM76} R.~Demkowicz-Dobrza\'nski and L.~Maccone,
Using entanglement against noise in quantum metrology,
\emph{Phys. Rev. Lett.} \textbf{113} (2014), 250801;
arXiv:1407.2934v3; doi:10.1103/PhysRevLett.113.250801.
\bibitem{KGAD76} S.~Kurdzia\l ek, W.~G\'orecki, F.~Albarelli
and R.~Demkowicz-Dobrza\'nski,
Using adaptiveness and causal superpositions against noise in quantum metrology,
\emph{Phys. Rev. Lett.} \textbf{131} (2023), 090801;
arXiv:2212.08106v2; doi:10.1103/PhysRevLett.131.090801.
\bibitem{YF76} H.~Yuan and C.-H.~F.~Fung,
Fidelity and Fisher information on quantum channels,
arXiv:1506.00819v3 (2017), Section~III, equation~(26).
\bibitem{SD76} S.~Sieniawski and R.~Demkowicz-Dobrza\'nski,
Adaptive quantum channel discrimination using methods of quantum metrology,
\emph{New Journal of Physics} \textbf{28} (2026), 024502;
arXiv:2510.15506v3; doi:10.1088/1367-2630/ae3edd.
\bibitem{Dittmann76} J.~Dittmann,
Explicit formulae for the Bures metric,
\emph{J. Phys. A: Math. Gen.} \textbf{32} (1999), 2663--2670;
doi:10.1088/0305-4470/32/14/007;
arXiv:quant-ph/9808044, entitled
\emph{Note on explicit formulae for the Bures metric}.
\bibitem{SZGeom76} H.-J.~Sommers and K.~\.{Z}yczkowski,
Bures volume of the set of mixed quantum states,
\emph{J. Phys. A: Math. Gen.} \textbf{36} (2003), 10083--10100;
arXiv:quant-ph/0304041v3; doi:10.1088/0305-4470/36/39/308.
\bibitem{ZRK76} L.~Zambrano, S.~Ramos-Calderer and R.~Kueng,
Fast quantum measurement tomography with optimal error bounds,
\emph{Quantum} \textbf{10} (2026), 2162;
arXiv:2507.04500v3; doi:10.22331/q-2026-07-15-2162.
\bibitem{BDPS76} A.~Bisio, G.~M.~D'Ariano, P.~Perinotti and M.~Sedl\'ak,
Quantum learning algorithms for quantum measurements,
\emph{Physics Letters A} \textbf{375} (2011), 3425--3434;
arXiv:1103.0480v2; doi:10.1016/j.physleta.2011.08.002.
\bibitem{HB76} B.~L.~Higgins, D.~W.~Berry, S.~D.~Bartlett,
M.~W.~Mitchell, H.~M.~Wiseman and G.~J.~Pryde,
Demonstrating Heisenberg-limited unambiguous phase estimation
without adaptive measurements,
\emph{New Journal of Physics} \textbf{11} (2009), 073023;
arXiv:0809.3308v4; doi:10.1088/1367-2630/11/7/073023.
\bibitem{KLY76} S.~Kimmel, G.~H.~Low and T.~J.~Yoder,
Robust calibration of a universal single-qubit gate set via robust phase estimation,
\emph{Phys. Rev. A} \textbf{92} (2015), 062315;
doi:10.1103/PhysRevA.92.062315;
erratum, \emph{Phys. Rev. A} \textbf{104} (2021), 069901,
doi:10.1103/PhysRevA.104.069901; arXiv:1502.02677v3.

\bibitem{Hyllus78} P.~Hyllus, W.~Laskowski, R.~Krischek,
C.~Schwemmer, W.~Wieczorek, H.~Weinfurter, L.~Pezz\'e and A.~Smerzi,
Fisher information and multiparticle entanglement,
\emph{Phys. Rev. A} \textbf{85} (2012), 022321;
arXiv:1006.4366v3; doi:10.1103/PhysRevA.85.022321.
Section~III.A, Observation~1, equation~(9).
\bibitem{Toth78} G.~T\'oth,
Multipartite entanglement and high-precision metrology,
\emph{Phys. Rev. A} \textbf{85} (2012), 022322;
arXiv:1006.4368v3; doi:10.1103/PhysRevA.85.022322.
Observation~3, equation~(6).
\bibitem{CZY78} K.~Chen, Z.~Zhang and N.~Yu,
Optimal lower bound for quantum channel tomography in away-from-boundary regime,
arXiv:2601.10683v1 (15 January 2026), Theorems~1.1, 3.2 and~4.2.
\bibitem{CLL78} S.~Chen, J.~Li and A.~Liu,
An optimal tradeoff between entanglement and copy complexity for state tomography,
arXiv:2402.16353v1 (26 February 2024).
\bibitem{NZ78} A.~Nayak and X.~Zhou,
Optimal Low-Rank Quantum State Tomography with Bounded-Sample Joint Measurements,
arXiv:2609.10514v1 (9 September 2026), Theorems~1.1--1.2.

\bibitem{BPR80} S.~Basu, R.~Pollack and M.-F.~Roy,
\emph{Algorithms in Real Algebraic Geometry}, second edition,
Algorithms and Computation in Mathematics, vol.~10,
Springer, Berlin, Heidelberg, 2006;
doi:10.1007/3-540-33099-2.
Theorems~2.77 and~2.80 and Chapter~14.
\bibitem{PSTW80} A.~Pelecanos, J.~Spilecki, E.~Tang and J.~Wright,
Mixed state tomography reduces to pure state tomography,
arXiv:2511.15806v1 (19 November 2025).
Theorems~1.1--1.3 and~4.2, Section~3.3, Figures~8--9.

\bibitem{ZJ84} S.~Zhou and L.~Jiang,
Asymptotic theory of quantum channel estimation,
\emph{PRX Quantum} \textbf{2} (2021), 010343;
arXiv:2003.10559v3; doi:10.1103/PRXQuantum.2.010343.
\bibitem{KL84} E.~Knill and R.~Laflamme,
A theory of quantum error-correcting codes,
\emph{Phys. Rev. A} \textbf{55} (1997), 900--911;
arXiv:quant-ph/9604034; doi:10.1103/PhysRevA.55.900.
\bibitem{SKBG84} R.~Saini, J.~Kiukas, D.~Burgarth and A.~Gilchrist,
Characterizing Fisher information of quantum measurement,
arXiv:2512.15428v1 (2025).
\bibitem{MY84} P.~N.~Mayer and H.~Yun,
Kernel embedding for operator-valued measures and its application to
quantum tomography, arXiv:2605.25146v1 (2026).
\bibitem{YOPG84} S.~Yoshida, K.~Okigami, P.~M.~Posta and D.~Grinko,
Optimal learning of covariant quantum states and channels,
arXiv:2609.39280v1 (2026), Theorem~7, Corollary~9 and
Theorems~12--13, 15 and~18; full text consulted 5 October 2026.

\bibitem{BCS87} J.~F.~Bonnans, R.~Cominetti and A.~Shapiro,
Second order optimality conditions based on parabolic second order
tangent sets, \emph{SIAM J. Optim.} \textbf{9} (1999), 466--492.
\bibitem{HHLW87} A.~W.~Harrow, A.~Hassidim, D.~W.~Leung and J.~Watrous,
Adaptive versus non-adaptive strategies for quantum channel
discrimination, \emph{Phys. Rev. A} \textbf{81} (2010), 032339.
\bibitem{SHW87} F.~Salek, M.~Hayashi and A.~Winter,
Usefulness of adaptive strategies in asymptotic quantum channel
discrimination, \emph{Phys. Rev. A} \textbf{105} (2022), 022419.
\bibitem{LHT87} Y.~Li, C.~Hirche and M.~Tomamichel,
Sequential quantum channel discrimination, arXiv:2210.11079.
\bibitem{HMNW88} Z.~Huang, J.~J.~Meyer, T.~Nuradha and M.~M.~Wilde,
Query complexities of quantum channel discrimination and estimation:
A unified approach, arXiv:2511.10832v2 (2026).
\bibitem{GHPS88} P.~Ghosal, P.~Halder, A.~Patra and A.~Sen,
Quantum hypothesis testing of non-mixed-unitarity: A multifaceted
hierarchy of quantum channel discrimination, arXiv:2609.40355v1 (2026).
\bibitem{OZNPQ89} T.-A.~Ohst, S.~Zhang, H.~C.~Nguyen,
M.~Pl\'avala and M.~T.~Quintino,
Characterising memory in quantum channel discrimination via
constrained separability problems,
\emph{Quantum} \textbf{10} (2026), 1988;
doi:10.22331/q-2026-01-28-1988; arXiv:2411.08110v2.
\bibitem{ZB90} M.~Zonnios and F.~C.~Binder,
Distinguishing quantum processes with bounded coherent memory,
Preprint, 2026. arXiv:2606.19511v1.
\bibitem{EQ90} C.~Eid and M.~T.~Quintino,
Post-measurement states are (very) useful for measurement discrimination,
arXiv:2602.12258v2 (2026).
\bibitem{GAE90} D.~Gross, K.~Audenaert and J.~Eisert,
Evenly distributed unitaries: on the structure of unitary designs,
arXiv:quant-ph/0611002 (2006).
\bibitem{ZHS90} M.~Ziman, T.~Heinosaari and M.~Sedl\'ak,
Unambiguous comparison of quantum measurements,
arXiv:0905.4445 (2009).
\bibitem{Uhlmann91} A.~Uhlmann, Roofs and convexity,
Entropy \textbf{12} (2010), 1799--1832.
\bibitem{BV91} S.~Boyd and L.~Vandenberghe, \emph{Convex Optimization},
Cambridge University Press, Cambridge, 2004.
\bibitem{FvdG93} C. A. Fuchs and J. van de Graaf,
Cryptographic distinguishability measures for quantum-mechanical states,
\emph{IEEE Trans. Inform. Theory} \textbf{45} (1999), 1216--1227;
arXiv:quant-ph/9712042, Theorem~1.
\end{thebibliography}

\end{document}
